\documentclass[11pt,a4paper,oneside,openany,headings=normal]{scrbook}
\usepackage{satbook}
\addbibresource{references.bib}

\begin{document}
\hypersetup{pageanchor=false}

\begin{titlepage}
\thispagestyle{empty}
\begin{tikzpicture}[remember picture,overlay]
  \fill[Ink] (current page.north west) rectangle
    ([yshift=-108mm]current page.north east);
  \fill[Blue] ([yshift=-108mm]current page.north west) rectangle
    ([yshift=-111mm]current page.north east);
  \draw[Rule,line width=.5pt]
    ([xshift=26mm,yshift=28mm]current page.south west) --
    ([xshift=-23mm,yshift=28mm]current page.south east);
\end{tikzpicture}
\vspace*{16mm}
{\sffamily\bfseries\color{white}\Large CHUYÊN KHẢO\par}
\vspace{10mm}
{\sffamily\bfseries\color{white}\fontsize{30}{35}\selectfont
Biểu diễn SAT tối ưu\par}
\vspace{4mm}
{\sffamily\color{white}\fontsize{19}{24}\selectfont
cho các bài toán tối ưu hóa tổ hợp\par}

\vspace{54mm}
{\sffamily\bfseries\Large\color{Ink}
Khảo cứu nền tảng và các họ mã hóa\par
Nguyên lý thiết kế và tái sử dụng cấu trúc\par
Nghiên cứu tình huống trong tối ưu tổ hợp\par}

\vfill
{\sffamily\bfseries\color{Ink}
Tô Văn Khánh \enspace·\enspace Kiều Văn Tuyên
\enspace·\enspace Trương Xuân Hiếu\par
Vũ Thanh Hương \enspace·\enspace Đào Xuân Nghĩa
\enspace·\enspace Nguyễn Kim Trung Đức\par}
\vspace{6mm}
{\sffamily\small\color{Ink}Bản hiệu chỉnh sau phản biện · 30/07/2026\par}
\end{titlepage}
\hypersetup{pageanchor=true}

\frontmatter
\chapter*{Lời nói đầu}
\addcontentsline{toc}{chapter}{Lời nói đầu}
Cuốn sách trình bày \emph{phép mã hóa SAT} (SAT encoding) như một thành phần
trung tâm của thuật toán giải chính xác. Phép mã hóa chuyển biến và ràng buộc
của bài toán gốc thành công thức dạng chuẩn hội \CNF. Hai công thức cùng biểu
diễn một ràng buộc có thể khác nhau đáng kể về kích thước, khả năng
\emph{lan truyền đơn vị} (unit propagation), cấu trúc mệnh đề và mức độ phù
hợp với \emph{giải tăng dần} (incremental solving). Vì thế người thiết kế cần
đồng thời trả lời bốn câu hỏi: phép mã hóa có đúng không, trạng thái nào được
biểu diễn và chia sẻ, bộ giải nhìn thấy cấu trúc gì, và kết quả tối ưu được
chứng nhận thế nào.

Ba phần của sách đi theo một đường liên tục. Phần I tổng hợp nền tảng SAT,
MaxSAT, phép chuyển CSP sang SAT, các tiêu chuẩn lan truyền, bộ công cụ ràng
buộc đếm (cardinality) và giả Boolean (pseudo-Boolean), cùng phương pháp thực
nghiệm. Phần II đặt bộ đếm dùng chung, thanh ghi thích nghi, phá đối xứng và
liên kết biểu diễn (channeling) vào các dòng nghiên cứu rộng hơn về ràng buộc
\textsc{Sequence}, phân rã ràng buộc toàn cục và cận dưới của phép mã hóa.
Phần III khảo sát bốn miền ứng dụng: lập lịch, đóng gói, các bài toán
bandwidth/antibandwidth và gán nhãn/phân bổ tần số.

Các công trình của nhóm tác giả được dùng làm \emph{nghiên cứu tình huống} để
minh họa trọn vẹn quá trình hình thành, chứng minh và đánh giá một phép mã hóa.
Mỗi chương bắt đầu từ hệ phân loại và các kết quả đại diện của cộng đồng, sau
đó kết nối chúng với bài toán, kỹ thuật và bằng chứng thực nghiệm của nhóm.

Ngôn ngữ chính là tiếng Việt. Ở lần xuất hiện đầu, sách ghi bản dịch trước và
thuật ngữ tiếng Anh trong ngoặc, chẳng hạn \emph{bộ giải} (solver),
\emph{bộ kiểm chuẩn} (benchmark) và \emph{bộ kiểm tra nghiệm} (checker).
Những tên riêng của phương pháp, ký hiệu chuẩn và từ khóa cần để tra cứu văn
liệu được giữ nguyên.

\begin{keyidea}{Cách dùng kết quả thực nghiệm}
Mỗi số liệu được đặt cạnh thiết kế thí nghiệm của công trình nguồn: tập bài
toán thử, bộ giải, thời hạn và tiêu chí kiểm chứng. Ký hiệu \OPT dành cho giá
trị có chứng nhận tối ưu; nghiệm tốt nhất đã biết được ký hiệu \BKS\
(\emph{best known solution}). Cách trình bày này cho phép đọc đúng ý nghĩa của
từng bảng và so sánh các cơ chế mã hóa trên cùng một hệ quy chiếu.
\end{keyidea}

\chapter*{Lời cảm ơn}
\addcontentsline{toc}{chapter}{Lời cảm ơn}

Nhóm tác giả trân trọng cảm ơn các đồng nghiệp, phản biện và cộng đồng
\SAT/\MaxSAT, Lập trình ràng buộc và Nghiên cứu vận trù đã công bố các kết quả
nền tảng, bộ kiểm chuẩn, bộ giải, bộ kiểm tra và gói phần mềm--dữ liệu tái lập
mà cuốn sách khảo cứu. Những nhận xét phản biện chi tiết đã giúp nhóm hoàn
thiện hệ thống hình minh họa, ví dụ tính tay, thuật toán, chỉ mục và chương kết.

Nhóm tác giả đặc biệt cảm ơn các thế hệ sinh viên và cựu sinh viên đã đồng
hành trong các seminar, trao đổi học thuật, đọc bản thảo, xây dựng ví dụ và
thử nghiệm phần mềm. Những câu hỏi từ góc nhìn người học cùng kinh nghiệm triển
khai của các bạn đã giúp nhiều khái niệm trong sách trở nên sáng rõ và gần với
thực hành nghiên cứu hơn.

\chapter*{Phạm vi khảo cứu và cách đọc}
\addcontentsline{toc}{chapter}{Phạm vi khảo cứu và cách đọc}

Sách dành cho người đọc quan tâm nghiên cứu SAT, từ học viên cao học bắt đầu
làm quen với phép mã hóa đến nhà nghiên cứu cần một khung so sánh để thiết kế
mô hình mới. Độc giả cần biết lôgic mệnh đề và các khái niệm cơ bản của tối ưu
tổ hợp; những kiến thức chuyên sâu về CDCL, lập trình ràng buộc hay quy hoạch
nguyên được giới thiệu tại nơi sử dụng.

Phạm vi được tổ chức theo ba trục:
\begin{enumerate}
  \item \textbf{Trục biểu diễn}: mã hóa trực tiếp, mã hóa hỗ trợ, mã hóa thứ
        tự, bộ đếm, cây tổng (totalizer), mạng sắp xếp/lựa chọn, BDD và các
        kiến trúc lai.
  \item \textbf{Trục suy luận}: tương đương theo phép chiếu, lan truyền đơn
        vị, mức nhất quán của miền, học mệnh đề, giải tăng dần và ghi chứng
        minh.
  \item \textbf{Trục ứng dụng}: chuỗi và cửa sổ, lập lịch, đóng gói, gán nhãn
        đồ thị, tô màu và phân bổ tần số.
\end{enumerate}

Mỗi bảng tổng hợp trả lời năm câu hỏi: bộ mã hóa biểu diễn trạng thái nào;
kích thước phụ thuộc tham số nào; lan truyền đơn vị suy ra được gì; giao diện
nào dùng được cho giải tăng dần hoặc ghép mô-đun; và cấu trúc bài toán nào phù
hợp. Năm câu hỏi này tạo thành khung lựa chọn phép mã hóa xuyên suốt cuốn sách.

Các hình TikZ nguyên bản trong sách gồm sơ đồ khái niệm và ví dụ nhỏ có thể
kiểm tra bằng tay. Các chú thích liên kết trực tiếp hình minh họa với khái
niệm, mô hình và công trình tương ứng.

\begin{summarybox}
\textbf{Lộ trình đọc nhanh.}
Người mới nên đọc Chương 1--4 rồi chọn một miền ở Phần III. Người thiết kế
ràng buộc đếm/PB nên đi từ Chương 2--3 sang Chương 5--6. Người làm mô hình
miền hữu hạn nên đọc Chương 2, 7 và một trong các Chương 8--11. Các hộp kết
quả tóm tắt bằng chứng nguồn; phần khảo cứu và bảng so sánh của từng chương
phát triển lập luận tổng quát.
\end{summarybox}

\chapter*{Quy ước đọc sơ đồ và các hộp nội dung}
\addcontentsline{toc}{chapter}{Quy ước đọc sơ đồ và các hộp nội dung}

Sơ đồ năm nút ở đầu mỗi chương là \emph{lộ trình lập luận}, đọc từ trái sang
phải. Mỗi nút biểu diễn một lớp quyết định hoặc một câu hỏi mà chương lần lượt
giải quyết; mũi tên biểu diễn quan hệ phụ thuộc trong lập luận. Các thủ tục
thực thi được trình bày bằng môi trường thuật toán và đánh số riêng.

\begin{center}
\captionof{table}{Phân cấp các môi trường dùng trong sách.}
\label{tab:box-hierarchy}
\begin{tabularx}{\textwidth}{@{}>{\bfseries\sffamily\raggedright\arraybackslash}p{.22\textwidth}YY@{}}
\toprule
Môi trường & Chức năng & Mức độ bằng chứng\\
\midrule
Định nghĩa, bổ đề, định lý & Khái niệm và kết quả hình thức &
Phải có giả thiết, phát biểu và chứng minh\\
Mệnh đề/Hệ quả & Kết quả trung gian hoặc hệ quả &
Được suy từ định nghĩa hay kết quả đã nêu\\
Ví dụ có lời giải & Tính tay trên bài toán nhỏ &
Kiểm tra trực giác và diễn giải từng bước\\
Ý tưởng cốt lõi & Ý niệm tổ chức hoặc sự trừu tượng hóa &
Lập luận khái niệm\\
Quy tắc thiết kế & Khuyến nghị kỹ thuật có điều kiện &
Nêu miền áp dụng và điểm chuyển thiết kế\\
Hộp kết quả & Số liệu hoặc kết luận từ công trình nguồn &
Gắn trích dẫn và điều kiện thí nghiệm của nguồn\\
Hộp tóm tắt & Tóm lược để tra cứu nhanh &
Tổng hợp các kết quả đã trình bày\\
\bottomrule
\end{tabularx}
\end{center}

\chapter*{Ký hiệu và quy ước}
\addcontentsline{toc}{chapter}{Ký hiệu và quy ước}
\begin{center}
\captionof{table}{Ký hiệu dùng xuyên suốt cuốn sách.}
\label{tab:notation}
\begin{tabularx}{\textwidth}{@{}>{\bfseries\sffamily}p{.24\textwidth}Y@{}}
\toprule
Ký hiệu & Ý nghĩa\\
\midrule
\(x,y,z\in\set{0,1}\) & Biến Boolean\\
\(C=\ell_1\lor\cdots\lor\ell_r\), \(F=\bigwedge C\) & Mệnh đề và công thức \CNF\\
\(X,Y\) & Biến chính và biến phụ\\
\(\operatorname{Enc}(\mathcal C;X,Y)\) & Phép mã hóa ràng buộc \(\mathcal C\)\\
\(F\equiv_X\mathcal C\) & Tương đương theo phép chiếu lên \(X\)\\
\(\operatorname{UP}(F,A)\) & Bao đóng lan truyền đơn vị dưới phép gán bộ phận \(A\)\\
\(\AMO,\ALO,\ExactlyOne\) & Nhiều nhất một, ít nhất một và đúng một\\
\(LB,UB,\OPT\) & Cận dưới, cận trên và giá trị tối ưu được chứng nhận\\
\bottomrule
\end{tabularx}
\end{center}

Chỉ số toán học bắt đầu từ \(1\), trừ khi phần cài đặt ghi rõ khác. Các miền
thời gian, tọa độ và nhãn là miền nguyên hữu hạn. Mỗi biến phụ được
đọc theo đặc tả đã nêu; một literal phụ có ngữ nghĩa dùng chung khi phép mã
hóa bảo đảm đặc tả tương ứng.

\tableofcontents
\listoffigures
\listoftables
\listofalgorithms

\mainmatter

\part{Nền tảng và tiêu chí thiết kế}
\chapter{SAT và MaxSAT như một nền tảng giải chính xác}
\chapterlead{Bộ giải SAT trả lời một câu hỏi quyết định. Để giải tối ưu tổ hợp,
ta cần thêm mô hình miền, phép biểu diễn, cơ chế tìm cận và một
\emph{model checker} (bộ kiểm tra độc lập). Bốn lớp này tạo thành một thuật
toán hoàn chỉnh.}

\index{SAT}\index{DPLL}\index{CDCL}\index{MaxSAT}
\index{incremental SAT}\index{proof logging}\index{model checker}

\flowdiagram{Mô hình miền}{Biểu diễn \CNF}{Bộ giải SAT/MaxSAT}
  {Giải mã nghiệm}{Kiểm tra và cận}

\section{Sự giao thoa của ba hướng nghiên cứu}

Phép biểu diễn SAT cho tối ưu tổ hợp là sự kết hợp của ba hướng tiếp cận cốt lõi. Thứ
nhất, \emph{propositional reasoning} (lập luận mệnh đề) đưa bài toán về \CNF
để dùng lan truyền đơn vị, phân tích xung đột và học mệnh đề. Thứ hai,
\emph{Constraint Programming} (CP, quy hoạch ràng buộc) nhìn mô hình qua biến miền
hữu hạn, giá trị hỗ trợ và các mức nhất quán. Thứ ba,
\emph{mathematical optimization} (tối ưu toán học) đặt hàm mục tiêu, mô hình nới lỏng, cận dưới và
chứng nhận tối ưu ở vị trí trung tâm.

Các công trình về \emph{arc consistency} (nhất quán cung) trong SAT chỉ ra
rằng chất lượng phép dịch còn phụ thuộc vào các suy luận miền mà bộ lan truyền
tạo được \parencite{gent2002arc,bacchus2007gac}. Một phép biểu diễn có thể đúng về
tập nghiệm nhưng làm mất suy luận vốn xuất hiện sớm trong bộ lan truyền của
ràng buộc toàn cục. SAT bù lại bằng khả năng học mệnh đề xuyên qua nhiều ràng
buộc. Vì vậy sách dùng hai cấp phân tích:
\begin{enumerate}
  \item ở cấp mô-đun, xét \emph{contract} (đặc tả ngữ nghĩa) và suy luận do
        lan truyền đơn vị;
  \item ở cấp mô hình, xét tương tác giữa các mô-đun trong CDCL và vòng tối
        ưu.
\end{enumerate}

Việc biên dịch CSP miền hữu hạn sang SAT đã phát triển thành một hướng độc lập,
với biểu diễn trực tiếp, hỗ trợ, thứ tự và lôgarit, cùng nhiều phép phân rã ràng
buộc toàn cục \parencite{tamura2009csp,petke2015bridging}. Trong góc nhìn này,
``biểu diễn SAT tối ưu'' là việc chọn một ngôn ngữ Boolean phù hợp với cấu trúc
miền, kiểu suy luận cần giữ và giao diện hàm mục tiêu của bài toán đang xét.

\begin{center}
\captionof{table}{Ba tiêu chí bổ trợ khi đánh giá một phép biểu diễn SAT.}
\label{tab:three-perspectives}
\begin{tabularx}{\textwidth}{@{}>{\bfseries\sffamily}p{.2\textwidth}YYY@{}}
\toprule
Góc nhìn & Đơn vị chính & Bằng chứng mạnh & Rủi ro khi đứng riêng\\
\midrule
SAT/CDCL & Literal và mệnh đề & Xung đột, mệnh đề học được, chứng minh &
Bỏ qua ngữ nghĩa miền\\
CP/CSP & Biến, miền, giá trị hỗ trợ & Nhất quán và lọc miền &
Không dự báo hết việc học mệnh đề\\
Tối ưu & Cận và hàm mục tiêu & Khoảng cận, cận dưới, giá trị tối ưu &
Mô hình Boolean có thể quá yếu\\
\bottomrule
\end{tabularx}
\end{center}

\section{Từ quyết định đến tối ưu}

Bài toán \SAT hỏi liệu một công thức Boolean có phép gán thỏa hay không
\parencite{cook1971complexity,biere2009handbook}. Trong ứng dụng tối ưu, ta
thường bắt đầu bằng một họ bài toán quyết định phụ thuộc cận \(B\):
\[
  P(B):\quad
  \text{tồn tại nghiệm khả thi }s\text{ sao cho }f(s)\le B?
\]
Với mục tiêu cực tiểu,
\[
  P(B)=\SAT\Longrightarrow P(B')=\SAT
  \quad\forall B'\ge B.
\]
Tính đơn điệu cho phép tìm nghiệm tối ưu trên đoạn \([LB,UB]\):
\[
  \OPT=\min\set{B\mid P(B)=\SAT}.
\]

Tìm kiếm tuyến tính từ cận trên thích hợp khi đã có nghiệm tốt và các cận kề
nhau chia sẻ nhiều cấu trúc. Tìm kiếm nhị phân giảm số lần gọi bộ giải nhưng
nhảy cận xa hơn. Với SAT tăng dần, công thức cơ sở được giữ nguyên; cận được
kích hoạt bằng các \emph{assumption} (literal giả định chỉ có hiệu lực trong
một lần gọi bộ giải) hoặc bằng các mệnh đề chỉ bổ sung
\parencite{een2003incremental}.

\begin{designrule}{Bất biến của vòng tìm cận}
Trước khi xây dựng thuật toán, cần xác định chính xác phát biểu \(P(B)\), chiều đơn điệu và
quy tắc cập nhật \(LB,UB\). Bộ giải chỉ trả \SAT hoặc \UNSAT; ý nghĩa của hai
trạng thái đối với cận thuộc về thuật toán bao ngoài.
\end{designrule}

\begin{workedexample}{tìm nghiệm tối ưu trên một dãy cận}
Giả sử \(P(B)\) hỏi liệu có lịch với thời điểm hoàn tất toàn lịch không quá \(B\), và ta biết
\(LB=4,UB=9\). Tìm kiếm nhị phân thử \(B=6\) và nhận \UNSAT, nên mọi cận
\(B\le6\) đều bị loại. Thử tiếp \(B=8\) và nhận một lịch khả thi, nên
\(UB\leftarrow8\). Cuối cùng \(P(7)=\SAT\), cho
\[
  P(6)=\UNSAT,\qquad P(7)=\SAT,\qquad \OPT=7.
\]
Mô hình thỏa ở cận \(7\) là bằng chứng cận trên; chứng minh \UNSAT của \(P(6)\)
là chứng nhận cận dưới. Hai bằng chứng gặp nhau tại giá trị tối ưu \(7\).
\end{workedexample}

\begin{figure}[tb]
\centering
\begin{tikzpicture}[satfig,x=12mm,y=10mm]
  \draw[very thick,Ink] (0,0)--(7,0);
  \foreach \x/\lab in {0/4,1/5,2/6,3/7,4/8,5/9}
    \draw[Ink] (\x,2pt)--(\x,-2pt) node[below=4pt] {\lab};
  \draw[very thick,Gold] (0,0)--(2,0);
  \draw[very thick,Teal] (3,0)--(5,0);
  \node[satwarn,rounded corners=2pt] (unsatbox) at (1.2,.95)
    {\(\UNSAT\) tại \(6\)};
  \node[sattrue,rounded corners=2pt] (satbox) at (3.8,.95)
    {\(\SAT\) tại \(7\)};
  \draw[satedge] (unsatbox.south)--(2,.08);
  \draw[satedge] (satbox.south)--(3,.08);
  \node[satlabel,below=12mm] at (2.5,0)
    {biên chuyển trạng thái\\xác định nghiệm tối ưu};
\end{tikzpicture}
\caption{Tính đơn điệu của bài toán quyết định theo cận \(B\).}
\label{fig:bound-transition}
\end{figure}

\section{CNF và biến phụ}

\subsection{Mệnh đề, literal và mô hình thỏa}

Một \emph{literal} là biến \(x\) hoặc phủ định \(\neg x\) của biến đó. Một
\emph{clause} (mệnh đề) là phép tuyển các literal; công thức \CNF là phép hội
của các mệnh đề:
\[
  F=\bigwedge_{i=1}^{m}C_i,
  \qquad
  C_i=\bigvee_{j=1}^{r_i}\ell_{ij}.
\]
Một mệnh đề được thỏa khi có ít nhất một literal đúng. Mệnh đề đơn vị có đúng
một literal; mệnh đề nhị phân thường biểu diễn phép kéo theo hoặc xung đột cục
bộ.

\subsection{Tseitin transformation (phép chuyển Tseitin)}

Biến phụ cho phép thay một biểu thức con bằng một tên ngắn. Với
\[
  z\leftrightarrow(x\land y),
\]
\CNF tương đương là
\[
  (\neg z\lor x)
  \land
  (\neg z\lor y)
  \land
  (z\lor\neg x\lor\neg y).
\]
Phép Tseitin giữ kích thước tuyến tính theo mạch Boolean
\parencite{tseitin1968}. Ba mệnh đề trên tạo một \emph{contract}
(đặc tả ngữ nghĩa): \(z\) đúng khi và chỉ khi cả \(x,y\) đúng. Nếu mô-đun sau chỉ dùng chiều
\(z\rightarrow x\land y\), có thể bỏ chiều ngược; nhưng khi đó \(z\) không còn
là tên tương đương của biểu thức.

\begin{keyidea}{Vai trò của biến phụ}
Biến phụ vừa làm gọn công thức vừa làm lộ trạng thái trung gian cho lan truyền
đơn vị và học mệnh đề. Thiết kế phép biểu diễn vì thế là thiết kế một không gian
trạng thái cho bộ giải.
\end{keyidea}

\begin{workedexample}{đọc đặc tả của biến Tseitin}
Với \(z\leftrightarrow(x\land y)\), đặt \(z=1\). Hai mệnh đề
\((\neg z\lor x)\) và \((\neg z\lor y)\) trở thành mệnh đề đơn vị, nên UP suy
\(x=y=1\). Ngược lại, đặt \(x=y=1\); mệnh đề
\((z\lor\neg x\lor\neg y)\) suy \(z=1\).

Nếu bỏ mệnh đề thứ ba, phép gán \(x=y=1,z=0\) vẫn thỏa. Khi đó hai mệnh đề còn
lại chỉ biểu diễn \(z\rightarrow x\land y\), không còn tương đương hai chiều.
Ví dụ ba biến này là phép kiểm nhanh nên làm trước khi dùng một literal phụ
trong mô-đun khác.
\end{workedexample}

\section{Từ DPLL đến CDCL}

DPLL tổ chức phép quyết định \SAT như một cây tìm kiếm có lan truyền
\parencite{biere2009handbook}. Ở mỗi nút, thuật toán chạy lan truyền đơn vị;
nếu xuất hiện mệnh đề rỗng thì nhánh hiện tại thất bại, nếu mọi mệnh đề đã thỏa
thì thu được mô hình thỏa. Nếu chưa kết luận, chọn một biến \(x\), thử \(x=1\), rồi
\(x=0\) nếu nhánh đầu thất bại. Hai nhánh bao phủ mọi phép gán của \(x\), nên
lập luận quy nạp theo số biến chưa gán cho tính đầy đủ của thuật toán.

\begin{algorithm}[H]
\caption{Khung DPLL cơ sở}
\label{alg:dpll}
\begin{minipage}{.94\textwidth}\small
\textbf{DPLL(\(F,\alpha\))}
\begin{enumerate}
  \item Lặp lan truyền đơn vị trên \(F\) dưới phép gán \(\alpha\).
  \item Nếu có mệnh đề rỗng, trả \(\UNSAT\); nếu mọi mệnh đề đã thỏa, trả
  \(\SAT\) và \(\alpha\).
  \item Chọn một biến chưa gán \(x\). Gọi
  \(\textbf{DPLL}(F,\alpha\cup\set{x=1})\); nếu nhận \(\SAT\), trả mô hình đó.
  \item Trả kết quả của
  \(\textbf{DPLL}(F,\alpha\cup\set{x=0})\).
\end{enumerate}
\end{minipage}
\end{algorithm}

DPLL cơ sở quay lui theo thứ tự quyết định và không lưu một mô tả tổng quát
của nhánh đã thất bại. CDCL giữ bộ khung quyết định--lan truyền--xung đột nhưng
ghi lý do của mỗi literal trên \emph{trail} (vết gán theo thứ tự), phân tích đồ thị kéo theo, học
một mệnh đề và thực hiện \emph{backjump} (nhảy lùi đến mức quyết định thích
hợp). \emph{Restart} (khởi động lại)
làm mới phần quyết định nhưng giữ các mệnh đề đã học. Bước chuyển này giải thích vì sao phép biểu diễn
không chỉ ảnh hưởng số nút: nó quyết định các cạnh lý do mà phân tích xung đột
có thể sử dụng.

Các bộ giải SAT hiện đại chủ yếu dựa trên
\emph{conflict-driven clause learning} (học mệnh đề theo xung đột, CDCL)
\parencite{marquessilva2009cdcl,een2003minisat}. Một vòng CDCL gồm:
\begin{enumerate}
  \item chọn một literal quyết định;
  \item chạy lan truyền đơn vị;
  \item khi có xung đột, phân tích đồ thị kéo theo;
  \item học mệnh đề và nhảy lùi;
  \item tiếp tục cho đến mô hình thỏa hoặc chứng minh \UNSAT.
\end{enumerate}

Phép biểu diễn ảnh hưởng trực tiếp ba lớp của vòng này. Mệnh đề nhị phân và tam
phân tạo đường lan truyền ngắn. Biến phụ quyết định hình dạng đồ thị kéo theo.
Cấu trúc mô-đun và tính địa phương ảnh hưởng chất lượng mệnh đề học được; chỉ số
LBD là một trong các thước đo thường dùng để quản lý chúng
\parencite{audemard2009lbd}.

Vì vậy hai phép biểu diễn có cùng số mệnh đề vẫn có thể khác về thời gian chạy.
Phép biểu diễn lớn hơn đôi khi chạy nhanh hơn nhờ các trạng thái trung gian giúp
bộ giải học mâu thuẫn có tính khái quát.

\begin{figure}[tb]
\centering
\begin{tikzpicture}[satfig,node distance=12mm and 13mm]
  \node[satbox] (d) {Quyết định\\\(x=1\)};
  \node[satbox,right=of d] (u) {Lan truyền\\đơn vị};
  \node[satwarn,rounded corners=2pt,right=of u] (c) {Xung đột};
  \node[satbox,below=of c] (a) {Phân tích\\đồ thị kéo theo};
  \node[sattrue,rounded corners=2pt,left=of a] (l) {Học mệnh đề};
  \node[satbox,left=of l] (b) {Nhảy lùi};
  \draw[satedge] (d)--(u);
  \draw[satedge] (u)--(c);
  \draw[satedge] (c)--(a);
  \draw[satedge] (a)--(l);
  \draw[satedge] (l)--(b);
  \draw[satedge] (b)--(d);
  \node[satlabel,below=5mm of l] {Phép biểu diễn quyết định độ ngắn của đường suy luận\\
    và phạm vi áp dụng của mệnh đề học được.};
\end{tikzpicture}
\caption{Một vòng CDCL nhìn từ góc độ thiết kế phép biểu diễn.}
\label{fig:cdcl-loop}
\end{figure}

\section{SAT tăng dần}

Một bộ giải tăng dần giữ các mệnh đề đã học qua nhiều lời gọi. Có hai cơ chế
chính.

\paragraph{Bổ sung mệnh đề đơn điệu.}
Công thức cơ sở \(F_0\) được mở rộng
\[
  F_0\subseteq F_1\subseteq\cdots\subseteq F_t.
\]
Đây là dạng tự nhiên cho tối ưu từ trên xuống: sau một mô hình có giá trị
\(\beta\), thêm mệnh đề buộc hàm mục tiêu tốt hơn.

\paragraph{Giả định.}
Một tập literal \(A_t\) chỉ có hiệu lực trong lần gọi \(t\). Các mệnh đề cơ sở
vẫn giữ nguyên; giả định cũ được thu hồi khi lời gọi kết thúc. Literal kích
hoạt \(a_B\) cho phép viết
\[
  a_B\rightarrow \operatorname{Bound}(B)
\]
và bật cận bằng giả định \(a_B\).

Mệnh đề học được hợp lệ khi bộ giải suy ra nó theo đúng quy trình giả định.
Công thức cơ sở, literal kích hoạt và thứ tự gọi là
một giao diện duy nhất.

\section{MaxSAT}

\emph{Partial MaxSAT} (MaxSAT bộ phận) chia công thức thành mệnh đề cứng \(H\) và
mệnh đề mềm \(S\). Mệnh đề cứng phải được thỏa; hàm mục tiêu tối thiểu số hoặc
tổng trọng số của các mệnh đề mềm bị vi phạm:
\[
  \min_{\alpha\models H}
  \sum_{C\in S}w_C[\alpha\not\models C].
\]
MaxSAT phù hợp khi mục tiêu có dạng tổng các quyết định Boolean, chẳng hạn số
thùng được dùng, số ràng buộc mềm bị vi phạm hoặc tổng mức phạt
\parencite{li2009maxsat,ansotegui2013maxsat}.

Một cách giải MaxSAT là chuyển hàm mục tiêu thành chuỗi bài toán \SAT với cận chi phí
\parencite{davies2011maxsat}. Vì thế khác biệt giữa SAT tăng dần và MaxSAT
thường nằm ở \emph{giao diện hàm mục tiêu}, còn phần biểu diễn tính khả thi có
thể được dùng chung.

Văn liệu MaxSAT hiện đại thường phân biệt ba kiến trúc. \emph{Linear/binary
search} (tìm kiếm tuyến tính/nhị phân) đặt cận lên tổng các literal nới lỏng.
\emph{Core-guided} (dẫn dắt bằng lõi) dùng một lõi bất khả thỏa để nhận diện
nhóm mệnh đề mềm không thể đồng thời thỏa, rồi nới chúng có kiểm soát.
\emph{Hitting-set} (tập đánh trúng) tách việc tìm lõi khỏi bài toán chọn các
mệnh đề mềm cần vi phạm.
Các bộ giải dạng mô-đun như Open-WBO và RC2 đặt phép biểu diễn đếm/PB ở trung tâm
thuật toán: mô-đun này được dựng,
mở rộng hoặc thay cận nhiều lần trong quá trình giải MaxSAT
\parencite{martins2014openwbo,ignatiev2019rc2}.

Ràng buộc đếm tăng dần giúp tránh xây lại hàm mục tiêu sau mỗi lõi hoặc cận
trên \parencite{martins2014incremental}. Vì thế khi chọn cây tổng, bộ đếm hay
mạng cho hàm mục tiêu, cần hỏi thêm: cấu trúc có hỗ trợ thêm đầu vào không; đầu
ra có đọc được ở nhiều cận không; và các mệnh đề đã học
có còn hợp lệ sau lần cập nhật không.

\begin{example}[Tối thiểu số tài nguyên]
Giả sử \(u_j=1\) khi tài nguyên \(j\) được sử dụng. Các ràng buộc gán bảo đảm mỗi tác
vụ chọn đúng một tài nguyên và
\[
  x_{ij}\rightarrow u_j.
\]
Cách SAT theo cận thêm
\[
  \sum_ju_j\le B.
\]
Cách MaxSAT giữ các mệnh đề cứng và đặt các mệnh đề mềm \(\neg u_j\) có trọng số
\(1\). Hai mô hình dùng chung các mô-đun gán và xung đột.
\end{example}

\section{Giải mã, ghi chứng minh và chứng nhận}

Mô hình thỏa của bộ giải chứa cả biến chính và biến phụ. Bộ giải mã chỉ đọc
biến chính để tạo nghiệm miền. Bộ kiểm tra sau đó xác minh:
\begin{enumerate}
  \item mọi giá trị nằm trong miền;
  \item mọi ràng buộc của bài toán gốc;
  \item hàm mục tiêu được tính lại từ nghiệm;
  \item cận và trạng thái bộ giải nhất quán.
\end{enumerate}
Bộ kiểm tra được cài đặt độc lập với bộ biểu diễn để hai thành phần tạo thành một phép
kiểm chéo.

Với giá trị tối ưu \(B^\star\), bằng chứng cận trên là một nghiệm có chi phí
\(B^\star\). Chứng cứ cận dưới có thể là chứng minh \UNSAT của
\(P(B^\star-1)\) hoặc một cận toán học bằng \(B^\star\). Mô hình SAT chứng
minh sự tồn tại; chứng minh \UNSAT loại mọi nghiệm ở cận tốt hơn.

Trong thực hành, bộ giải có thể xuất chứng minh theo một định dạng mệnh đề như
DRAT; bộ kiểm tra độc lập kiểm lại chuỗi phép thêm/xóa mệnh đề
\parencite{wetzler2014drat}. Điều này tạo một chuỗi tin cậy ngắn hơn so với
việc tin toàn bộ bộ giải, nhưng chỉ khi công thức được chứng minh đúng là công
thức của cận cần xét. Với giả định, ta vật hóa các giả định cuối thành mệnh đề
đơn vị, ghi chính xác mã kiểm tra của \CNF và kiểm chứng minh trên
\emph{artifact} (gói tái lập) đó.

\begin{center}
\captionof{table}{Vật chứng cần có cho các mức kết luận.}
\label{tab:evidence-levels}
\begin{tabularx}{\textwidth}{@{}>{\bfseries\sffamily}p{.22\textwidth}YY@{}}
\toprule
Kết luận & Vật chứng & Bộ kiểm tra độc lập\\
\midrule
Khả thi tại \(B\) & Mô hình/\emph{witness solution} (nghiệm làm chứng) &
Kiểm ràng buộc và chi phí miền\\
Vô nghiệm tại \(B\) & Chứng minh mệnh đề & Bộ kiểm chứng trên đúng \CNF\\
Tối ưu \(B^\star\) & Nghiệm tại \(B^\star\) và chứng nhận cận dưới &
Bộ kiểm nghiệm và bộ kiểm chứng cận\\
\bottomrule
\end{tabularx}
\end{center}

\section{Quy trình triển khai}

\begin{summarybox}
\begin{enumerate}
  \item Đọc bài toán đầu vào và chuẩn hóa miền.
  \item Tính cận dưới và cận trên.
  \item Sinh biến, ánh xạ và \CNF cơ sở.
  \item Chạy SAT/MaxSAT theo một quy trình tìm cận.
  \item Giải mã mô hình sau mỗi lần \SAT.
  \item Kiểm nghiệm và hàm mục tiêu bằng bộ kiểm tra độc lập.
  \item Lưu nghiệm cùng chứng nhận cận dưới ở cận cuối.
\end{enumerate}
\end{summarybox}

PySAT cung cấp môi trường thuận tiện để thử nghiệm nhiều bộ giải và phép biểu diễn
\parencite{ignatiev2018pysat}; ánh xạ biến và bộ kiểm tra vẫn thuộc về ứng
dụng. Một dây chuyền xử lý tốt có thể thay bộ giải mà giữ nguyên ngữ nghĩa.

\section{Bài học thiết kế}

\begin{enumerate}
  \item SAT là tầng quyết định; tối ưu cần quy trình tìm cận hoặc MaxSAT.
  \item Biến phụ phải có đặc tả ngữ nghĩa, không chỉ có một số DIMACS.
  \item Giải tăng dần hiệu quả khi công thức cơ sở ổn định và cận thay qua
        một giao diện rõ.
  \item Bộ giải mã và bộ kiểm tra là một phần của thuật toán chính xác.
  \item Chứng nhận tối ưu cần cả nghiệm và chứng nhận cận dưới.
\end{enumerate}

\subsection*{Câu hỏi nghiên cứu}
\begin{enumerate}
  \item Viết \(P(B)\) và quy tắc cập nhật cận cho một bài toán cực đại.
  \item Bỏ từng mệnh đề trong phép biểu diễn \(z\leftrightarrow(x\land y)\) và tìm
        phép gán làm mất một chiều tương đương.
  \item Mô tả hai cách biểu diễn ``tối thiểu số nhóm được dùng'': SAT theo cận và
        \emph{Partial MaxSAT} (MaxSAT bộ phận).
\end{enumerate}

\chapter{Thế nào là một phép mã hóa SAT tốt?}
\chapterlead{Chất lượng phép mã hóa được đánh giá trên nhiều trục. Đúng đắn là
điều kiện bắt buộc; các trục còn lại gồm kích thước, sức lan truyền, tính địa
phương, chi phí tích hợp và bằng chứng thực nghiệm.}

\index{encoding}\index{soundness}\index{completeness}
\index{unit propagation}\index{arc consistency}\index{channeling}

\flowdiagram{Tính đúng}{Kích thước}{Lan truyền}
  {Cấu trúc bộ giải}{Thực nghiệm}

\section{Phép mã hóa là một quan hệ có giao diện}

Cho ràng buộc \(\mathcal C\) trên tập biến chính \(X\), tập biến phụ \(Y\) và
\[
  F=\operatorname{Enc}(\mathcal C;X,Y).
\]
Tiêu chuẩn cơ bản là tương đương theo phép chiếu:
\[
  F\equiv_X\mathcal C
  \iff
  \forall\alpha_X:\quad
  \alpha_X\models\mathcal C
  \Longleftrightarrow
  \exists\beta_Y:\ \alpha_X\cup\beta_Y\models F.
\]
Hai chiều có tên riêng:
\begin{itemize}
  \item \textbf{soundness (tính đúng)}: mọi mô hình thỏa của \(F\) chiếu thành
        nghiệm của \(\mathcal C\);
  \item \textbf{completeness (tính đầy đủ)}: mọi nghiệm của \(\mathcal C\) mở
        rộng được thành mô hình thỏa của \(F\).
\end{itemize}

Tính đúng thường được chứng minh bằng cách lấy một mô hình \CNF và giải mã.
Tính đầy đủ thường được chứng minh bằng cách xây biến phụ từ một nghiệm miền.
Với bộ đếm, cách xây tự nhiên là đặt \emph{register} (thanh ghi ngưỡng, tức
biến phụ lưu một mức đếm) đúng theo tổng tiền tố thật.

\begin{proposition}[Mẫu chứng minh bằng phép chiếu]
Giả sử bộ giải mã \(D\) ánh mọi mô hình thỏa của \(F\) thành nghiệm của
\(\mathcal C\), và với mọi nghiệm \(s\) của \(\mathcal C\) tồn tại một phép gán biến phụ
\(\beta(s)\) làm \(F\) đúng. Khi đó \(F\equiv_X\mathcal C\).
\end{proposition}

\begin{proof}
Bộ giải mã cho tính đúng. Hàm \(\beta\) cho tính đầy đủ. Hai chiều đúng với
mọi phép gán trên \(X\), nên tương đương theo phép chiếu.
\end{proof}

\section{Hệ phân loại theo cách biểu diễn miền}

Trước khi so hai bộ mã hóa, cần xác định chúng có cùng ``từ vựng Boolean'' hay
không. Với biến miền hữu hạn \(v\in D=\set{d_1,\ldots,d_m}\), bốn họ thường
gặp là:
\begin{description}[style=nextline]
  \item[Direct encoding (mã hóa trực tiếp)]
  \(x_{v,a}\) mang nghĩa \(v=a\). Ràng buộc miền cần đúng-một; xung đột giữa
  hai giá trị được viết trực tiếp.

  \item[Support encoding (mã hóa hỗ trợ)]
  Mỗi lựa chọn \(v=a\) kéo theo một phép tuyển các giá trị hỗ trợ của biến bên
  kia. Cách này thường phản ánh nhất quán cung tốt nhưng các mệnh đề hỗ trợ có
  thể dài.

  \item[Order encoding (mã hóa thứ tự)]
  \(o_{v,t}\) mang nghĩa \(v\ge t\) hoặc \(v\le t\). Bất đẳng thức, khoảng và
  quan hệ trước--sau trở thành phép kéo theo giữa các ngưỡng
  \parencite{tamura2009csp}.

  \item[Log/binary encoding (mã hóa lôgarit/nhị phân)]
  Giá trị dùng \(\lceil\log_2m\rceil\) bit. Số biến nhỏ, nhưng một giá trị hoặc
  khoảng thường không còn là một literal đơn; lan truyền miền có thể đi qua
  mạch sâu hơn.
\end{description}

Mã hóa lai giữ hai góc nhìn, chẳng hạn biến giá trị để chọn một giá trị và biến
thứ tự để biểu diễn một khoảng. Chi phí liên kết hai biểu diễn chỉ đáng bỏ ra
khi góc nhìn thứ hai rút ngắn một họ ràng buộc lớn hoặc tạo suy luận mới. Đây
là kiến trúc được dùng lại trong Chương 7, 9--11.

\begin{center}
\captionof{table}{Các họ biểu diễn miền hữu hạn.}
\label{tab:domain-encodings}
\begin{tabularx}{\textwidth}{@{}>{\bfseries\sffamily}p{.18\textwidth}YYY@{}}
\toprule
Họ & Literal tự nhiên & Phù hợp & Điểm phải kiểm tra\\
\midrule
Trực tiếp & \(v=a\) & Bảng cấm, miền nhỏ & Đúng-một và nhiều cặp xung đột\\
Hỗ trợ & ``\(v=a\) có hỗ trợ'' & Ràng buộc nhị phân dạng bảng & Độ dài mệnh đề hỗ trợ\\
Thứ tự & \(v\ge t\) & Cận, khoảng, hiệu & Miền rộng và liên kết với biến giá trị\\
Lôgarit & Bit của giá trị & Miền lớn, mạch số học & Giá trị ngoài miền, sức lan truyền\\
Lai & Giá trị + ngưỡng/quan hệ & Nhiều kiểu ràng buộc & Chi phí và chiều liên kết\\
\bottomrule
\end{tabularx}
\end{center}

\begin{figure}[tb]
\centering
\begin{tikzpicture}[satfig]
  \node[satlabel,anchor=west] at (0,2.1) {Trực tiếp: một literal cho từng giá trị};
  \foreach \x/\v in {0/1,1.2/2,2.4/3,3.6/4}{
    \node[satbox,minimum width=10mm] at (\x,1.25) {\(v=\v\)};
  }
  \node[satlabel,anchor=west] at (0,.35) {Thứ tự: một chuỗi ngưỡng đơn điệu};
  \foreach \x/\v in {0/2,1.4/3,2.8/4}{
    \node[satstate] (o\v) at (\x,-.55) {\(\ge\v\)};
  }
  \draw[satedge] (o4)--(o3);
  \draw[satedge] (o3)--(o2);
  \draw[decorate,decoration={brace,mirror,amplitude=4pt},Ink]
    (-.5,-1.2)--(3.3,-1.2)
    node[midway,below=5pt,satlabel]{chuyển trạng thái xác định giá trị duy nhất};
\end{tikzpicture}
\caption{Hai từ vựng Boolean cho biến \(v\in\set{1,2,3,4}\).}
\label{fig:direct-order-domain}
\end{figure}

\begin{workedexample}{một bất đẳng thức trong mã hóa trực tiếp và thứ tự}
Cho \(u,v\in\set{0,\ldots,4}\) và \(u+2\le v\). Mã hóa trực tiếp cấm từng cặp
\((a,b)\) với \(a+2>b\):
\[
  \neg x_{u,a}\lor\neg x_{v,b}.
\]
Mã hóa thứ tự dùng \(o_{z,t}\leftrightarrow(z\ge t)\) và các mệnh đề
\[
  \neg o_{u,t}\lor o_{v,t+2},
  \qquad t=1,2.
\]
Nếu một ràng buộc khác suy \(u\ge2\), UP đặt \(o_{u,2}=1\), rồi suy
\(o_{v,4}=1\), tức \(v\ge4\). Một mệnh đề thứ tự đã thay toàn bộ các cặp trực tiếp
có \(u\ge2,v\le3\). Đổi lại, giá trị chính xác của \(v\) chỉ xuất hiện qua điểm chuyển
của chuỗi.
\end{workedexample}

\section{Contract: đặc tả ngữ nghĩa của biến phụ}

Một thanh ghi \(r_{i,j}\) có thể được định nghĩa theo ba mức:
\begin{center}
\captionof{table}{Ba mức đặc tả của literal ngưỡng.}
\label{tab:threshold-contracts}
\begin{tabularx}{\textwidth}{@{}>{\bfseries\sffamily}p{.22\textwidth}YY@{}}
\toprule
Đặc tả & Công thức & Cách sử dụng\\
\midrule
Điều kiện đủ & \(r_{i,j}\rightarrow S_i\ge j\) &
\(r=1\) chứng minh đạt ngưỡng\\
Điều kiện cần & \(S_i\ge j\rightarrow r_{i,j}\) &
\(r=0\) chứng minh chưa đạt ngưỡng\\
Tương đương & \(r_{i,j}\leftrightarrow(S_i\ge j)\) &
Cả hai cực tính đều có ngữ nghĩa\\
\bottomrule
\end{tabularx}
\end{center}

Một \emph{connector} (bộ nối giữa các mô-đun) dùng \(\neg r_{i,j}\) như mệnh
đề ``tổng dưới \(j\)'' phải có chiều cần. Nếu bộ đếm chỉ cung cấp chiều đủ,
bộ nối đó không đúng dù khi đứng riêng hai công thức vẫn có thể
\emph{equisatisfiable} (đồng khả thỏa, tức cùng có nghiệm hoặc cùng vô nghiệm).

\begin{designrule}{Quy tắc ghép mô-đun}
Trước khi nối hai mô-đun, đối chiếu cực tính mà bên sử dụng cần với chiều đặc
tả mà bên cung cấp bảo đảm. Nếu hai phía cần đọc cùng một biến theo hai cách,
phải có mệnh đề liên kết tương ứng.
\end{designrule}

\begin{figure}[tb]
\centering
\begin{tikzpicture}[satfig,node distance=16mm]
  \node[satbox] (p) {Bộ đếm\\cung cấp};
  \node[satstate,right=of p] (r) {\(r_{i,j}\)};
  \node[satbox,right=of r] (c1) {Mô-đun A\\đọc \(r\)};
  \node[satbox,below=11mm of c1] (c2) {Mô-đun B\\đọc \(\neg r\)};
  \draw[satedge] (p)--node[above,satlabel]{``tổng \(\ge j\)''}(r);
  \draw[satedge] (r)--(c1);
  \draw[satedge] (r)--(c2);
  \node[satwarn,rounded corners=2pt,below=11mm of r] (w)
    {Mô-đun B cần chiều\\
     \(\displaystyle S_i<j\Rightarrow\neg r_{i,j}\)};
  \draw[satdash] (w)--(c2);
\end{tikzpicture}
\caption{Đặc tả của literal phụ quyết định cách các mô-đun khác được sử dụng nó.}
\label{fig:provider-client-contract}
\end{figure}

\section{Đo kích thước}

Bốn đại lượng nền tảng là:
\[
  n_v=\card{Y},\qquad
  n_c=\card{F},\qquad
  n_\ell=\sum_{C\in F}\card C,\qquad
  w_{\max}=\max_{C\in F}\card C.
\]
Ngoài độ lớn tiệm cận, cần xem phân bố mệnh đề nhị phân/tam phân, số lần xuất
hiện của biến và bộ nhớ của bộ sinh. Với kỹ thuật hai literal được theo dõi,
tổng số literal thường gần chi phí bộ nhớ hơn số mệnh đề.

Một công thức đóng về \(n_v,n_c\) phải chỉ rõ:
\begin{itemize}
  \item miền tham số, đặc biệt \(k=0,1,n\);
  \item thanh ghi biên nào được bỏ;
  \item khối cuối có thể ngắn;
  \item mệnh đề lôgic nào cần Tseitin hóa thêm.
\end{itemize}

Bộ sinh nên tự xuất thống kê và so với công thức lý thuyết trên các bài toán
nhỏ. Đây là một kiểm thử đơn vị cho phép đếm, đồng thời phát hiện mệnh đề bị lặp hoặc
thanh ghi vô dụng.

\section{Unit propagation: lan truyền đơn vị}

Cho phép gán bộ phận \(A\). \emph{Unit propagation} (lan truyền đơn vị, UP)
lặp việc thỏa một mệnh đề
đơn vị cho
đến điểm cố định hoặc xung đột. Có ba mức phát biểu thường gặp:

\begin{description}[style=nextline]
  \item[UP phát hiện vi phạm]
  Nếu \(A\) không mở rộng được thành nghiệm của \(\mathcal C\), UP trên
  \(F\land A\) tạo xung đột.

  \item[Unit-refutation completeness (URC)]
  Mọi phép gán làm \(F\) không thỏa được đều bị UP chứng minh bất khả thỏa
  trên toàn bộ biến; sách gọi tính chất này là \emph{đầy đủ bác bỏ bằng lan
  truyền đơn vị}.

  \item[Propagation completeness (PC)]
  Mọi literal bị \(F\land A\) kéo theo trong phạm vi quan sát đều được UP suy
  ra, hoặc UP chứng minh bất khả thỏa; đây là \emph{đầy đủ lan truyền}.
\end{description}

Trong một phép mã hóa miền hữu hạn, \emph{generalized arc consistency}
(nhất quán cung tổng quát, GAC) loại mọi giá trị không còn giá trị hỗ trợ.
Phát biểu GAC luôn đi cùng cách Boolean hóa miền và tập literal đại diện cho
giá trị.

Kết quả của Gent và Bacchus cung cấp một cầu nối chính xác: nếu literal Boolean
đại diện trực tiếp cho việc một giá trị còn trong miền, có thể so sánh bao đóng
của lan truyền đơn vị với bao đóng của một bộ lan truyền GAC
\parencite{gent2002arc,bacchus2007gac}. Khi thêm liên kết biểu diễn, đổi từ mã
hóa trực tiếp sang thứ tự hoặc chỉ giữ một chiều Tseitin, tập literal quan sát
thay đổi; phát biểu về GAC vì thế cũng được xác định lại theo tập này.

\begin{example}[AMO từng cặp]
Với
\[
  \operatorname{AMO}(x_1,\ldots,x_n)=
  \bigwedge_{i<j}(\neg x_i\lor\neg x_j),
\]
nếu \(x_i=1\), UP lập tức suy \(\neg x_j\) cho mọi \(j\ne i\). Nếu hai biến
đúng, UP tạo xung đột. AMO từng cặp vì thế có lan truyền rất trực tiếp trên
biến chính, đổi lại \(O(n^2)\) mệnh đề.
\end{example}

\section{Cận dưới kích thước và mặt Pareto}

Yêu cầu lan truyền đặt ra cận dưới cho kích thước phép mã hóa. Với AMO, các
cận dưới đã biết cho thấy một phép mã hóa đầy đủ lan truyền cần số mệnh đề
tuyến tính với hệ số bị chặn, kể cả khi dùng biến phụ
\parencite{kucera2019lower}. Kết quả này dẫn đến ba nguyên tắc:
\begin{enumerate}
  \item xác định rõ lớp mã hóa khi phát biểu về kích thước, biến phụ và sức lan
        truyền;
  \item một cải tiến vài mệnh đề có thể đáng kể nếu nó tiến gần cận dưới;
  \item so sánh kích thước phải giữ cùng tiêu chuẩn lan truyền và cùng ngữ
        nghĩa của đầu ra.
\end{enumerate}

Các kết quả về phép phân rã \textsc{AllDifferent} và ràng buộc đếm toàn cục
còn liên hệ yêu cầu lan truyền với cận biểu diễn của mạch đơn điệu
\parencite{bessiere2009decompositions}. Khi một ràng buộc toàn cục có
bộ lan truyền đạt mức nhất quán mạnh hơn phép mã hóa đa thức đang xét, phép so
sánh tập trung vào phần suy luận được chuyển sang học mệnh đề và chi phí
để đạt phần suy luận đó.

\begin{keyidea}{Mặt Pareto của phép mã hóa}
Một bộ mã hóa được định vị trên bốn trục: kích thước, sức lan truyền, chi phí
xây dựng và khả năng tăng dần. Điểm Pareto là điểm mà cải thiện thêm một trục
sẽ đánh đổi ít nhất một trục khác. \emph{Benchmark} (bộ kiểm chuẩn) đo vị trí
thực nghiệm của từng thiết kế trên mặt Pareto này.
\end{keyidea}

\section{Mức độ phù hợp với bộ giải}

Hai phép mã hóa có cùng tính đúng và kích thước vẫn có thể khác về:
\begin{itemize}
  \item \textbf{tính địa phương}: một ràng buộc cục bộ dùng các biến gần nhau
        hay đi qua một cấu trúc toàn cục;
  \item \textbf{độ sâu kéo theo}: số bước từ quyết định đến xung đột;
  \item \textbf{khả năng học}: xung đột có tạo mệnh đề tái sử dụng được không;
  \item \textbf{đối xứng}: trạng thái phụ có tạo nhiều mô hình tương đương;
  \item \textbf{khả năng tăng dần}: đầu ra có đọc được ở nhiều cận không.
\end{itemize}

Không có một công thức giải tích đơn giản cho các thuộc tính này. Ta dùng
đại lượng thay thế như số mệnh đề nhị phân, đồ thị kéo theo trên bài toán nhỏ,
hoặc số lần lan truyền/xung đột; kết luận cuối cùng cần bộ kiểm chuẩn có kiểm
soát.

\section{Tính mô-đun và mã hóa lại có cấu trúc}

Một mô-đun mã hóa tốt có bốn thành phần:
\[
  \texttt{variables}
  +\texttt{clauses}
  +\texttt{decoder}
  +\texttt{contract}.
\]
Khi nhiều ràng buộc chia sẻ biểu thức con, mã hóa lại có cấu trúc tạo biến phụ
cho phần chung thay vì Tseitin hóa từng công thức riêng
\parencite{haberlandt2023aux}. Lợi ích xuất hiện khi mô-đun thực sự được dùng
chung trong \CNF và bên sử dụng đọc đúng đặc tả.

Ví dụ, nhiều cửa sổ cùng đọc ngưỡng của một tiền tố. Một bộ đếm tiền tố dùng
chung giảm lặp; bộ nối chỉ ghép các ngưỡng ở hai biên. Đây là chủ đề trung tâm
của Chương 5.

\section{Ma trận phát biểu--bằng chứng}

\begin{center}
\captionof{table}{Ma trận phát biểu--bằng chứng cho một phép mã hóa SAT.}
\label{tab:claim-evidence}
\begin{tabularx}{\textwidth}{@{}>{\bfseries\sffamily}p{.24\textwidth}YY@{}}
\toprule
Phát biểu & Bằng chứng lý thuyết & Bằng chứng thực nghiệm\\
\midrule
Mã hóa đúng & Chứng minh hai chiều & Kiểm thử vét cạn/vi sai trên miền nhỏ\\
Mã hóa nhỏ hơn & Công thức đếm & Kích thước \CNF thực tế\\
Lan truyền mạnh hơn & Bổ đề UP/PC/GAC & Bộ đếm lan truyền/xung đột\\
Nhanh hơn & Phân tích cơ chế & Bộ kiểm chuẩn, nhiều \emph{random seed}
(hạt giống ngẫu nhiên), \emph{performance profile} (hồ sơ hiệu năng)\\
Chứng nhận tối ưu & Định lý về cận & Nghiệm và chứng minh ở hai cận kề\\
\bottomrule
\end{tabularx}
\end{center}

Ma trận đưa mỗi ý tưởng vào đúng loại đóng góp: định lý, thiết kế cấu trúc hay
quan sát thực nghiệm.

\section{Một quy trình đánh giá phép mã hóa}

\begin{summarybox}
\begin{enumerate}
  \item Viết rõ ngữ nghĩa của biến chính và biến phụ.
  \item Chứng minh tính đúng và tính đầy đủ theo phép chiếu.
  \item Đếm biến, mệnh đề, literal và độ rộng mệnh đề theo đúng biến thể.
  \item Phân tích một số suy luận cốt lõi của lan truyền đơn vị (UP).
  \item Xây dựng \emph{oracle} (bộ đối chiếu nhỏ trên miền đã biết) để so sánh
        tập nghiệm sau phép chiếu.
  \item Đánh giá trên bộ kiểm chuẩn theo quy trình của Chương 4.
  \item Kiểm tra khả năng kết hợp mô-đun và giải tăng dần.
\end{enumerate}
\end{summarybox}

\section{Bài học thiết kế}

\begin{enumerate}
  \item Tính đúng là quan hệ theo phép chiếu, không phải số mô hình của biến phụ.
  \item Đặc tả ngữ nghĩa quyết định biến phụ có thể được tái sử dụng thế nào.
  \item Đếm literal và cấu trúc mệnh đề quan trọng ngang với số mệnh đề.
  \item Phát biểu về sức lan truyền phải nêu rõ phạm vi literal được quan sát.
  \item Thời gian chạy là kết quả của tương tác giữa phép mã hóa, bộ giải và bài toán thử.
\end{enumerate}

\subsection*{Câu hỏi nghiên cứu}
\begin{enumerate}
  \item Cho một phép mã hóa AMO có biến phụ, viết bộ giải mã và chứng
        minh theo phép chiếu.
  \item Thiết kế một phép thử phân biệt đặc tả điều kiện đủ với đặc tả tương
        đương của thanh ghi ngưỡng.
  \item Chọn ba thước đo để giải thích vì sao một phép mã hóa lớn hơn lại chạy
        nhanh hơn trên các bài toán thử \UNSAT.
\end{enumerate}

\chapter{Bộ công cụ ràng buộc đếm và giả Boolean}
\chapterlead{\emph{Cardinality constraint} (ràng buộc đếm) giới hạn số literal
đúng trong một tập; \emph{pseudo-Boolean constraint} (ràng buộc giả Boolean)
mở rộng mô hình này
bằng các trọng số nguyên. Hai lớp này xuất hiện trong ràng buộc gán, sức chứa,
chuỗi, đối xứng và hàm mục tiêu. Không có một bộ biểu diễn tốt nhất cho mọi miền;
cần chọn theo \(n\), ngưỡng \(k\), các đầu ra cần đọc và cách bộ giải thay cận.}

\index{cardinality constraint}\index{at-most-one}
\index{sequential counter}\index{totalizer}\index{sorting network}
\index{binary decision diagram}\index{pseudo-Boolean constraint}

\flowdiagram{AMO/EXO}{AMK/ALK}{Đúng \(K\)}
  {Giả Boolean}{Chọn bộ biểu diễn}

\section{Bản đồ phát triển của phép biểu diễn ràng buộc đếm}

Các phép biểu diễn ràng buộc đếm không phát triển theo một chuỗi trong đó bộ mã
hóa mới thay thế hoàn toàn bộ biểu diễn cũ. Mỗi họ giải quyết một kiểu đánh đổi:
\begin{enumerate}
  \item \emph{pairwise encoding} (biểu diễn từng cặp) và các phép phân rã AMO ưu tiên
        phát hiện xung đột trực tiếp;
  \item \emph{Sequential Counter} (bộ đếm tuần tự) hiện thực hóa (reification) tổng trên từng tiền tố;
  \item \emph{Totalizer} (cây tổng) hiện thực hóa (reification) tổng theo cấu trúc cây;
  \item \emph{Sorting Network / Selection Network} (mạng sắp xếp / lựa chọn) ưu tiên cấu
        trúc đều và độ sâu nhỏ;
  \item \emph{binary decision diagram} (biểu đồ quyết định nhị phân, BDD)
        hợp nhất các trạng thái tổng tương đương;
  \item Totalizer tổng quát và các biến thể PB xử lý trọng số;
  \item danh mục phương pháp hoặc bộ chọn đã học chọn phép biểu diễn theo đặc
        trưng của bài toán thử.
\end{enumerate}

Các khảo cứu về ràng buộc đếm nhấn mạnh rằng phải tách ba thuộc tính: kích
thước tiệm cận, mức nhất quán do lan truyền đơn vị duy trì và hiệu năng trong
một bộ giải cụ thể \parencite{bailleux2010survey}. Một Sequential Counter
\(O(nk)\) có thể nhỏ khi \(k\) bé; một mạng đếm có thể lớn hơn nhưng có độ sâu
tốt; Totalizer lại thuận lợi khi cận thay đổi. Không có quan hệ thứ tự tuyến
tính giữa ba lựa chọn đó.

\section{Các dạng ràng buộc}

Với \(X=\set{x_1,\ldots,x_n}\), ba tên tiếng Anh lần lượt là
\emph{at-most-\(k\)} (nhiều nhất \(k\), AMK), \emph{at-least-\(k\)}
(ít nhất \(k\), ALK) và \emph{exactly-\(k\)} (đúng \(k\), EXK):
\[
  \operatorname{AMK}(X,k):\quad \sum_{i=1}^{n}x_i\le k,
\]
\[
  \operatorname{ALK}(X,k):\quad \sum_{i=1}^{n}x_i\ge k,
\]
\[
  \operatorname{EXK}(X,k):\quad \sum_{i=1}^{n}x_i=k.
\]
\emph{At-most-one} (nhiều nhất một, AMO) là AMK với \(k=1\);
\emph{at-least-one} (ít nhất một, ALO) là một mệnh đề dài
\(x_1\lor\cdots\lor x_n\); \emph{exactly-one} (đúng một, EXO) là AMO kết hợp ALO.

Đối ngẫu hữu ích:
\[
  \sum_i x_i\ge k
  \iff
  \sum_i(1-x_i)\le n-k.
\]
Tuy nhiên, đối ngẫu có thể đổi ngưỡng từ nhỏ sang lớn. Vì vậy, lựa chọn giữa
Direct Encoding ALK và biểu diễn AMK trên các literal phủ định phụ thuộc vào
\(k\) và những đầu ra đã có.

\section{At-most-one (AMO): ràng buộc nhiều nhất một}

\subsection{Biểu diễn từng cặp}

Biểu diễn từng cặp sinh
\[
  \neg x_i\lor\neg x_j
  \qquad\forall i<j.
\]
Phương pháp này không dùng biến phụ, có \(\binom n2\) mệnh đề nhị phân và lan truyền trực
tiếp. Đây là lựa chọn tốt khi miền nhỏ.

\subsection{Biểu diễn tuần tự}

Đặt \(s_i=1\) khi ít nhất một trong \(x_1,\ldots,x_i\) đúng. Một biến thể
biểu diễn AMO tuần tự dùng:
\[
  \neg x_i\lor s_i,
  \qquad
  \neg s_{i-1}\lor s_i,
  \qquad
  \neg x_i\lor\neg s_{i-1}.
\]
Mệnh đề cuối cấm \(x_i\) khi tiền tố trước đó đã có một giá trị đúng. Kích thước
tuyến tính theo \(n\) \parencite{sinz2005cardinality}.

\subsection{Biểu diễn tích và biểu diễn chỉ huy}

\emph{Product encoding} (biểu diễn tích) đặt biến vào ma trận rồi dùng AMO theo
hàng và cột. \emph{Commander encoding} (biểu diễn chỉ huy) chia biến thành nhóm,
tạo một biến đại diện cho mỗi nhóm rồi biểu diễn AMO trên các đại diện. Hai cách
này tạo tính cục bộ theo khối và phù hợp với miền vừa hoặc lớn
\parencite{chen2010amo}.

\begin{center}
\captionof{table}{So sánh các bộ biểu diễn AMO cơ bản.}
\label{tab:amo-encoders}
\begin{tabularx}{\textwidth}{@{}>{\bfseries\sffamily}p{.19\textwidth}YYY@{}}
\toprule
Phép biểu diễn AMO & Biến phụ & Mệnh đề & Điểm mạnh\\
\midrule
Từng cặp & \(0\) & \(O(n^2)\) & Nhị phân, đơn giản\\
Tuần tự & \(O(n)\) & \(O(n)\) & Lan truyền theo tiền tố\\
Tích & \(O(\sqrt n)\) & \(O(n)\) & Cân bằng hai chiều\\
Chỉ huy & \(O(n/g)\) & Phụ thuộc nhóm \(g\) & Cục bộ, dễ kết hợp\\
\bottomrule
\end{tabularx}
\end{center}

\section{Sequential Counter cho AMK}

Đặt
\[
  s_{i,j}=1
  \iff
  \sum_{q=1}^{i}x_q\ge j,
  \qquad 1\le j\le k.
\]
Bộ đếm cập nhật theo hệ thức truy hồi
\[
  s_{i,j}
  \leftrightarrow
  s_{i-1,j}\lor(x_i\land s_{i-1,j-1}),
\]
với \(s_{i,0}=\top\). Cận AMK cấm trạng thái tràn:
\[
  \neg x_i\lor\neg s_{i-1,k}.
\]
Số thanh ghi và mệnh đề là \(O(nk)\). Với \(k\) nhỏ, bộ đếm tạo các quan hệ
kéo theo cục bộ và cho phép đọc nhiều ngưỡng tiền tố
\parencite{sinz2005cardinality}.

\begin{proposition}[Ngữ nghĩa của bộ đếm]
Nếu các hệ thức truy hồi được biểu diễn hai chiều thì, với mọi \(i,j\),
\[
  s_{i,j}=1\iff\sum_{q=1}^{i}x_q\ge j.
\]
\end{proposition}

\begin{proof}
Quy nạp theo \(i\). Tiền tố \(i\) đạt \(j\) khi tiền tố \(i-1\) đã đạt \(j\),
hoặc \(x_i=1\) và tiền tố trước đạt \(j-1\). Đây chính là hệ thức truy hồi.
\end{proof}

\begin{figure}[tb]
\centering
\begin{tikzpicture}[satfig,x=16mm,y=10mm]
  \foreach \i in {1,...,5}
    \node[satlabel] at (\i,2.7) {\(x_\i\)};
  \foreach \i in {1,...,5}{
    \foreach \j in {1,2}{
      \node[satstate] (s\i\j) at (\i,2-\j) {\(\scriptstyle s_{\i,\j}\)};
    }
  }
  \foreach \i in {2,...,5}{
    \draw[satedge] (s\the\numexpr\i-1\relax1)--(s\i1);
    \draw[satedge] (s\the\numexpr\i-1\relax2)--(s\i2);
  }
  \foreach \i in {2,...,5}
    \draw[satdash] (\i,2.35)--(s\i1);
  \foreach \i in {2,...,5}
    \draw[satdash] (\i,2.35) to[bend right=18] (s\i2);
  \node[satwarn,rounded corners=2pt,right=10mm of s52]
    {\(\neg s_{5,3}\)\\(trạng thái tràn bị cắt)};
\end{tikzpicture}
\caption{Lưới trạng thái của Sequential Counter cho \(n=5,k=2\). Mũi tên liền
truyền ngưỡng; mũi tên đứt đưa đầu vào mới vào hệ thức truy hồi.}
\label{fig:seq-counter-grid}
\end{figure}

\begin{workedexample}{AMK với năm đầu vào và cận hai}
Xét \(x_1+\cdots+x_5\le2\). Sau khi \(x_1=x_3=1\), bộ đếm suy ra
\[
  s_{3,2}=1.
\]
Mệnh đề chống tràn tại \(i=4\),
\[
  \neg x_4\lor\neg s_{3,2},
\]
trở thành mệnh đề đơn vị và suy ra \(\neg x_4\). Tương tự, tính đơn điệu của
tiền tố đưa \(s_{3,2}\) đến \(s_{4,2}\), rồi mệnh đề chống tràn tại \(i=5\)
suy ra \(\neg x_5\). Bộ đếm không cần đợi đầu vào đúng thứ ba mới báo xung đột; cấu trúc sẽ
loại mọi vị trí còn lại ngay khi cận đã đầy.
\end{workedexample}

\section{Totalizer}

Totalizer xây dựng một cây nhị phân. Mỗi nút trả về một vectơ đơn phân
\[
  o_j=1\iff \text{tổng trong cây con}\ge j.
\]
Khi ghép hai nút con có đầu ra \(a,b\), nút cha \(r\) thỏa
\[
  a_p\land b_q\rightarrow r_{p+q}.
\]
Chiều ngược bảo đảm đặc tả đầy đủ nếu các đầu ra được dùng bên ngoài. Cận
\(\sum x_i\le k\) chỉ cần \(\neg o_{k+1}\).

Totalizer là lựa chọn tự nhiên khi:
\begin{itemize}
  \item cần nhiều ngưỡng trên cùng một tập biến;
  \item cận được thay đổi tăng dần;
  \item cần ràng buộc đúng \(K\) bằng \(o_k\land\neg o_{k+1}\).
\end{itemize}
Totalizer cân bằng tạo tính cục bộ theo cây; kích thước có thể được cắt ở
\(k+1\) thay vì giữ toàn bộ tổng \parencite{bailleux2003cardinality}.

\begin{figure}[tb]
\centering
\begin{tikzpicture}[satfig,level distance=13mm,
  level 1/.style={sibling distance=45mm},
  level 2/.style={sibling distance=22mm}]
  \node[satstate] (r) {\(r_{1:4}\)}
    child {node[satstate] (a) {\(a_{1:2}\)}
      child {node[satbox] {\(x_1\)}}
      child {node[satbox] {\(x_2\)}}}
    child {node[satstate] (b) {\(b_{1:2}\)}
      child {node[satbox] {\(x_3\)}}
      child {node[satbox] {\(x_4\)}}};
  \node[satlabel,right=11mm of r] {\(r_j=1\iff\sum_i x_i\ge j\)};
\end{tikzpicture}
\caption{Totalizer cân bằng: mỗi nút trả về các ngưỡng của tổng trong cây con.}
\label{fig:totalizer-tree}
\end{figure}

\section{Mạng đếm}

Sorting Network sắp các đầu vào Boolean thành các đầu ra đơn phân:
\[
  o_1\ge o_2\ge\cdots\ge o_n,
  \qquad
  o_j=1\iff\sum_i x_i\ge j.
\]
Bộ so sánh
\[
  (a,b)\mapsto(a\lor b,\ a\land b)
\]
được chuyển đổi theo phương pháp Tseitin. Cận AMK là \(\neg o_{k+1}\). Mạng có cấu trúc đều, độ sâu
đa lôgarit và thuận lợi cho lan truyền song song
\parencite{batcher1968sorting,asin2011networks}.

So với Sequential Counter, mạng dùng nhiều mệnh đề hơn khi \(k\) nhỏ nhưng có độ
sâu suy diễn thấp hơn và các đầu ra toàn cục rõ ràng.

\begin{figure}[tb]
\centering
\begin{tikzpicture}[satfig,x=17mm,y=9mm]
  \node[satbox] (a) at (0,1) {\(a\)};
  \node[satbox] (b) at (0,0) {\(b\)};
  \node[satstate] (c) at (2,1) {\(\lor\)};
  \node[sataux] (d) at (2,0) {\(\land\)};
  \node[satbox] (hi) at (4,1) {\(\max(a,b)\)};
  \node[satbox] (lo) at (4,0) {\(\min(a,b)\)};
  \draw[satedge] (a)--(c);
  \draw[satedge] (b)--(c);
  \draw[satedge] (a)--(d);
  \draw[satedge] (b)--(d);
  \draw[satedge] (c)--(hi);
  \draw[satedge] (d)--(lo);
\end{tikzpicture}
\caption{Bộ so sánh Boolean là thành phần cơ sở của Sorting Network và mạng đếm.}
\label{fig:boolean-comparator}
\end{figure}

\section{Biểu diễn bằng BDD}

\emph{Binary decision diagram} (biểu đồ quyết định nhị phân, BDD) biểu diễn
trạng thái tổng tích lũy. Nút ở mức \(i\) và trạng thái \(s\) phân nhánh theo
\(x_i\). Các trạng thái vượt \(k\) đi vào nút kết thúc sai. Một biến phụ cho
mỗi nút BDD cùng các mệnh đề liên kết theo hai cung tạo thành phép biểu diễn
\parencite{bryant1986bdd,abio2012bdd}.

Ưu điểm của BDD:
\begin{itemize}
  \item xử lý tự nhiên các ràng buộc giả Boolean có trọng số;
  \item có thể đạt sức lan truyền mạnh;
  \item hợp nhất các trạng thái tương đương.
\end{itemize}
Nhược điểm là kích thước phụ thuộc thứ tự biến và số trạng thái tổng khác nhau.

\section{Ràng buộc đúng \(K\) và đầu ra chia sẻ}

EXK có thể viết thành AMK và ALK, nhưng dựng hai bộ đếm độc lập sẽ lặp trạng
thái. Nếu bộ biểu diễn trả về các đầu ra đơn phân \(o_j\), chỉ cần
\[
  o_k\land\neg o_{k+1}.
\]
Sequential Counter cũng có thể chia sẻ thanh ghi: biên \(s_{n,k}\) bảo đảm ít
nhất \(k\), còn các mệnh đề chống tràn bảo đảm nhiều nhất \(k\).

\begin{keyidea}{Định nghĩa một lần, đọc nhiều cận}
Tách \texttt{CounterDefinition} khỏi \texttt{Bound}. Cùng một cấu trúc đầu ra
có thể phục vụ AMK, ALK, EXK và vòng giải tăng dần. Đây là bước đầu của tư duy
mô-đun được phát triển ở Phần II.
\end{keyidea}

\section{Pseudo-Boolean có trọng số}

Ràng buộc tổng quát
\[
  \sum_{i=1}^{n}a_ix_i\le B,
  \qquad a_i\in\mathbb Z_{\ge0},
\]
không còn là cardinality khi trọng số khác nhau. Các hướng chính gồm:
\begin{itemize}
  \item biểu diễn bằng bộ cộng hoặc mạch cho biểu diễn nhị phân của tổng;
  \item BDD cho các trạng thái tổng có thể đạt;
  \item Totalizer tổng quát gom và cắt các tổng;
  \item phân rã bằng Sorting Network hoặc mạng đếm sau khi tách trọng số;
  \item thư viện kết hợp như PBLib chọn bộ biểu diễn theo cấu trúc
        \parencite{een2006pb,philipp2015pblib}.
\end{itemize}

Nếu hàm mục tiêu thay cận nhiều lần, các đầu ra hoặc giao diện kích hoạt quan
trọng hơn vài phần trăm kích thước ở một cận duy nhất.

\subsection{Totalizer tổng quát}

\emph{Generalized totalizer} (Totalizer tổng quát) thay vectơ đơn phân liên tiếp
bằng tập các tổng có thể đạt tại mỗi nút. Khi ghép hai nút có tập tổng \(A,B\),
nút cha giữ các giá trị
\[
  \min(a+b,B+1),\qquad a\in A,\ b\in B,
\]
trong đó \(B+1\) gom mọi trạng thái tràn. Phép biểu diễn đặc biệt hiệu quả khi
trọng số có cấu trúc và số tổng phân biệt sau phép cắt nhỏ hơn tổng số học thô
\parencite{joshi2015gte}. Hình dạng cây không còn là chi tiết vô hại: hai cách
ghép các lá có thể tạo số trạng thái trung gian rất khác nhau.

\subsection{Khai thác các nhóm AMO trong PB}

Trong nhiều mô hình, các hạng tử của ràng buộc PB đến từ một biến miền hữu hạn:
trong mỗi nhóm, nhiều nhất một literal có thể đúng. Nếu trải phẳng mọi hạng tử
như các biến độc lập, bộ biểu diễn sẽ bỏ mất thông tin này. Các phép biểu diễn
PB(AMO) đưa nhóm AMO trực tiếp vào \emph{multi-valued decision diagram}
(biểu đồ quyết định đa trị, MDD), Totalizer
tổng quát, Sequential Counter có trọng số hoặc Totalizer theo môđun
\parencite{bofill2022pbamo}. Hai lợi ích có thể xuất hiện:
\begin{itemize}
  \item ít trạng thái tổng hơn, vì tổ hợp chọn hai hạng tử cùng nhóm không bao
        giờ đạt được;
  \item lan truyền dùng được cả cận PB lẫn tính loại trừ lẫn nhau trong nhóm.
\end{itemize}

Đây là một ví dụ tổng quát của \emph{structure-aware encoding}
(phép biểu diễn có xét cấu trúc): trước khi chọn bộ biểu diễn cho một biểu thức
tuyến tính, cần giữ nguồn gốc cấu trúc của các hạng tử thay vì chỉ truyền vào
mảng hệ số.

\section{Từ lựa chọn thủ công đến lựa chọn theo bài toán thử}

Quy tắc ``\(k\) nhỏ thì dùng Sequential Counter, nhiều cận thì dùng Totalizer'' là
một đối chứng hữu ích, nhưng chưa bao quát tương tác giữa hệ số, nhóm AMO, độ
rộng miền và bộ giải. Nghiên cứu về \emph{algorithm selection}
(lựa chọn thuật toán) đã dùng các đặc trưng của ràng buộc và bài toán
thử để chọn giữa nhiều phép biểu diễn SAT, kể cả trên những lớp bài toán chưa xuất
hiện trong dữ liệu huấn luyện \parencite{ulrich2022selection}. Kết quả này thay
đổi cách xây dựng thư viện bộ biểu diễn: mỗi mô-đun cần công bố đủ siêu dữ liệu để
một bộ chọn ra quyết định.

\begin{center}
\captionof{table}{Nhóm đặc trưng để lựa chọn bộ biểu diễn theo bài toán thử.}
\label{tab:selector-features}
\begin{tabularx}{\textwidth}{@{}>{\bfseries\sffamily}p{.22\textwidth}YY@{}}
\toprule
Nhóm đặc trưng & Ví dụ & Tác động tiềm năng\\
\midrule
Kích thước & \(n,k,B\), số hạng tử & Loại bộ biểu diễn chắc chắn quá lớn\\
Trọng số & Giá trị lớn nhất, ước chung lớn nhất, số tổng phân biệt & Không gian trạng thái BDD/GT\\
Cấu trúc & Nhóm AMO, hệ số lặp & Khả năng nén trạng thái\\
Giao diện & Số cận, thêm đầu vào, hiện thực hóa (reification) & Giải tăng dần và kết hợp mô-đun\\
Mô hình & Mật độ, đối xứng, đồ thị ràng buộc & Tương tác với học mệnh đề\\
\bottomrule
\end{tabularx}
\end{center}

Bộ chọn không loại bỏ nhu cầu chứng minh tính đúng: công cụ chỉ lựa chọn giữa các
mô-đun đã được kiểm chứng. Việc đánh giá cũng phải tách các họ dùng để huấn
luyện khỏi các họ dùng để kiểm thử; nếu không, bộ chọn có thể chỉ nhận dạng tên
bộ kiểm chuẩn thay vì học cấu trúc của phép biểu diễn.

\section{Bản đồ lựa chọn}

\begin{center}
\captionof{table}{Gợi ý lựa chọn bộ biểu diễn ràng buộc đếm/PB.}
\label{tab:cardinality-choice}
\begin{tabularx}{\textwidth}{@{}>{\bfseries\sffamily}p{.26\textwidth}YY@{}}
\toprule
Tình huống & Lựa chọn khởi đầu & Lý do\\
\midrule
AMO, \(n\) nhỏ & Từng cặp & Không biến phụ, mệnh đề nhị phân\\
AMO/AMK, \(k\) nhỏ & Tuần tự & \(O(nk)\), lan truyền theo tiền tố\\
Nhiều ngưỡng & Totalizer/mạng đếm & Tái sử dụng đầu ra\\
Đúng \(K\) & Một vectơ đầu ra chung & Không lặp AMK/ALK\\
PB có trọng số, ít trạng thái & BDD & Hợp nhất trạng thái tổng\\
PB có trọng số tổng quát & PBLib/kết hợp & Chọn theo bài toán thử\\
PB có nhóm AMO & Bộ biểu diễn PB(AMO) & Không hiện thực hóa (reification) trạng thái bất khả\\
Cận thay đổi tăng dần & Totalizer/đầu ra thứ tự & Giao diện cận ổn định\\
\bottomrule
\end{tabularx}
\end{center}

Đây là cấu hình khởi đầu; Chương 4 trình bày cách đánh giá bộ biểu diễn trên phân
phối bài toán mục tiêu; Chương 5--6 phát triển các bộ đếm khi nhiều
ràng buộc chia sẻ cấu trúc.

\section{Bài học thiết kế}

\begin{enumerate}
  \item Chọn bộ biểu diễn theo cận và đầu ra cần đọc, không chỉ theo tên ràng buộc.
  \item Biểu diễn AMO từng cặp vẫn là đối chứng mạnh khi miền nhỏ.
  \item Sequential Counter phù hợp với ngưỡng nhỏ và cấu trúc tiền tố.
  \item Totalizer/mạng đếm thuận lợi khi cần nhiều cận hoặc giải tăng dần.
  \item Ràng buộc đúng \(K\) nên chia sẻ một cấu trúc tổng nếu có thể.
\end{enumerate}

\subsection*{Câu hỏi nghiên cứu}
\begin{enumerate}
  \item Dựng Sequential Counter cho \(n=5,k=2\) và liệt kê các thanh ghi.
  \item So sánh EXK bằng hai bộ đếm độc lập với một Totalizer chung.
  \item Với trọng số \((1,2,3,5)\), vẽ các trạng thái BDD cho cận \(B=5\).
\end{enumerate}

\chapter{Thiết kế thực nghiệm cho phép biểu diễn SAT}
\chapterlead{Đánh giá trên bộ kiểm chuẩn không chỉ trả lời bộ biểu diễn nào nhanh
hơn. Một thực nghiệm tốt còn phải giải thích hiệu ứng đến từ kích thước, sức
lan truyền, bộ giải, đặc trưng hay cấu trúc của bài toán thử.}

\index{benchmark}\index{timeout}\index{PAR-2}
\index{performance profile}\index{cactus plot}
\index{virtual best solver}\index{shifted geometric mean}

\flowdiagram{Câu hỏi}{Ma trận thiết kế}{Bộ kiểm chuẩn}
  {Thước đo}{Phân tích}

\section{Thực nghiệm là một phép đo có mô hình}

Một bảng thời gian chạy không tự nói bộ biểu diễn nào tốt hơn. Quan sát được tạo
bởi cả
chuỗi
\[
  \text{bài toán thử}
  \rightarrow\text{mô hình}
  \rightarrow\text{bộ biểu diễn}
  \rightarrow\text{bộ tiền xử lý}
  \rightarrow\text{cấu hình bộ giải}
  \rightarrow\text{phần cứng}.
\]
Thay bất kỳ mắt xích nào cũng có thể đổi thứ hạng. Vì vậy kết luận cần nêu rõ
\emph{population} (quần thể mà kết luận hướng đến): một họ bài toán thử, một
miền tham số, một bộ giải hay một danh mục bộ giải. Cần phân biệt
\emph{internal validity} (hiệu lực nội tại: hiệu ứng thật sự do yếu tố
đang xét) với \emph{external validity} (hiệu lực khái quát: hiệu ứng
có còn giữ ngoài tập kiểm chuẩn hay không).

Trong nghiên cứu phép biểu diễn, ba loại bằng chứng bổ sung nhau:
\begin{itemize}
  \item \textbf{bằng chứng cơ chế}: công thức đếm, vết lan truyền và xung đột
        nhỏ giải thích vì sao hiệu ứng xuất hiện;
  \item \textbf{bằng chứng có kiểm soát}: \emph{ablation}
        (thí nghiệm phân tích thành phần (ablation study)) giữ cố định các yếu tố khác;
  \item \textbf{bằng chứng so sánh}: đánh giá trên nhiều họ bài toán và nhiều
        bộ giải để kiểm tra độ bền của kết luận.
\end{itemize}
Một kết quả có giá trị tham khảo cao cần ít nhất hai loại; thời gian chạy đơn
lẻ chỉ là bằng chứng so sánh yếu nếu thiếu cơ chế và biến kiểm soát.

\section{Bắt đầu từ một câu hỏi nhân quả}

Một nghiên cứu về phép biểu diễn thường có một trong bốn câu hỏi:
\begin{enumerate}
  \item phép biểu diễn mới giảm kích thước bao nhiêu;
  \item sức lan truyền thay đổi như thế nào;
  \item thời gian chạy và số bài giải được thay đổi ra sao;
  \item một mô-đun có tương tác với phá đối xứng, giải tăng dần hoặc bộ giải không.
\end{enumerate}

Câu hỏi càng rõ, ma trận thiết kế càng nhỏ. Nếu mục tiêu là đo bộ đếm dùng
chung, hãy giữ cố định bộ giải, tiền xử lý, chiến lược tìm kiếm và quy trình
xử lý hàm mục tiêu; chỉ thay mô-đun bộ đếm. Nếu đồng thời thay bộ giải và cách
phá đối xứng, kết quả không phân định được đóng góp riêng lẻ của từng thành phần.

\begin{designrule}{Đơn vị của thí nghiệm phân tích thành phần (ablation study)}
Một đặc trưng không có hiệu ứng độc lập khi làm thay đổi giao diện của phép biểu diễn.
Trong trường hợp đó, dùng cấu hình
\(\text{phép biểu diễn}\times\text{đặc trưng}\) làm đơn vị so sánh.
\end{designrule}

\section{Ma trận thiết kế}

Một ma trận điển hình gồm:
\[
  \text{bộ biểu diễn}
  \times
  \text{bộ giải}
  \times
  \text{đặc trưng}
  \times
  \text{bài toán thử}
  \times
  \text{random seed}.
\]
Ở đây \emph{random seed} (hạt giống ngẫu nhiên) là giá trị khởi tạo bộ sinh
số giả ngẫu nhiên; giữ lại giá trị này giúp tái lập đúng chuỗi lựa chọn ngẫu
nhiên của một lần chạy.
Không nhất thiết chạy tích Descartes đầy đủ. Có thể dùng hai giai đoạn:
\begin{enumerate}
  \item \textbf{sàng lọc}: chọn vài bộ giải đại diện và miền tham số rộng;
  \item \textbf{xác nhận}: giữ các cấu hình có triển vọng, tăng số bài toán
        thử và số hạt giống ngẫu nhiên.
\end{enumerate}

Mọi cấu hình phải dùng cùng ngữ nghĩa và bộ kiểm tra. Nếu một bộ biểu diễn có tiền
xử lý riêng, thời gian tiền xử lý phải được tính vào tổng thời gian chạy hoặc
được báo thành một cột riêng.

\section{Bộ kiểm chuẩn và phân tầng}

Một bộ kiểm chuẩn nên có siêu dữ liệu cấu trúc: số biến miền, mật độ ràng buộc,
độ rộng cửa sổ, ngưỡng, lớp đối xứng, khoảng \(LB\)--\(UB\), hoặc các đặc trưng
riêng của bài toán. Phân tầng theo những đặc trưng này giúp giải thích tại sao
một phép biểu diễn thắng trong một vùng.

Ba nguồn bài toán thử thường dùng:
\begin{itemize}
  \item bài toán chuẩn từ văn liệu;
  \item bài toán thực tế có nguồn gốc rõ ràng;
  \item bài toán sinh tổng hợp để quét có kiểm soát một tham số.
\end{itemize}

Các bài toán sinh tổng hợp cho phép quan sát điểm chuyển theo \(n,k,w\), mật độ
hoặc độ rộng miền; dữ liệu thực cung cấp bối cảnh sử dụng, còn bài toán chuẩn
tạo khả năng so sánh với kết quả trước. Báo riêng ba nguồn giúp người đọc hiểu
đúng vai trò của từng tập.

\section{Tách tinh chỉnh và đánh giá}

Tùy chọn bộ giải, độ rộng khối, kích thước nhóm và ngưỡng biểu diễn là các siêu
tham số. Nếu chọn chúng trên chính tập kiểm thử, thời gian chạy báo cáo sẽ lạc
quan quá mức. Một quy trình đơn giản:
\begin{enumerate}
  \item chia bộ kiểm chuẩn thành tập tinh chỉnh và tập đánh giá;
  \item khóa cấu hình sau khi tinh chỉnh;
  \item chỉ dùng tập đánh giá cho kết luận cuối;
  \item nếu dữ liệu ít, dùng \emph{cross-validation} (kiểm định chéo) theo họ
        thay vì chia ngẫu nhiên từng bài toán.
\end{enumerate}

Tinh chỉnh theo họ tránh việc các biến thể gần như trùng nhau xuất hiện ở cả
hai tập.

\section{Các thước đo}

\subsection{Kích thước và xây dựng}

Ghi ít nhất:
\[
  n_v,\quad n_c,\quad n_\ell,\quad
  \#\text{binary},\quad
  \#\text{ternary},\quad
  t_{\text{encode}},\quad
  \text{bộ nhớ cực đại}.
\]
Thời gian của bộ giải nên đo từ sau khi dựng công thức; đồng thời báo thời gian
toàn trình nếu bộ sinh công thức chiếm một phần đáng kể.

\subsection{Quá trình tìm kiếm}

Số quyết định, lượt lan truyền, xung đột, lần khởi động lại và mệnh đề đã học
giúp giải thích cơ chế tạo ra thời gian chạy quan sát được:
\[
  \text{lượt lan truyền/xung đột},\qquad
  \text{xung đột/giây},\qquad
  \text{bộ nhớ/biến}.
\]

\subsection{Right-censored runtime: thời gian chạy có kiểm duyệt}

Khi một lần chạy hết thời hạn, ta chỉ biết thời gian cần thiết lớn hơn giới hạn;
dữ liệu như vậy được gọi là \emph{right-censored} (dữ liệu bị giới hạn thời gian).
Vì vậy, thống kê thời gian chạy phải tính cả bài đã giải và bài hết thời hạn
bằng một trong các cách:
\begin{itemize}
  \item số bài giải được tại một giới hạn thời gian cố định;
  \item PAR-2 hoặc PAR-10, tức thời gian trung bình có phạt khi hết thời hạn;
  \item \emph{cactus plot} (biểu đồ Cactus (Cactus plot));
  \item \emph{performance profile} (hồ sơ hiệu năng);
  \item thời gian đạt mục tiêu cho bài toán tối ưu.
\end{itemize}

Hồ sơ hiệu năng của Dolan--Moré chuẩn hóa thời gian chạy theo cấu hình tốt nhất
trên từng bài toán thử. Với tỷ số
\[
  r_{p,s}=
  \frac{t_{p,s}}{\min_{s'}t_{p,s'}},
\]
đường
\[
  \rho_s(\tau)=
  \frac{1}{|P|}
  \card{\set{p\in P\mid r_{p,s}\le\tau}}
\]
cho biết tỷ lệ bài mà cấu hình \(s\) nằm trong hệ số \(\tau\) so với tốt nhất
\parencite{dolan2002profiles}. Giá trị tại \(\tau=1\) đo số lần thắng; phần
đuôi đo độ bền. Bài hết thời hạn phải được gán một tỷ số lớn nhất nhất quán và
được mô tả rõ.

\begin{keyidea}{Đọc số bài giải được và thời gian chạy cùng nhau}
Một cấu hình giải thêm các bài khó có thể có tổng thời gian chạy lớn hơn cấu
hình hết thời hạn sớm. Vì vậy không xếp hạng bằng một cột thời gian đơn lẻ.
\end{keyidea}

\begin{workedexample}{Hết thời hạn làm thay đổi cách xếp hạng}
Với giới hạn \(30\) giây, ta có:
\[
\begin{array}{c|ccc}
 & p_1&p_2&p_3\\\hline
\text{Đối chứng}&2&5&\text{hết hạn}\\
\text{Mới}&1&10&20
\end{array}
\]
Đối chứng nhanh hơn trên \(p_2\), nhưng chỉ giải \(2/3\) bài; cấu hình mới giải
cả ba. Nếu PAR-2 gán \(60\) giây cho lần chạy hết thời hạn:
\[
  \operatorname{PAR2}(\text{Đối chứng})=\frac{2+5+60}{3}=22.33,
  \qquad
  \operatorname{PAR2}(\text{Mới})=\frac{1+10+20}{3}=10.33.
\]
Trung bình chỉ trên bài đã giải sẽ cho đối chứng \(3.5\) giây và cấu hình mới
\(10.33\) giây, từ đó tạo kết luận ngược vì đã bỏ chính bài khó nhất.
\end{workedexample}

\begin{figure}[tb]
\centering
\begin{tikzpicture}[satfig,x=18mm,y=11mm]
  \draw[->,thick,Ink] (0,0)--(5.4,0) node[right,satlabel] {\(\tau\)};
  \draw[->,thick,Ink] (0,0)--(0,3.5) node[above,satlabel] {\(\rho_s(\tau)\)};
  \foreach \y/\lab in {1/0.33,2/0.67,3/1.00}{
    \draw[Rule] (0,\y)--(5,\y);
    \node[left=4pt,satlabel] at (0,\y) {\lab};
  }
  \draw[very thick,Blue]
    (0.5,0)--(0.5,1)--(1.2,1)--(1.2,2)--(3.8,2)--(3.8,3)--(5,3);
  \draw[very thick,Teal]
    (0.5,0)--(0.5,1)--(.9,1)--(.9,2)--(1.8,2)--(1.8,3)--(5,3);
  \node[satlabel,text=Blue] at (4.5,2.6) {Đối chứng};
  \node[satlabel,text=Teal] at (2.5,3.25) {Mới};
  \draw[Ink] (.5,2pt)--(.5,-2pt) node[below=3pt,satlabel]{1};
  \draw[Ink] (1.8,2pt)--(1.8,-2pt) node[below=3pt,satlabel]{2};
  \draw[Ink] (3.8,2pt)--(3.8,-2pt) node[below=3pt,satlabel]{4};
\end{tikzpicture}
\caption{Hồ sơ hiệu năng (performance profile): phần đầu cho thấy số bài giải nhanh nhất, phần đuôi
cho độ bền so với cấu hình tốt nhất trên từng bài toán thử.}
\label{fig:performance-profile}
\end{figure}

\subsection{Biểu đồ Cactus (Cactus plot) và biểu đồ phân tán}

Với một cấu hình \(s\), sắp tăng thời gian của các bài đã giải:
\[
  t_{(1),s}\le t_{(2),s}\le\cdots\le t_{(q_s),s}.
\]
Biểu đồ Cactus (Cactus plot) đặt số bài đã giải \(j\) trên trục ngang và
\(t_{(j),s}\) trên trục dọc. Đường đi xa hơn sang phải biểu thị độ bền; đường
thấp hơn biểu thị giải nhanh hơn ở cùng số bài đã giải. Đường dừng ở \(q_s\),
nên các bài hết thời hạn không bị âm thầm loại khỏi kết luận.

Biểu đồ phân tán dùng so sánh ghép cặp: mỗi điểm ứng với cùng một bài toán thử,
có tọa độ \((t_{p,A},t_{p,B})\), thường trên thang lôgarit. Điểm phía trên đường chéo
\(t_A=t_B\) nghĩa là \(A\) nhanh hơn; điểm phía dưới nghĩa là \(B\) nhanh hơn.
Bài hết thời hạn cần được đặt ở biên \(T\) bằng ký hiệu riêng hoặc mũi tên, và
SAT/UNSAT nên dùng dấu điểm khác nhau. Biểu đồ phân tán cho biết \emph{bài toán
nào} tạo khác biệt; biểu đồ Cactus (Cactus plot) cho biết \emph{toàn tập} được bao phủ
đến đâu.

\begin{figure}[tb]
\centering
\begin{tikzpicture}[satfig,x=8mm,y=8mm]
  \begin{scope}
    \draw[->,Ink] (0,0)--(5.2,0) node[right,satlabel]{số bài giải được};
    \draw[->,Ink] (0,0)--(0,4.2) node[above,satlabel]{thời gian};
    \draw[Blue,very thick] (.5,.35)--(1.3,.55)--(2.1,.9)--(2.9,1.45)--(3.7,2.4);
    \draw[Teal,very thick] (.5,.45)--(1.3,.7)--(2.1,1.2)--(2.9,2.1)--(4.5,3.5);
    \node[satlabel,text=Blue] at (3.55,2.05) {A};
    \node[satlabel,text=Teal] at (4.55,3.15) {B};
    \node[satlabel,font=\sffamily\bfseries\small] at (2.5,4.6) {Cactus plot};
  \end{scope}
  \begin{scope}[xshift=62mm]
    \draw[->,Ink] (0,0)--(4.6,0) node[right,satlabel]{\(t_A\)};
    \draw[->,Ink] (0,0)--(0,4.2) node[above,satlabel]{\(t_B\)};
    \draw[Rule,dashed] (.25,.25)--(4,4);
    \foreach \x/\y in {.6/1.2,1.1/.8,1.5/2.7,2.1/1.5,2.6/3.4,3.4/2.8}
      \fill[Blue] (\x,\y) circle (2.2pt);
    \draw[Gold,very thick] (4,2.4)--(4.35,2.4);
    \draw[Gold,very thick] (2.1,4)--(2.1,4.35);
    \node[satlabel,text=Blue,rotate=45] at (3.25,3.55) {\(t_A=t_B\)};
    \node[satlabel,font=\sffamily\bfseries\small] at (2.2,4.6) {Phân tán ghép cặp};
  \end{scope}
\end{tikzpicture}
\caption{Biểu đồ Cactus (Cactus plot) (trái) giữ số bài giải được trong hình; biểu đồ
phân tán (phải) giữ ghép cặp theo bài toán thử. Các dấu ở biên minh họa lần
chạy hết thời hạn của một cấu hình.}
\label{fig:cactus-scatter}
\end{figure}

\section{SAT và UNSAT là hai chế độ khác nhau}

Các bài toán SAT cần tìm một nghiệm làm chứng; các bài toán UNSAT cần đóng toàn bộ
không gian tìm kiếm. Một phép biểu diễn có thể mạnh ở SAT nhưng yếu ở UNSAT hoặc
ngược lại.
Bảng kết quả nên tách:
\[
  \#SAT,\quad \#UNSAT,\quad t_{SAT},\quad t_{UNSAT}.
\]

Trong tối ưu, các lần gọi bộ giải gần giá trị tối ưu thường có cả hai trạng
thái: mô hình tại cận khả thi và \UNSAT tại cận tốt hơn. Vì thế, hồ sơ theo
toàn bộ dãy cận chứa nhiều thông tin hơn thời gian chạy cuối cùng.

\section{Random seed và biến thiên giữa các lần chạy}

\emph{Random seed}, được gọi là hạt giống ngẫu nhiên sau lần giới thiệu đầu,
không phải dữ liệu của bài toán; tham số này là giá trị khởi tạo (seed) cho bộ sinh số giả ngẫu nhiên
của thuật toán. CDCL có yếu tố ngẫu nhiên trong phá hòa và tìm kiếm. Với bài
toán khó, cần chạy nhiều hạt giống ngẫu nhiên và báo trung vị, các phân vị hoặc
khoảng tin cậy. So sánh ghép cặp phải dùng cùng bài toán, cùng chỉ số hạt giống
ngẫu nhiên và cùng giới hạn thời gian, bộ nhớ.

Nếu bộ biểu diễn là tất định nhưng bộ giải có yếu tố ngẫu nhiên, hạt giống ngẫu
nhiên thuộc về bộ giải. Nếu bộ sinh dùng thứ tự ngẫu nhiên, cần tách hạt giống
ngẫu nhiên của bộ sinh khỏi hạt giống ngẫu nhiên của bộ giải.

\section{So sánh ghép cặp và độ lớn hiệu ứng}

Mỗi bài toán thử tạo một cặp quan sát tự nhiên giữa hai bộ biểu diễn. Vì phân bố
thời gian chạy thường lệch và có dữ liệu hết thời hạn, chỉ báo trung bình và
một kiểm định trên các lần chạy độc lập là chưa phù hợp. Nên báo đồng thời:
\begin{itemize}
  \item số lần thắng/hòa/thua trên cùng bài toán và cùng hạt giống ngẫu nhiên;
  \item trung vị của lôgarit hệ số tăng tốc
        \(\log(t_{\mathrm{base}}/t_{\mathrm{new}})\);
  \item khoảng bất định bằng \emph{bootstrap} (lấy mẫu lại) theo \emph{họ}, không chỉ
        theo từng bài toán;
  \item chênh lệch số bài giải được tại nhiều mốc thời gian.
\end{itemize}

Độ lớn hiệu ứng quan trọng hơn một giá trị \(p\) đứng riêng. Một bộ biểu diễn nhanh
hơn \(3\%\) trên hàng nghìn bài dễ nhưng làm mất năm bài khó có ý nghĩa khác
với bộ biểu diễn chỉ giảm ít trung vị nhưng mở rộng vùng giải được. Với tối ưu,
nên so cả thời gian đạt nghiệm tốt nhất đã biết, thời gian hoàn tất chứng minh
và khoảng cách tối ưu theo thời gian.

\section{Phân tích theo vùng cấu trúc}

Một kết luận toàn cục dễ che mất điểm chuyển ưu thế. Bộ đếm tuần tự có thể tốt
khi \(k/n\) nhỏ; mạng đếm có thể cải thiện khi cần nhiều ngưỡng; biểu diễn thứ tự
có thể thắng khi các khoảng cấm rộng. Vì vậy, trước khi vẽ kết quả, cần xác
định trước các tầng theo đặc trưng có ý nghĩa thuật toán:
\[
  \frac{k}{n},\quad
  \text{độ chồng lấn cửa sổ},\quad
  \text{độ rộng miền},\quad
  \text{mật độ ràng buộc},\quad
  \text{quy mô đối xứng}.
\]
Sau đó dùng bản đồ nhiệt hoặc bảng theo tầng để tìm \emph{miền ưu thế} thay vì
chỉ tìm ``người thắng''. Cách phân tích này cũng cung cấp dữ liệu cho bộ chọn
phép biểu diễn ở Chương 3.

\section{Kiểm đúng trước khi đo nhanh}

Trước khi chạy bộ kiểm chuẩn lớn:
\begin{enumerate}
  \item kiểm thử vét cạn trên các bài toán nhỏ;
  \item kiểm thử vi sai giữa hai bộ biểu diễn;
  \item giải mã và kiểm tra mọi mô hình SAT dùng trong bảng;
  \item tính lại hàm mục tiêu từ nghiệm, không đọc trực tiếp chi phí nội bộ;
  \item giữ một số công thức \CNF và chứng minh nhỏ làm phép thử hồi quy.
\end{enumerate}

Với \(n\) nhỏ, bộ liệt kê sinh mọi phép gán cho biến chính. So sánh phép chiếu
của các mô hình \CNF với tập nghiệm miền. Cách này phát hiện lỗi biên hiệu quả
hơn việc chỉ chạy các bài toán kiểm chuẩn lớn.

\section{Gói tái lập và khả năng truy vết}

\emph{Artifact} (gói phần mềm--dữ liệu tái lập) tối thiểu gồm:
\begin{itemize}
  \item mã nguồn và mã định danh bản sửa đổi chính xác;
  \item bản kê bài toán thử và mã kiểm tra toàn vẹn;
  \item môi trường cùng phiên bản bộ giải;
  \item lệnh chạy và giới hạn tài nguyên;
  \item nhật ký thô, bộ phân tích nhật ký và chương trình tạo bảng;
  \item bộ giải mã và bộ kiểm tra.
\end{itemize}
Gói tái lập tạo một đường đi kiểm tra được từ phát biểu đến bằng chứng thô.

\begin{figure}[tb]
\centering
\begin{tikzpicture}[satfig,node distance=8mm]
  \node[satbox] (claim) {Phát biểu\\``nhanh hơn''};
  \node[satbox,right=of claim] (table) {Bảng/biểu đồ};
  \node[satbox,right=of table] (parsed) {Nhật ký\\đã phân tích};
  \node[satbox,right=of parsed] (raw) {Nhật ký thô};
  \node[satbox,right=of raw] (run) {Lệnh chạy\\+ bản sửa đổi};
  \draw[satedge] (claim)--(table);
  \draw[satedge] (table)--(parsed);
  \draw[satedge] (parsed)--(raw);
  \draw[satedge] (raw)--(run);
  \node[satlabel,below=6mm of parsed]
    {Mỗi mũi tên phải có chương trình, mã kiểm tra hoặc siêu dữ liệu để truy vết.};
\end{tikzpicture}
\caption{Chuỗi truy vết tối thiểu của một phát biểu thực nghiệm.}
\label{fig:claim-artifact-chain}
\end{figure}

Số liệu thực nghiệm có thể được kiểm tra ở ba tầng:
\begin{enumerate}
  \item \textbf{báo cáo}: bài viết mô tả tập bài toán thử, giới hạn thời gian
        và các cấu hình đối chứng;
  \item \textbf{tái tạo}: mã nguồn, nhật ký và dữ liệu cho phép dựng lại các
        bảng và biểu đồ;
  \item \textbf{tái lập độc lập}: một môi trường khác thực hiện lại thí
        nghiệm theo cùng thiết kế.
\end{enumerate}
Tầng báo cáo hỗ trợ việc diễn giải kết quả; tầng tái tạo kiểm tra đường đi từ
nhật ký đến bảng; tầng tái lập độc lập đánh giá độ bền khi môi trường thực thi
thay đổi.

\section{Mẫu quy trình thực nghiệm}

\begin{summarybox}
\textbf{Câu hỏi.} Mô-đun nào được đánh giá?\\
\textbf{Giả thuyết.} Kỳ vọng tác động lên kích thước, sức lan truyền hay thời gian chạy?\\
\textbf{Yếu tố.} Bộ biểu diễn, bộ giải, đặc trưng, họ bài toán, hạt giống ngẫu nhiên.\\
\textbf{Biến kiểm soát.} Ngữ nghĩa, bộ kiểm tra, cận, tiền xử lý, phần cứng.\\
\textbf{Thước đo.} Kích thước, thời gian dựng, bộ đếm tìm kiếm, số bài giải được/PAR/hồ sơ.\\
\textbf{Kiểm đúng.} Kiểm thử vét cạn, kiểm thử vi sai và bộ kiểm tra mô hình.\\
\textbf{Đầu ra.} Nhật ký thô, biểu đồ, bảng tổng hợp và bản kê gói tái lập.
\end{summarybox}

\section{Bài học thiết kế}

\begin{enumerate}
  \item Đánh giá thực nghiệm bắt đầu từ một câu hỏi, không bắt đầu từ một bảng cấu hình.
  \item Tách tinh chỉnh khỏi đánh giá và phân tầng theo họ bài toán.
  \item Báo cả kích thước, bộ đếm tìm kiếm và thời gian chạy có dữ liệu hết hạn.
  \item Tách SAT/UNSAT và dùng nhiều hạt giống ngẫu nhiên ở vùng khó.
  \item Bộ kiểm tra mô hình là điều kiện trước của mọi bảng hiệu năng.
\end{enumerate}

\subsection*{Câu hỏi nghiên cứu}
\begin{enumerate}
  \item Thiết kế thí nghiệm phân tích thành phần (ablation study) \(2\times2\) cho bộ biểu diễn và phá đối xứng.
  \item Với ba lần hết thời hạn trong mười bài toán thử, tính PAR-2 và so với
        trung bình trên các bài đã giải.
  \item Lập một bản kê gói tái lập tối thiểu cho một thí nghiệm SAT tăng dần.
\end{enumerate}


\part{Tái sử dụng cấu trúc trong phép mã hóa SAT}
\chapter{Bộ đếm dùng chung cho các cửa sổ chồng lấn}
\chapterlead{Khi một ràng buộc trượt trên chuỗi, hai cửa sổ kề nhau chia sẻ gần
toàn bộ biến. Mã hóa từng cửa sổ độc lập sẽ lặp trạng thái; bộ đếm dùng chung
(\emph{shared counter}) biến phần giao đó thành một mô-đun có thể tái sử dụng.}

\index{SEQUENCE constraint}\index{shared counter}
\index{Sequential Counter Ladder}\index{Ladder AMK}
\index{sliding window}

\flowdiagram{Chuỗi Boolean}{Phân hoạch khối}{Bộ đếm cục bộ}
  {Bộ nối}{Họ cửa sổ}

\section{SEQUENCE trong văn liệu về ràng buộc}

Ràng buộc cửa sổ thuộc một họ lâu đời hơn các bộ đếm dùng chung.
\textsc{Sequence}
yêu cầu mọi đoạn liên tiếp độ dài cố định có số biến nhận một giá trị trong
khoảng \([\ell,u]\). Nó xuất hiện trong xếp thứ tự xe, lập lịch ca, xác định
cỡ lô và các mô hình chuỗi. Vấn đề cốt lõi không chỉ là số literal đúng trong
một cửa sổ, mà còn là sự phụ thuộc giữa toàn bộ các cửa sổ chồng lấn.

Văn liệu đã phát triển ít nhất bốn cách nhìn:
\begin{description}[style=nextline]
  \item[Phân rã theo cửa sổ]
  Tạo một ràng buộc đếm cho mỗi cửa sổ. Cách này có tính mô-đun và dễ dùng,
  nhưng nếu các trạng thái phụ không được chia sẻ thì phần giao bị biểu diễn
  nhiều lần.

  \item[Tổng tiền tố]
  Đặt \(S_i=\sum_{j\le i}x_j\), khi đó mỗi cửa sổ là
  \(\ell\le S_{i+w}-S_i\le u\). Mọi cửa sổ dùng chung một dãy trạng thái tiền tố.

  \item[Ôtômát/ràng buộc chính quy]
  Trạng thái ghi đủ lịch sử cần thiết để quyết định ký hiệu tiếp theo. Một
  ràng buộc \textsc{Regular} có thể biểu diễn việc chuỗi có thuộc một ngôn ngữ
  hữu hạn trạng thái hay không \parencite{pesant2004regular}.

  \item[Luồng hoặc đồ thị]
  Biểu diễn các hiệu tiền tố và cận cửa sổ bằng mạng luồng, từ đó xây dựng
  bộ lan truyền toàn cục.
\end{description}

So sánh có hệ thống các phép mã hóa \textsc{Sequence} cho thấy cách phân rã,
ôtômát và phương pháp dựa trên tổng khác nhau rõ về kích thước và mức nhất quán
\parencite{brand2007sequence}. Các bộ lan truyền dựa trên luồng tiếp tục khai
thác cấu trúc toàn cục để đạt khả năng lọc mạnh \parencite{maher2008sequence};
những khung tổng quát hơn còn xem ràng buộc trên chuỗi như tổ hợp của một
\emph{chữ ký} và một \emph{đặc trưng} cần đếm
\parencite{monette2014sequences}.

\begin{center}
\captionof{table}{Các cách tái sử dụng trạng thái trong ràng buộc cửa sổ.}
\label{tab:window-reuse}
\begin{tabularx}{\textwidth}{@{}>{\bfseries\sffamily}p{.19\textwidth}YYY@{}}
\toprule
Họ & Trạng thái dùng chung & Điểm mạnh & Chi phí/rủi ro\\
\midrule
Mỗi cửa sổ & Không hoặc ít & Đơn giản, dùng bộ mã hóa sẵn & Lặp trạng thái phụ\\
Tổng tiền tố & Tổng từ đầu chuỗi & Mọi cửa sổ là hiệu hai tiền tố & Miền tổng toàn cục\\
Ôtômát & Hậu tố/lịch sử & Xử lý mẫu phức tạp & Bùng nổ trạng thái theo \(w\)\\
Luồng & Nút và cung & Lọc toàn cục & Bộ lan truyền/phép mã hóa phức tạp\\
Bộ đếm theo khối & Tiền tố/hậu tố trong khối & Cục bộ và tái sử dụng có kiểm soát &
Cần chọn phân hoạch\\
\bottomrule
\end{tabularx}
\end{center}

Bộ đếm dùng chung theo khối của chương này nằm giữa hai cực. Nó không tạo một
trạng thái toàn cục cho mọi tiền tố, nhưng cũng không xây bộ đếm độc lập cho
từng cửa sổ. Phân hoạch quyết định phạm vi tái sử dụng, còn bộ nối giữ ngữ
nghĩa của từng cửa sổ.

\section{Họ ràng buộc cửa sổ}

Cho chuỗi
\[
  \Omega=\langle x_1,\ldots,x_n\rangle.
\]
Cửa sổ độ rộng \(w\) bắt đầu tại \(i\):
\[
  W_i(w)=\langle x_i,\ldots,x_{i+w-1}\rangle.
\]
AMK bậc thang là họ
\[
  \sum_{x\in W_i(w)}x\le k
  \qquad i=1,\ldots,n-w+1.
\]
Các ràng buộc chuỗi tổng quát cho phép cả cận dưới và cận trên:
\[
  \ell\le
  \sum_{x\in W_i(w)}x
  \le u.
\]

Nếu mỗi cửa sổ dùng một bộ đếm tuần tự riêng, kích thước xấp xỉ
\[
  O((n-w+1)wk).
\]
Trong khi đó, hai cửa sổ kề nhau chỉ khác một biến đi ra và một biến đi vào.
Mục tiêu của phép mã hóa dùng chung là đưa chi phí về gần
\[
  O(nk)
\]
hoặc \(O(nk+n_c)\), với \(n_c\) là số bộ nối.

\begin{figure}[tb]
\centering
\begin{tikzpicture}[satfig,x=12mm,y=7mm]
  \foreach \i in {1,...,8}{
    \node[satbox,minimum width=10mm,minimum height=8mm] (x\i) at (\i,0)
      {\(x_{\i}\)};
  }
  \draw[Blue,very thick] (0.55,-.8)--(4.45,-.8)
    node[midway,below,satlabel,text=Blue]{\(W_1=\{x_1,\ldots,x_4\}\)};
  \draw[Teal,very thick] (1.55,-1.65)--(5.45,-1.65)
    node[midway,below,satlabel,text=Teal]{\(W_2=\{x_2,\ldots,x_5\}\)};
  \draw[Gold,very thick] (2.55,-2.5)--(6.45,-2.5)
    node[midway,below,satlabel,text=Gold]{\(W_3=\{x_3,\ldots,x_6\}\)};
  \draw[satdim] (2.05,.72)--(4.95,.72)
    node[midway,above,satlabel]{phần trạng thái có thể dùng lại};
\end{tikzpicture}
\caption{Ba cửa sổ độ rộng \(4\) chỉ dịch một vị trí: cửa sổ kề nhau giữ lại
ba trong bốn biến. Xây lại một bộ đếm cho từng hàng sẽ lặp phần giao này.}
\label{fig:overlapping-windows}
\end{figure}

\section{Phân hoạch thành khối}

Chia chuỗi thành các khối liên tiếp
\[
  B_1,B_2,\ldots,B_m,
  \qquad
  \Omega=B_1\cdot B_2\cdots B_m.
\]
Trong mỗi khối, xây hai họ biến ngưỡng:
\[
  F^b_{i,q}=1
  \iff
  \sum_{j=1}^{i}B_b[j]\ge q,
\]
\[
  R^b_{i,q}=1
  \iff
  \sum_{j=i}^{\card{B_b}}B_b[j]\ge q.
\]
\(F\) đếm tiền tố từ trái sang phải; \(R\) đếm hậu tố từ phải sang trái. Một
cửa sổ cắt biên giữa \(B_b,B_{b+1}\) là hợp của một hậu tố bên trái và một
tiền tố bên phải. Số literal đúng trong cửa sổ được đọc từ \(R^b\) và
\(F^{b+1}\).

\begin{figure}[tb]
\centering
\begin{tikzpicture}[satfig,x=11mm,y=9mm]
  \foreach \i in {1,...,8}{
    \node[satbox,minimum width=9mm] (x\i) at (\i,0) {\(x_{\i}\)};
  }
  \node[satlabel,text=Blue] at (2.5,.75) {\(B_1\)};
  \node[satlabel,text=Teal] at (6.5,.75) {\(B_2\)};
  \draw[Blue,thick,rounded corners=2pt] (.55,-.48) rectangle (4.45,.48);
  \draw[Teal,thick,rounded corners=2pt] (4.55,-.48) rectangle (8.45,.48);
  \draw[Gold,ultra thick] (2.55,-.75)--(6.45,-.75)
    node[midway,below,satlabel,text=Gold]{cửa sổ \(W_3=x_3,\ldots,x_6\)};
  \node[satstate] (r) at (3,-2.25) {\(R^1\)};
  \node[satstate] (f) at (6,-2.25) {\(F^2\)};
  \node[satbox,fill=PaleGold,draw=Gold] (c) at (4.5,-3.55)
    {bộ nối\\\(\neg R^1_{3,q}\lor\neg F^2_{2,k-q+1}\)};
  \draw[satedge] (x3.south)--(r);
  \draw[satedge] (x4.south)--(r);
  \draw[satedge] (x5.south)--(f);
  \draw[satedge] (x6.south)--(f);
  \draw[satedge] (r)--(c);
  \draw[satedge] (f)--(c);
\end{tikzpicture}
\caption{Một cửa sổ cắt biên được phân rã thành hậu tố của khối trái và tiền tố
của khối phải; bộ nối chỉ đọc các ngưỡng cần thiết.}
\label{fig:block-connector}
\end{figure}

\begin{example}
Giả sử cửa sổ gồm hậu tố \(A\) của khối trái và tiền tố \(B\) của khối phải.
Đặt \(T(S,q)\) là literal ``\(\sum_{x\in S}x\ge q\)'', với hai quy ước biên
\(T(S,q)=\top\) khi \(q\le 0\) và \(T(S,q)=\bot\) khi \(q>\card{S}\).
Điều kiện at-most-\(k\) tương đương với việc không tồn tại
\(q\in\set{0,\ldots,k+1}\) sao cho
\[
  \sum_{x\in A}x\ge q
  \quad\text{và}\quad
  \sum_{x\in B}x\ge k+1-q.
\]
Bộ nối vì thế chứa
\[
  \neg T(A,q)
  \lor
  \neg T(B,k+1-q),
  \qquad q=0,\ldots,k+1.
\]
Hai giá trị biên \(q=0\) và \(q=k+1\) không phải chi tiết hình thức: chúng
phát hiện trường hợp toàn bộ \(k+1\) giá trị đúng nằm ở chỉ một phía của biên.
\end{example}

\begin{workedexample}{Đọc một cửa sổ qua hai thanh ghi}
Lấy \(n=8,w=4,k=2\), chia
\(B_1=(x_1,\ldots,x_4)\), \(B_2=(x_5,\ldots,x_8)\), và xét
\[
  x=(1,0,1,0\mid 1,0,0,1).
\]
Cửa sổ \(W_3=(x_3,x_4\mid x_5,x_6)=(1,0\mid1,0)\) có tổng \(2\).
Hậu tố trái đạt ngưỡng \(1\) nhưng không đạt \(2\):
\(R^1_{3,1}=1,R^1_{3,2}=0\). Tiền tố phải cũng có
\(F^2_{2,1}=1,F^2_{2,2}=0\). Với \(k=2\), bốn bộ nối ứng với
\(q=0,1,2,3\) cấm các cặp ngưỡng
\((\top,F_3),(R_1,F_2),(R_2,F_1),(R_3,\top)\).
Tất cả đều thỏa: hai mệnh đề biên xác nhận mỗi phía riêng lẻ không đạt \(3\),
còn ở hai mệnh đề giữa luôn có một ngưỡng bằng \(0\).

Nếu đổi \(x_6\) thành \(1\), tiền tố phải đạt ngưỡng \(2\). Khi đó
\(\neg R^1_{3,1}\lor\neg F^2_{2,2}\) sai, phát hiện ngay tổng cửa sổ bằng
\(3\). Ví dụ cho thấy bộ nối không cộng lại toàn cửa sổ; nó chỉ ghép hai
chứng cứ ngưỡng.
\end{workedexample}

\section{Bộ nối như một giao diện}

Bộ đếm cục bộ không biết cửa sổ nào sẽ đọc nó. Nó chỉ cung cấp đặc tả của các
literal ngưỡng. Bộ nối biết cách một ràng buộc cắt biên và chọn các ngưỡng cần
ghép.

Với at-least-\(k\):
\[
  \sum_{x\in A\cup B}x\ge k
\iff
  \bigwedge_{q=1}^{k}
  \left[
    \left(\sum_{x\in A}x\ge q\right)
    \lor
    \left(\sum_{x\in B}x\ge k-q+1\right)
  \right].
\]
Nếu thanh ghi có đặc tả tương đương, mỗi ngoặc vuông trở thành
\[
  R^b_{a,q}\lor F^{b+1}_{c,k-q+1}.
\]

\begin{designrule}{Bên cung cấp và bên sử dụng}
Bộ đếm là bên cung cấp literal ngưỡng; bộ nối là bên sử dụng. Đặc tả hai
chiều cần thiết khi bộ nối dùng cả thanh ghi dương và phủ định để diễn giải
``đạt'' và ``chưa đạt'' ngưỡng.
\end{designrule}

\section{Staircase AMO và SCL}

Với \(k=1\), mỗi cửa sổ thỏa AMO. Thang bộ đếm dùng chung
(\emph{Shared Counter Ladder}, SCL) xây các trạng thái AMO/AMZ trong khối và
nối chúng qua các biên. So với việc tạo xung đột từng cặp trong từng cửa sổ,
SCL không lặp các cặp nằm trong phần giao.

Giả sử \(m=\lceil n/w\rceil\) khối có kích thước gần \(w\). Chi phí bộ đếm cục
bộ tuyến tính theo tổng kích thước khối, còn các bộ nối tuyến tính theo số biên và
ngưỡng. Vì vậy kích thước toàn họ có thể giảm từ bậc hai về tuyến tính theo
\(n\) ở các AMO bậc thang điển hình.

Mã hóa Duplex dùng BDD tiến/lùi trong mỗi cửa sổ và liên kết các trạng thái
\parencite{fazekas2020duplex}. SCL dùng trạng thái tuần tự, nhờ đó giảm số
biến phụ và mệnh đề trong phân tích của
\textcite{truong2025scamo}.

Trong hệ phân loại trên, SCL phục vụ các tập xung đột có dạng khoảng dịch
chuyển và phần giao lớn giữa các khoảng kề nhau. Mã hóa từng cặp và bộ đếm tuần
tự độc lập phù hợp hơn với các tập AMO rời rạc hoặc thiếu một thứ tự tuyến tính
ổn định. Độ chồng lấn vì thế là đặc trưng đầu tiên để chọn giữa trạng thái dùng
chung và bộ đếm cục bộ.

\begin{resultbox}{SCAMO và Anti-Bandwidth}
Trên họ SCAMO với \(n=1000\) và độ rộng cửa sổ thay đổi,
\textcite{truong2025scamo} báo SCL dùng ít biến và mệnh đề hơn Duplex cùng các
phép mã hóa từng cửa sổ. Khi đưa vào bài toán Anti-Bandwidth, SCL đạt kết quả
cạnh tranh trên 24 đồ thị kiểm chuẩn, đặc biệt ở các bài toán lớn nơi chi phí
lặp cửa sổ trở nên rõ rệt.
\end{resultbox}

\section{Mở rộng sang Ladder AMK}

Khi \(k>1\), mỗi khối cần nhiều mức ngưỡng. Ta dùng bộ đếm tuần tự:
\[
  F^b_{i,q}
  \leftrightarrow
  F^b_{i-1,q}\lor
  (B_b[i]\land F^b_{i-1,q-1}),
\]
và hệ thức truy hồi đối xứng cho \(R^b\).

Bộ nối cho cận trên cấm mọi cặp ngưỡng có tổng \(k+1\):
\[
  \neg T(A,q)
  \lor
  \neg T(B,k+1-q),
  \qquad q=0,\ldots,k+1.
\]
Vì thế các thanh ghi thực phải được xây đến ngưỡng \(k+1\), trừ những mức có
thể giản lược thành \(\top\) hoặc \(\bot\) theo kích thước phần cửa sổ.
Nếu một cửa sổ đi qua nhiều hơn một biên, có hai lựa chọn:
\begin{enumerate}
  \item tăng độ rộng khối để mỗi cửa sổ cắt nhiều nhất một biên;
  \item xây bộ nối nhiều ngôi hoặc một bộ đếm trung gian trên tổng của các khối.
\end{enumerate}
Thiết kế thứ nhất đơn giản hơn; thiết kế thứ hai linh hoạt với nhiều width.

Ladder AMK trong \textcite{truong2026ladderamk} phát triển bộ đếm tuần tự dùng
chung cho một tập ràng buộc đếm dạng bậc thang. Giá trị cốt lõi không nằm ở
một mệnh đề riêng lẻ mà ở việc bộ đếm được tạo cho cả họ ràng buộc.

\subsection{Quan hệ với tổng tiền tố và ràng buộc chuỗi tổng quát}

Ladder AMK có thể được mô tả bằng ma trận liên thuộc: mỗi hàng là một ràng
buộc, mỗi cột là một đầu vào, và các giá trị \(1\) trong mỗi hàng tạo thành
một khoảng; các biên khoảng thay đổi đơn điệu theo hàng. Mã hóa tổng tiền tố
dùng một trạng thái toàn cục cho mọi cột. Ladder AMK chỉ vật hóa những ngưỡng
tiền tố/hậu tố mà các hàng thật sự đọc. Do đó tiêu chí lựa chọn không chỉ là
\(n,k\), mà còn là độ phức tạp của hai đường biên bậc thang.

Với nhiều độ rộng và nhiều cận, một ôtômát có thể nắm toàn bộ lịch sử nhưng số
trạng thái tăng theo tổ hợp các bộ đếm đang hoạt động. Số ngưỡng dùng chung tăng
theo số biên và mức đếm được yêu cầu. Hai hướng vì thế tạo một đánh đổi:
ôtômát hợp nhất các ràng buộc trong trạng thái, còn Ladder AMK giữ các ràng
buộc tách biệt nhưng dùng chung đại lượng trung gian. Đây là câu hỏi thực
nghiệm và lý thuyết đáng quan tâm hơn việc chỉ so số mệnh đề trên một cấu hình.

\section{Đúng đắn}

\begin{theorem}[Tính đúng của phép mã hóa cửa sổ dùng chung]
\label{thm:shared-window-correct}
Giả sử:
\begin{enumerate}
  \item các thanh ghi tiến/lùi đúng theo đặc tả ngưỡng;
  \item mọi cửa sổ nằm trong một khối được ràng buộc bởi bộ đếm cục bộ;
  \item mỗi cửa sổ hoặc nằm trong một khối, hoặc cắt đúng một biên khối;
  \item với cửa sổ \(A\cup B\) cắt biên, phép mã hóa chứa đủ các mệnh đề
  \[
    \neg T(A,q)\lor\neg T(B,k+1-q),
    \qquad q=0,\ldots,k+1.
  \]
\end{enumerate}
Khi đó \CNF tương đương theo phép chiếu với họ staircase AMK.
\end{theorem}

\begin{proof}
Để chứng minh tính đúng, xét một cửa sổ. Nếu nằm trong khối, bộ đếm cục bộ loại mọi tổng
vượt \(k\). Nếu cắt biên, đặt \(a=\sum_{x\in A}x\),
\(b=\sum_{x\in B}x\). Khi \(a+b\ge k+1\), chọn
\(q=\min\{a,k+1\}\). Ta có \(T(A,q)=1\); nếu \(q<k+1\) thì
\(b\ge k+1-a=k+1-q\), còn nếu \(q=k+1\) thì ngưỡng bên phải bằng \(0\)
và đúng theo quy ước. Bộ nối tương ứng bị vi phạm.

Để chứng minh tính đầy đủ, lấy một chuỗi thỏa mọi cửa sổ và gán các thanh ghi
bằng đúng tổng tiền tố/hậu tố. Các mệnh đề của bộ đếm đều đúng. Nếu một bộ nối sai thì tồn
tại \(q\) sao cho \(a\ge q\) và \(b\ge k+1-q\), suy ra \(a+b\ge k+1\),
mâu thuẫn với giả thiết cửa sổ nhiều nhất \(k\). Do đó phép gán trên biến gốc
luôn mở rộng được thành mô hình của \CNF; kết hợp với tính đúng cho tương đương
theo phép chiếu.
\end{proof}

\subsection{Cửa sổ cắt nhiều biên}

Định lý~\ref{thm:shared-window-correct} bao phủ trực tiếp cửa sổ trong một khối
và cửa sổ cắt một biên. Với cửa sổ đi qua ba khối trở lên, tổng được ghép từ
nhiều hơn hai thành phần.
Có thể ghép tuần tự các khối bằng một bộ đếm trung gian, hoặc dùng một cây ghép
có đầu ra bị cắt ở \(k+1\). Với các tổng bộ phận
\(S_1,\ldots,S_r\), bất biến cần giữ là
\[
  U_{j,q}\iff \min\!\left(k+1,\sum_{h=1}^{j}S_h\right)\ge q .
\]
Sau lần ghép cuối, mệnh đề \(\neg U_{r,k+1}\) đóng AMK. Chi phí tăng theo số
khối mà từng cửa sổ chạm tới; một cây ghép riêng cho mỗi cửa sổ làm giảm tỷ
lệ trạng thái dùng chung. Điều kiện ``mỗi cửa sổ cắt nhiều nhất một biên'' do
đó là tiêu chí phân hoạch, còn bộ đếm ghép là phương án cho phân hoạch tổng
quát.

\section{Chọn độ rộng khối}

Khối nhỏ tạo nhiều biên và bộ nối nhưng bộ đếm cục bộ ngắn. Khối lớn giảm số
biên nhưng giữ nhiều thanh ghi không được mọi cửa sổ dùng. Vì vậy độ rộng là
một siêu tham số cấu trúc.

Ba cách chọn:
\begin{itemize}
  \item bằng độ rộng cửa sổ khi mọi ràng buộc có cùng \(w\);
  \item theo ước chung lớn nhất hoặc cụm của nhiều độ rộng;
  \item phân hoạch thích nghi, đặt biên tại nơi có ít cửa sổ cắt qua.
\end{itemize}

Phân hoạch thích nghi có thể xem là bài toán tối ưu trên đồ thị khoảng: mỗi cửa
sổ là một khoảng, chi phí của một biên bằng số ràng buộc cắt biên nhân với số
ngưỡng cần nối.

\section{Triển khai}

Một giao diện lập trình (API) cho mô-đun có thể là:
\[
  \texttt{SharedCounter(sequence, partition, maxThreshold)}
\]
và trả về:
\begin{itemize}
  \item các literal ngưỡng tiền tố/hậu tố;
  \item hàm \texttt{prefix(block,pos,q)};
  \item hàm \texttt{suffix(block,pos,q)};
  \item ánh xạ thanh ghi để gỡ lỗi và kiểm tra nội bộ.
\end{itemize}
Sau đó \texttt{addAtMostWindow(i,w,k)} chỉ sinh các bộ nối. Nhiều mô-đun sử
dụng có thể đọc cùng một bộ đếm: AMK, ALK, EXK hoặc hàm mục tiêu.

\begin{algorithm}[H]
\caption{Sinh AMK cửa sổ dùng chung khi mỗi cửa sổ cắt nhiều nhất một biên}
\label{alg:shared-window}
\begin{minipage}{.94\textwidth}\small
\textbf{Đầu vào:} chuỗi \(x_1,\ldots,x_n\); một phân hoạch khối; tập cửa sổ
\(\mathcal W\); cận \(k(W)\) của từng cửa sổ.\\
\textbf{Đầu ra:} một \CNF tương đương theo phép chiếu.
\begin{enumerate}
  \item Trong mỗi khối, sinh các thanh ghi ngưỡng tiến/lùi đến mức lớn
  nhất \(\max_{W\in\mathcal W}(k(W)+1)\) thật sự được đọc.
  \item Với cửa sổ nằm trọn trong một khối, đọc bộ đếm cục bộ và cấm
  ngưỡng \(k(W)+1\).
  \item Với cửa sổ \(W=A\cup B\) cắt một biên, lần lượt
  \(q=0,\ldots,k(W)+1\), sinh
  \(\neg T(A,q)\lor\neg T(B,k(W)+1-q)\); giản lược các literal hằng.
  \item Nếu phát hiện một cửa sổ cắt từ hai biên trở lên, dừng mô-đun này
  và chuyển sang bộ đếm ghép; không áp dụng bộ nối hai ngôi.
  \item Trả \CNF cùng ánh xạ từ mỗi thanh ghi về
        \((\text{khối},\text{vị trí},q)\).
\end{enumerate}
\end{minipage}
\end{algorithm}

\subsection{Miền hiệu quả của bộ đếm dùng chung}

Bộ đếm dùng chung đạt lợi thế rõ nhất khi các cửa sổ chồng lấn mạnh, cận \(k\)
nhỏ và mỗi cửa sổ chỉ cắt ít khối. Với các cửa sổ gần rời nhau hoặc bài toán
nhỏ, mã hóa từng cặp và bộ đếm độc lập thường có hằng số thấp hơn. Một bộ chọn
có thể quyết định trước khi sinh \CNF từ ba đại lượng: tổng số ngưỡng cục bộ,
số bộ nối sau giản lược và số bước ghép của các cửa sổ nhiều biên.

\section{Bài học thiết kế}

\begin{summarybox}
\begin{enumerate}
  \item Cửa sổ chồng lấn là tín hiệu nên dùng trạng thái chung.
  \item Bộ đếm cục bộ cung cấp các ngưỡng; bộ nối mô tả ràng buộc.
  \item Thanh ghi tiến và lùi tách một cửa sổ cắt biên thành hai phần.
  \item Chứng minh tính đúng tách cửa sổ nội khối khỏi cửa sổ cắt biên.
  \item Độ rộng khối cân bằng chi phí thanh ghi với chi phí bộ nối.
\end{enumerate}
\end{summarybox}

\subsection*{Câu hỏi nghiên cứu}
\begin{enumerate}
  \item Dựng các bộ nối cho một cửa sổ AMK với \(k=2\) cắt biên thành hai phần.
  \item Tính số thanh ghi khi \(n=20,w=5,k=3\) và độ rộng khối bằng \(5\).
  \item Đề xuất hàm mục tiêu cho phân hoạch thích nghi dựa trên số cửa sổ cắt biên.
\end{enumerate}

\chapter{Thanh ghi đếm thích nghi cho
  \texorpdfstring{ràng buộc đúng \(K\)}{ràng buộc đúng K}}
\chapterlead{Bộ đếm tuần tự thông thường dành cùng một độ rộng cho mọi
hàng. Thanh ghi thích nghi chỉ tạo những trạng thái đếm có thể đạt tới, nhờ
đó giữ được cấu trúc tuần tự nhưng loại bỏ phần tam giác trạng thái bất khả.}

\index{exactly-k}\index{New Sequential Counter}
\index{adaptive counter}\index{Social Golfer}

\flowdiagram{Dãy Boolean}{Miền trạng thái khả đạt}{Truy hồi ngưỡng}
  {Điều kiện biên}{Đúng \(K\)}

\section{Ràng buộc đúng \texorpdfstring{\(K\)}{K} trong không gian thiết kế}

\emph{Exactly-\(K\)} (ràng buộc đúng \(K\)) có nhiều cách phân rã đúng:
\[
  \sum X=K
  \iff
  \left(\sum X\le K\right)\land
  \left(\sum X\ge K\right),
\]
nhưng hai vế không nhất thiết phải dùng hai bộ biểu diễn độc lập. Một vectơ đơn
phân đầy đủ cho phép đọc cùng lúc \(o_K\) và \(\neg o_{K+1}\); BDD giữ các
trạng thái tổng; mạng sắp xếp các đầu vào; bộ đếm tuần tự theo dõi tiền tố. Câu
hỏi của chương là cụ thể: nếu đã chọn các trạng thái ngưỡng tuần tự thì những
ô nào thật sự có thể đạt và cần được hiện thực hóa (reification)?

\begin{center}
\captionof{table}{Các hướng biểu diễn ràng buộc đúng \(K\).}
\label{tab:exactly-k-designs}
\begin{tabularx}{\textwidth}{@{}>{\bfseries\sffamily}p{.19\textwidth}YYY@{}}
\toprule
Họ & Trạng thái & Giao diện EXK & Đặc trưng\\
\midrule
Hai AMK & Hai bộ đếm & Hai cận tách & Dễ ghép nhưng có thể lặp trạng thái\\
Tuần tự & Tổng tiền tố & Chống tràn + đạt \(K\) & Tốt khi \(K\) nhỏ\\
Cây tổng & Tổng cây con & \(o_K\land\neg o_{K+1}\) & Dùng lại nhiều cận\\
Mạng & Vị trí đã sắp & Hai đầu ra kề nhau & Cấu trúc đều, độ sâu thấp\\
BDD & \((i,s)\) & Chấp nhận trạng thái \(K\) & Hợp nhất trạng thái theo tổng\\
NSC & Ngưỡng tiền tố khả đạt & Biên trên + dưới & Cắt vùng tam giác bất khả\\
\bottomrule
\end{tabularx}
\end{center}

Các cận dưới cho AMO và những khảo cứu về ràng buộc đếm đặt số biến bên cạnh
số mệnh đề, literal và mức lan truyền
\parencite{kucera2019lower,bailleux2010survey}. NSC giữ bậc \(O(nK)\), giảm
một phần xác định của bảng và thay cách đóng biên Exactly-\(K\). Đóng góp cốt
lõi là một \emph{không gian trạng thái gọn hơn trong họ tuần tự}; vị trí trên
biên Pareto (Pareto frontier) được đo bằng số mệnh đề, số literal và sức lan truyền.

\section{Từ bảng chữ nhật đến miền trạng thái khả đạt}

Xét dãy Boolean \(X=\langle x_1,\ldots,x_n\rangle\) và ràng buộc
\[
  \sum_{i=1}^{n}x_i=k.
\]
Một bộ đếm tuần tự có thể dành \(k\) cột cho mọi tiền tố. Tuy nhiên, sau
\(i\) biến đầu tiên, tổng không thể lớn hơn \(i\). Các ô \(j>i\) vì thế không
mang trạng thái hợp lệ.

\emph{New Sequential Counter} (bộ đếm tuần tự mới, NSC) dùng thanh ghi tam giác
\[
  R_{i,j}=1
  \quad\Longleftrightarrow\quad
  \sum_{q=1}^{i}x_q\ge j,
  \qquad
  \begin{matrix}
    1\le i\le n-1,\\[-1mm]
    1\le j\le \min(i,k).
  \end{matrix}
\]
Chỉ \(n-1\) tiền tố được hiện thực hóa (reification); biến \(x_n\) được xử lý trong khi giải (in-processing) khi giải (in-processing) điều kiện cuối.
Với \(n>k\), số biến phụ là
\[
  \sum_{i=1}^{n-1}\min(i,k)
  =
  k(n-1)-\frac{k(k-1)}{2}.
\]
So với bảng chữ nhật \(k(n-1)\), phần tiết kiệm chính xác là
\(k(k-1)/2\) biến. Bậc tiệm cận được giữ nguyên, còn hệ số hữu hạn giảm đúng các
trạng thái bất khả mà bộ giải phải duy trì.

\begin{figure}[tb]
\centering
\begin{tikzpicture}[satfig,x=7mm,y=7mm]
  \node[satlabel] at (2.5,1.25) {bảng chữ nhật};
  \node[satlabel] at (9.5,1.25) {miền khả đạt của NSC};
  \foreach \i in {1,...,5}{
    \foreach \j in {1,...,3}{
      \draw[Rule,fill=SoftGray] (\i,-\j) rectangle ++(.82,.82);
      \node[font=\sffamily\tiny] at (\i+.41,-\j+.41) {\(R_{\i,\j}\)};
      \ifnum\j>\i
        \fill[Gold!28] (\i,-\j) rectangle ++(.82,.82);
        \draw[Gold,thick] (\i,-\j)--++(.82,.82);
        \draw[Gold,thick] (\i+.82,-\j)--++(-.82,.82);
      \fi
    }
  }
  \foreach \i in {1,...,5}{
    \foreach \j in {1,...,3}{
      \ifnum\j>\i
      \else
        \draw[Blue,fill=PaleBlue] ({7+\i},-\j) rectangle ++(.82,.82);
        \node[font=\sffamily\tiny] at ({7+\i+.41},-\j+.41) {\(R_{\i,\j}\)};
      \fi
    }
  }
  \node[satlabel,text=Gold,align=center] at (2.8,-4)
    {các ô \(j>i\)\\không thể đạt};
  \node[satlabel,text=Blue,align=center] at (10.1,-4)
    {\(12\) ô thay vì \(15\)\\khi \(n=6,k=3\)};
\end{tikzpicture}
\caption{Bộ đếm tuần tự hình chữ nhật và miền trạng thái tam giác. NSC không
tạo ba ô bất khả \(R_{1,2},R_{1,3},R_{2,3}\).}
\label{fig:nsc-triangular-domain}
\end{figure}

\begin{keyidea}{``Thích nghi'' ở đây có nghĩa gì?}
Độ rộng hàng \(i\) là \(\min(i,k)\), được quyết định bởi số trạng thái có thể
đạt sau \(i\) đầu vào. ``Thích nghi'' ở đây mô tả hình dạng tĩnh của bộ đếm,
được xác định trước khi bộ giải chạy.
\end{keyidea}

\section{Đặc tả ngữ nghĩa và hệ thức truy hồi}

Dùng hai quy ước biên
\[
  R_{i,0}\equiv\top,
  \qquad
  R_{0,j}\equiv\bot\quad(j\ge1).
\]
Khi đó mọi ô hợp lệ tuân theo
\begin{equation}
  R_{i,j}
  \;\longleftrightarrow\;
  R_{i-1,j}\lor
  \bigl(x_i\land R_{i-1,j-1}\bigr).
  \label{eq:nsc-recurrence}
\end{equation}
Vế phải mô tả đúng hai cách để tiền tố mới đạt ngưỡng \(j\): tiền tố cũ đã đạt
ngưỡng, hoặc \(x_i=1\) và tiền tố cũ đã đạt \(j-1\).

Trong CNF, hệ thức truy hồi được phân rã thành các quan hệ kéo theo theo hai chiều. Phần
tiến gồm
\[
  R_{i-1,j}\rightarrow R_{i,j},
  \qquad
  (x_i\land R_{i-1,j-1})\rightarrow R_{i,j}.
\]
Phần lùi gồm
\[
  (\neg x_i\land\neg R_{i-1,j})\rightarrow\neg R_{i,j},
  \qquad
  \neg R_{i-1,j-1}\rightarrow\neg R_{i,j}.
\]
Sau khi thay các hằng biên và bỏ literal thừa, ta thu được các mệnh đề hai hoặc
ba literal. Cách sinh mệnh đề trực tiếp theo miền tam giác tránh tạo biến rồi
giao cho bộ tiền xử lý xóa đi.

\begin{designrule}{Thanh ghi được đọc như một đại lượng ngữ nghĩa}
Nếu \(R_{i,j}\) xuất hiện trong điều kiện biên, hàm mục tiêu hoặc một ràng buộc
khác, hệ thức truy hồi phải bảo đảm đặc tả hai chiều. Chỉ các quan hệ kéo theo
tiến không đủ để diễn giải \(\neg R_{i,j}\) là ``tiền tố chưa đạt \(j\)''.
\end{designrule}

\section{Đóng ràng buộc \texorpdfstring{đúng \(K\)}{đúng K}}

Phần nhiều nhất \(k\) cấm một đầu vào mới làm tràn bộ đếm:
\begin{equation}
  x_i\rightarrow\neg R_{i-1,k},
  \qquad i=k+1,\ldots,n.
  \label{eq:nsc-overflow}
\end{equation}
Phần ít nhất \(k\) yêu cầu tổng cuối đạt ngưỡng \(k\):
\begin{equation}
  R_{n-1,k}\lor
  \bigl(x_n\land R_{n-1,k-1}\bigr).
  \label{eq:nsc-lower}
\end{equation}
Biểu thức \cref{eq:nsc-lower} được phân phối thành CNF trước khi xuất. Khi
\(k=1\), \(R_{n-1,0}\equiv\top\); khi \(k=n\), có thể đơn giản hóa toàn bộ
ràng buộc thành các mệnh đề đơn vị \(x_i\).

\begin{theorem}[Đúng đắn của NSC]
Hệ thức truy hồi \cref{eq:nsc-recurrence} cùng các điều kiện
\cref{eq:nsc-overflow,eq:nsc-lower} tương đương theo phép chiếu với
\(\sum_{i=1}^{n}x_i=k\).
\end{theorem}

\begin{proof}
Hệ thức truy hồi và các hằng biên cho, bằng quy nạp theo \(i\),
\[
  R_{i,j}=1
  \Longleftrightarrow
  \sum_{q=1}^{i}x_q\ge j.
\]
Nếu có hơn \(k\) đầu vào đúng, xét đầu vào đúng thứ \(k+1\), ký hiệu \(x_i\).
Ngay tại vị trí trước đó, \(R_{i-1,k}=1\), trái với
\cref{eq:nsc-overflow}. Do đó tổng không vượt \(k\).

Ngược lại, \cref{eq:nsc-lower} chỉ đúng khi \(k\) đầu vào đầu tiên đã xuất hiện
trong \(n-1\) vị trí đầu, hoặc \(x_n=1\) và tiền tố trước đó chứa ít nhất
\(k-1\) đầu vào đúng. Vì vậy tổng ít nhất \(k\). Hai cận cho tổng đúng bằng
\(k\). Với mọi vectơ \(X\) có tổng \(k\), gán \(R_{i,j}\) theo tổng tiền tố
thật thỏa hệ thức truy hồi và cả hai điều kiện biên.
\end{proof}

\begin{workedexample}{Exactly-\(3\) trên sáu biến}
Với \(n=6,k=3\), NSC tạo
\[
  \min(1,3)+\min(2,3)+\cdots+\min(5,3)=12
\]
biến thanh ghi, thay vì \(5\cdot3=15\) ô của bảng chữ nhật. Xét
\[
  x=(1,0,1,0,0,1).
\]
Sau năm biến, tổng tiền tố bằng \(2\), nên
\[
  (R_{5,1},R_{5,2},R_{5,3})=(1,1,0).
\]
Điều kiện cuối
\(
R_{5,3}\lor(x_6\land R_{5,2})
\)
đúng chính vì \(x_6=1\) và tiền tố đã đạt ngưỡng \(2\). Nếu đặt thêm
\(x_5=1\), ta có \(R_{5,3}=1\); mệnh đề chống tràn
\(x_6\rightarrow\neg R_{5,3}\) lập tức buộc \(x_6=0\). Hai tình huống này
minh họa lần lượt biên dưới và biên trên của cùng một thanh ghi.
\end{workedexample}

\section{Cơ chế suy diễn lan truyền hai chiều}

Bộ đếm hữu ích vì các cận xuất hiện sớm trong tìm kiếm. Hai tình huống quan
trọng có thể kiểm tra trực tiếp:

\begin{enumerate}
  \item Khi đã có \(k\) đầu vào đúng trong một tiền tố, thanh ghi mức \(k\) của
  tiền tố đó trở thành đúng. Các mệnh đề chống tràn buộc mọi đầu vào còn lại về \(0\).
  \item Khi số vị trí có thể còn nhận giá trị \(1\) vừa đủ để đạt \(k\), điều
  kiện cận dưới và các quan hệ kéo theo lùi lần lượt loại nhánh đặt một vị trí còn
  lại về \(0\).
\end{enumerate}

Hai hướng lan truyền này tương ứng với hai biên ``đã đủ'' và ``phải dùng hết
phần còn lại''. Chúng nên được kiểm tra bằng phép thử đơn vị trên mọi phép gán
bộ phận nhỏ, thay vì suy luận hiệu quả chỉ từ số biến hoặc số mệnh đề.

\section{Dùng một thanh ghi cho nhiều giao diện}

Khi thanh ghi đã có đặc tả ngưỡng, nhiều ràng buộc có thể đọc cùng một mô-đun:
\[
  \sum X\le k
  \quad\Longrightarrow\quad
  \text{các mệnh đề chống tràn},
\]
\[
  \sum X\ge k
  \quad\Longrightarrow\quad
  \text{đầu ra đạt ngưỡng cuối},
\]
\[
  \sum X=k
  \quad\Longrightarrow\quad
  \text{kết hợp hai giao diện trên}.
\]
Nếu cần nhiều cận \(k_1<k_2<\cdots\), bộ đếm nên được tạo đến ngưỡng lớn nhất
rồi cung cấp các đầu ra nhỏ hơn. Cách này tránh xây hai mạng đếm độc lập cho
ràng buộc và hàm mục tiêu.

\section{Nghiên cứu tình huống ngắn: xếp nhóm chơi golf}

\subsection{Từ thiết kế tổ hợp đến mô hình SAT}

\emph{Social Golfer} (bài toán xếp nhóm người chơi golf) vừa là bài toán lập
lịch vừa là bài toán thiết kế tổ hợp. Ba
nguồn khó đan xen với nhau: ràng buộc đúng kích thước nhóm, cấm lặp cặp người
chơi và một nhóm đối xứng lớn do hoán vị tuần, nhóm, người. Vì vậy, hiệu quả
của một phép biểu diễn ràng buộc đếm trên SGP không thể tách khỏi lựa chọn biến
chính và cách phá đối xứng.

Các mô hình SAT trước đó từng dùng thêm chỉ số vị trí trong nhóm, sau đó được
sửa và rút gọn bằng biến người--nhóm--tuần cùng cấu trúc bậc thang
\parencite{triska2012golfer}. Dòng phát triển này minh họa một quy tắc quan
trọng: bỏ một chỉ số không mang nghĩa thường tạo lợi ích lớn hơn tối ưu vài
mệnh đề trong chính bộ biểu diễn ràng buộc đếm. Phép đánh giá NSC vì thế bắt đầu sau
khi phần biểu diễn đối tượng đã được chuẩn hóa.

Một bài toán thử \(g\)-\(p\)-\(w\) có \(n=gp\) người chơi, \(g\) nhóm mỗi tuần,
mỗi nhóm có \(p\) người và \(w\) tuần. Biến ba chỉ số
\[
  G_{i,r,t}=1
  \quad\Longleftrightarrow\quad
  \text{người chơi \(i\) thuộc nhóm \(r\) ở tuần \(t\)}
\]
loại bỏ chỉ số vị trí bên trong nhóm, bởi thứ tự các thành viên trong một nhóm
không mang ý nghĩa.

Mô hình gồm ba họ ràng buộc:
\begin{align*}
  &\sum_{r=1}^{g}G_{i,r,t}=1
    &&\text{với mọi người chơi \(i\), tuần \(t\)},\\
  &\sum_{i=1}^{n}G_{i,r,t}=p
    &&\text{với mọi nhóm \(r\), tuần \(t\)},\\
  &\neg G_{i,r,t}\lor\neg G_{j,r,t}\lor
    \neg G_{i,s,u}\lor\neg G_{j,s,u}
    &&\text{với \(i<j,\ t<u\)}.
\end{align*}
Họ thứ ba cấm một cặp người chơi gặp nhau ở cả nhóm \(r\), tuần \(t\) và nhóm
\(s\), tuần \(u\). NSC được dùng ở họ thứ hai; mỗi cặp nhóm--tuần có một
bộ đếm tam giác riêng \parencite{kieu2026golfer}.

Hai phép chuẩn hóa rẻ được thêm vào mô hình. Tuần đầu được cố định thành các
khối người chơi liên tiếp; ở các tuần sau, một số người chơi neo được gán vào
các nhóm đã định. Chúng loại các hoán vị hiển nhiên của người chơi, nhóm và
tuần mà không thêm biến phụ.

\begin{figure}[tb]
\centering
\begin{tikzpicture}[satfig,x=17mm,y=9mm]
  \node[satlabel] at (0,1.0) {tuần};
  \foreach \g in {1,2,3}
    \node[satlabel] at (\g,1.0) {nhóm \(\g\)};
  \node[satlabel] at (0,0) {\(1\)};
  \node[satlabel] at (0,-1) {\(2\)};
  \foreach \g/\txt in {1/{1,2},2/{3,4},3/{5,6}}{
    \node[satbox,minimum width=15mm,fill=PaleBlue] at (\g,0) {\(\{\txt\}\)};
  }
  \foreach \g/\txt in {1/{1,3},2/{2,5},3/{4,6}}{
    \node[satbox,minimum width=15mm,fill=PaleTeal] at (\g,-1) {\(\{\txt\}\)};
  }
  \draw[decorate,decoration={brace,amplitude=4pt},Gold,thick]
    (3.6,.42)--(3.6,-1.42)
    node[midway,right=5pt,satlabel,text=Gold,align=left]
      {mỗi người đúng một nhóm/tuần;\\mỗi cặp gặp nhiều nhất một lần};
\end{tikzpicture}
\caption{Một lịch khả thi cho bài toán nhỏ \(3\)-\(2\)-\(2\). Mỗi ô nhóm là
một ràng buộc đúng \(2\); tuần đầu có thể được cố định làm lát chuẩn.}
\label{fig:social-golfer-small}
\end{figure}

\begin{resultbox}{Kết quả trên 66 bài toán Social Golfer}
Trong thí nghiệm của \textcite{kieu2026golfer}, cấu hình dùng NSC giải được
56/66 bài toán; các cấu hình ADDER, BDD và CARD cùng giải được 54, còn
BINOMIAL giải được 32. NSC có số lần hết thời hạn thấp nhất; thứ hạng trên từng
bài toán thay đổi theo cấu trúc của công thức. Kết quả làm rõ lợi ích của việc
cân bằng số biến, số mệnh đề và cấu trúc lan truyền.
\end{resultbox}

\subsection{Phân tích kết quả}

Mô hình ba chỉ số loại đối xứng vị trí, còn trạng thái tuần tự tam giác cung
cấp một phép biểu diễn EXK cạnh tranh cho ràng buộc kích thước nhóm. Có thể đo
riêng tác động của NSC bằng cách giữ cố định mô hình SGP, chỉ thay mô-đun kích
thước nhóm, phân tầng theo \(p\), \(g\), \(w\) và tách các bài SAT/UNSAT.

\section{Cài đặt và kiểm thử}

Một bộ sinh NSC nên nhận dãy literal, cận \(k\) và chế độ
\texttt{AMK}/\texttt{ALK}/\texttt{EXK}. Bảng ánh xạ
\((i,j)\mapsto\texttt{varID}\) chỉ chứa ô hợp lệ. Ba lớp kiểm thử là đủ để bắt
phần lớn lỗi biên:

\begin{enumerate}
  \item liệt kê toàn bộ vectơ đầu vào khi \(n\le8\) và so sánh với tổng thật;
  \item với mỗi vectơ hợp lệ, kiểm tra có tồn tại phần mở rộng cho biến \(R\);
  \item với mỗi phép gán một phần, so sánh literal đầu vào mà lan truyền đơn vị
  suy ra với các hệ quả trực tiếp của hai cận.
\end{enumerate}

Các trường hợp \(k=0\), \(k=1\), \(k=n-1\), \(k=n\) cần được kiểm thử riêng. Đây là nơi
quy ước hằng biên và phạm vi chỉ số thường gây lỗi nhất.

\begin{algorithm}[H]
\caption{Sinh NSC cho ràng buộc đúng \(K\)}
\label{alg:nsc}
\begin{minipage}{.94\textwidth}\small
\textbf{Đầu vào:} các literal \(x_1,\ldots,x_n\) và \(0\le k\le n\).\\
\textbf{Đầu ra:} \CNF và ánh xạ ngữ nghĩa của các thanh ghi.
\begin{enumerate}
  \item Xử lý trực tiếp \(k=0\) bằng \(\bigwedge_i\neg x_i\) và \(k=n\)
  bằng \(\bigwedge_i x_i\). Nếu \(k>n/2\), có thể thay
  \(y_i=\neg x_i\) và đếm Exactly-\((n-k)\).
  \item Chỉ cấp phát \(R_{i,j}\) với \(1\le i\le n-1\) và
  \(1\le j\le\min(i,k)\); ghi ánh xạ
  \(R_{i,j}\leftrightarrow[\sum_{q=1}^{i}x_q\ge j]\).
  \item Với từng ô hợp lệ, phân rã tương đương
  \(R_{i,j}\leftrightarrow
    R_{i-1,j}\lor(x_i\land R_{i-1,j-1})\),
  thay \(R_{i,0}=\top\), \(R_{0,j}=\bot\), rồi giản lược các literal hằng.
  \item Sinh các mệnh đề chống tràn của~\eqref{eq:nsc-overflow}.
  \item Sinh điều kiện đạt cận dưới~\eqref{eq:nsc-lower}; trả \CNF và ánh xạ.
\end{enumerate}
\end{minipage}
\end{algorithm}

\section{Miền lựa chọn của NSC}

NSC loại chính xác \(k(k-1)/2\) ô so với bảng chữ nhật, nên tỷ lệ tiết kiệm
biến chỉ là
\[
  \frac{k(k-1)/2}{k(n-1)}
  =\frac{k-1}{2(n-1)}.
\]
Khi \(n\) lớn nhưng \(k\) cố định, tỷ lệ này tiến về \(0\); hằng số của ánh xạ
và các mệnh đề hai chiều có thể quan trọng hơn phần tam giác bị bỏ. Với
\(k\) gần \(n\), đếm các literal bù \(\neg x_i\) đến \(n-k\) thường gọn hơn.
Nếu cùng một tập đầu vào cần nhiều cận hoặc cần thay cận tăng dần, cây tổng và
mạng đếm có đầu ra chia sẻ tự nhiên hơn một NSC chuyên biệt. BDD cũng có thể
thắng khi phép rút gọn hợp nhất được nhiều trạng thái tương đương.

Bộ chọn cho ràng buộc đúng \(K\) dùng \(n\), \(\min(k,n-k)\), số cận sẽ được
truy vấn, nhu cầu giải tăng dần và đặc tả lan truyền làm đặc trưng. NSC phù hợp
với một cận cố định và vùng tam giác bị cắt đủ lớn; cây tổng/mạng đếm phù hợp
khi cần chia sẻ nhiều đầu ra. Kết quả Social Golfer minh họa cách các đặc
trưng này dẫn đến một lựa chọn hiệu quả trên họ bài toán tương ứng.

\section{Bài học thiết kế}

\begin{summarybox}
\begin{enumerate}
  \item Miền trạng thái khả đạt quyết định hình dạng bộ đếm.
  \item Thanh ghi \(R_{i,j}\) có đặc tả ngưỡng rõ ràng, không chỉ là biến
  Tseitin vô danh.
  \item Ràng buộc đúng \(K\) là sự kết hợp của một biên tràn và một điều kiện đạt
  ngưỡng cuối.
  \item Hiệu quả cần được đọc đồng thời qua kích thước, lan truyền và hành vi
  giải trên bài toán khó.
  \item Social Golfer cho thấy lựa chọn biến chính và phép biểu diễn ràng buộc đếm phải
  được thiết kế cùng nhau.
\end{enumerate}
\end{summarybox}

\subsection*{Câu hỏi nghiên cứu}
\begin{enumerate}
  \item Vẽ miền biến \(R_{i,j}\) cho \(n=7,k=3\) và tính phần tiết kiệm so với
  bảng chữ nhật.
  \item Viết các mệnh đề chống tràn cho \(\sum_{i=1}^{6}x_i\le2\).
  \item So sánh hai mô hình Social Golfer có và không có chỉ số vị trí trong
  nhóm về số biến chính và nhóm đối xứng còn lại.
\end{enumerate}

\chapter{Đối xứng, biến quan hệ và liên kết biểu diễn}
\chapterlead{Một phép biểu diễn tốt không chỉ mô tả nghiệm; mô hình còn quyết định cách
bộ giải nhìn thấy các nghiệm tương đương và các xung đột cục bộ.
\emph{Symmetry breaking} (phá đối xứng) thu gọn không gian tìm kiếm, còn
\emph{channeling} (liên kết biểu diễn) nối các tiêu chí bổ trợ của cùng một
quyết định.}

\index{symmetry breaking}\index{lex-leader}
\index{relation variable}\index{order encoding}\index{channeling}

\flowdiagram{Quyết định gốc}{Nhóm đối xứng}{Đại diện chuẩn}
  {Biến quan hệ}{Mệnh đề cục bộ}

\section{Hai cấp đối xứng: mô hình và công thức}

Một đối xứng của bài toán tác động lên đối tượng miền: đổi tên khung chứa, máy,
nhóm, tuần, màu hoặc phản xạ hình học. Một đối xứng của \CNF là hoán vị biến
và literal bảo toàn tập mệnh đề. Hai nhóm này liên quan nhưng không trùng nhau.
Phép biểu diễn có thể:
\begin{itemize}
  \item làm một đối xứng miền biến mất vì không biểu diễn chỉ số thừa;
  \item làm một đối xứng miền trở nên khó phát hiện do biến phụ;
  \item tạo đối xứng phụ mới giữa các trạng thái tương đương;
  \item có đối xứng cú pháp không tương ứng với một phép biến đổi nghiệm dễ
        diễn giải.
\end{itemize}

Các vị từ phá đối xứng tổng quát xây một điều kiện chọn đại diện cho các quỹ
đạo của bài toán tìm kiếm \parencite{crawford1996symmetry}. Ở cấp \CNF, những
hệ như Shatter xây đồ thị tô màu của công thức, tìm tự đẳng cấu rồi sinh các
ràng buộc phá đối xứng \parencite{aloul2003shatter}. Ở cấp mô hình, người thiết
kế đã biết ngữ nghĩa nên có thể viết ràng buộc ngắn hơn, chẳng hạn tiền tố
khung chứa đã dùng. Hai cách bổ sung nhau: phá đối xứng miền được áp dụng trước;
phát hiện trên \CNF tìm phần đối xứng còn sót.

\begin{center}
\captionof{table}{Bốn cấp khai thác đối xứng.}
\label{tab:symmetry-levels}
\begin{tabularx}{\textwidth}{@{}>{\bfseries\sffamily}p{.2\textwidth}YYY@{}}
\toprule
Cấp & Đầu vào & Lợi thế & Hạn chế\\
\midrule
Mô hình & Đối tượng và phép đổi tên & Ràng buộc ngắn, có chứng minh miền &
Cần tri thức bài toán\\
\CNF tĩnh & Đồ thị công thức & Tự động, dùng cho mô hình bất kỳ &
Biến phụ có thể che đối xứng\\
Động & Trạng thái tìm kiếm & Khai thác đối xứng còn lại & Chi phí tích hợp bộ giải\\
Chuẩn hóa & Toàn quỹ đạo & Có thể giữ một đại diện & Ràng buộc lớn/lan truyền sâu\\
\bottomrule
\end{tabularx}
\end{center}

Các lớp ràng buộc phá đối xứng được ghép theo thứ tự. Ràng buộc miền làm thay
đổi nhóm tự đẳng cấu của đồ thị; bộ phá đối xứng tự động vì thế chạy trên công
thức đã có mệnh đề đơn vị và liên kết biểu diễn. Mỗi lớp bảo toàn ít nhất một
đại diện của các quỹ đạo còn lại, còn thí nghiệm phân tích thành phần (ablation study) đo tác động của từng
lớp lên hiệu năng
\parencite{gent2006symmetry}.

\section{Đối xứng như tác động lên nghiệm}

Gọi \(S\) là tập nghiệm khả thi và \(\Gamma\) là một nhóm các phép biến đổi
bảo toàn tính khả thi và giá trị mục tiêu. Hai nghiệm \(s,s'\in S\) tương
đương khi tồn tại \(\gamma\in\Gamma\) sao cho \(s'=\gamma(s)\). Tập
\[
  \operatorname{Orb}(s)=\set{\gamma(s)\suchthat\gamma\in\Gamma}
\]
là quỹ đạo của \(s\).

Phá đối xứng thêm công thức \(B\) sao cho mỗi quỹ đạo còn ít nhất một
đại diện:
\[
  \forall s\in S,\qquad
  \operatorname{Orb}(s)\cap\operatorname{Mod}(B)\ne\varnothing.
\]
Không nhất thiết phải giữ đúng một đại diện. Trong thực hành, một tập ràng buộc
phá đối xứng nhỏ, rẻ và phù hợp với mô hình thường hữu ích hơn việc cố gắng biểu diễn
một dạng chuẩn hoàn chỉnh \parencite{walsh2006symmetry,gent2006symmetry}.

\begin{figure}[tb]
\centering
\begin{tikzpicture}[satfig,node distance=17mm]
  \node[satstate,sattrue] (s1) {\(s_1\)};
  \node[satstate,right=of s1] (s2) {\(s_2\)};
  \node[satstate,below=of s2] (s3) {\(s_3\)};
  \node[satstate,left=of s3] (s4) {\(s_4\)};
  \draw[satdash] (s1)--node[above,satlabel]{đổi tên}(s2);
  \draw[satdash] (s2)--node[right,satlabel]{phản xạ}(s3);
  \draw[satdash] (s3)--node[below,satlabel]{đổi tên}(s4);
  \draw[satdash] (s4)--node[left,satlabel]{phản xạ}(s1);
  \node[draw=Teal,very thick,rounded corners=4pt,fit=(s1),inner sep=4pt,
        label={[satlabel,text=Teal]above left:đại diện chuẩn}] {};
  \node[draw=Rule,dashed,rounded corners=7pt,fit=(s1)(s2)(s3)(s4),
        inner sep=8pt,label={[satlabel]below right:\(\operatorname{Orb}(s_1)\)}] {};
\end{tikzpicture}
\caption{Một quỹ đạo gồm bốn cách biểu diễn tương đương. Ràng buộc phá đối
xứng hợp lệ giữ ít nhất một nút; phép chuẩn hóa mạnh hơn sẽ cố giữ đúng một nút.}
\label{fig:symmetry-orbit}
\end{figure}

\begin{designrule}{Chứng minh theo phép biến đổi}
Chứng minh một ràng buộc phá đối xứng bằng phép hoán vị hoặc phép phản xạ đưa mọi nghiệm bị
loại sang một nghiệm giữ lại. Ánh xạ này đồng thời xác định nhóm đối xứng và
đại diện được bảo toàn.
\end{designrule}

\section{Loại đối xứng ngay từ lựa chọn biến}

Biến chính có thể tạo hoặc loại đối xứng. Trong Social Golfer, biến
\(G_{i,r,t}\) chỉ ghi người chơi, nhóm và tuần. Nếu thêm chỉ số vị trí trong
nhóm, mỗi nhóm \(p\) người tạo thêm \(p!\) cách biểu diễn cùng một tập.
Trong bài toán xếp khối (packing problem), biến chiếm chỗ theo từng ô có thể tạo nhiều cách diễn
giải một hình chữ nhật; biến tọa độ và hướng đặt biểu diễn hình học trực tiếp hơn.

Nguyên tắc chung là:
\[
  \text{không biểu diễn thứ tự nếu thứ tự không thuộc nghiệm.}
\]
Sau bước chọn biến, các ràng buộc phá đối xứng xử lý nhóm đối xứng còn lại.
Cách tiếp cận đó giảm cả số biến lẫn số mệnh đề; bước phá đối xứng tiếp tục
giảm số mô hình.

\section{Bốn mẫu phá đối xứng}

\subsection{Cố định một lát chuẩn}

Khi các lớp thời gian hoặc tài nguyên hoán vị tự do, có thể cố định lát đầu.
Với Social Golfer, tuần đầu được chuẩn hóa:
\[
  G_{i,r,1}=1
  \quad\text{khi}\quad
  (r-1)p<i\le rp.
\]
Mọi lịch khả thi có thể đổi tên người chơi và nhóm để thỏa cách sắp này. Những
Các mệnh đề đơn vị đó loại một phần lớn quỹ đạo mà hầu như không tăng chi phí giải
\parencite{kieu2026golfer}.

\subsection{Tiền tố sử dụng và thứ tự ưu tiên giá trị}

Giả sử \(u_b\) cho biết khung chứa \(b\) được dùng. Vì tên các khung chứa có thể hoán vị,
áp đặt
\[
  u_{b+1}\rightarrow u_b.
\]
Các khung chứa được dùng tạo thành một tiền tố. Với các giá trị đối xứng, thứ
tự ưu tiên giá trị yêu cầu lần xuất hiện đầu của giá trị nhỏ hơn đứng trước lần
xuất hiện đầu của giá trị lớn hơn. Hai mẫu này cùng chọn thứ tự theo ``lần sử
dụng đầu'', nhưng tiền tố sử dụng chỉ cần một dãy Boolean nên thường rẻ hơn.

\begin{proposition}[Tiền tố sử dụng không làm mất lớp nghiệm]
\label{prop:usage-prefix}
Giả sử các khung chứa có cùng miền, cùng sức chứa và hàm mục tiêu không phụ
thuộc vào tên khung. Với biến gán \(a_{ib}\) và biến sử dụng \(u_b\), mọi
nghiệm đều có một nghiệm cùng giá trị mục tiêu thỏa
\(u_{b+1}\rightarrow u_b\).
\end{proposition}

\begin{proof}
Cho \(U=\set{b:u_b=1}\) và \(q=\card{U}\). Chọn một song ánh
\(\pi:U\to\set{1,\ldots,q}\), rồi mở rộng \(\pi\) thành một hoán vị trên toàn
bộ tập khung chứa. Định nghĩa
\[
  a'_{i,\pi(b)}=a_{ib},
  \qquad
  u'_{\pi(b)}=u_b .
\]
Mọi biến thuộc bản thân vật \(i\), như tọa độ, phép quay hay kích thước, được
giữ nguyên. Vì sức chứa và các ràng buộc gán được phát biểu đồng nhất cho mọi
khung, đổi tên theo \(\pi\) bảo toàn tính khả thi; vì hàm mục tiêu không đọc
chỉ số khung, giá trị mục tiêu cũng không đổi. Trong nghiệm mới,
\(u'_1=\cdots=u'_q=1\) và \(u'_{q+1}=\cdots=0\), nên toàn bộ implication
\(u'_{b+1}\rightarrow u'_b\) đúng.
\end{proof}

Chứng minh trên cũng chỉ rõ điều kiện không thỏa: nếu khung chứa có sức chứa, chi phí hoặc
miền khác nhau thì hoán vị tùy ý không còn là đối xứng. Khi đó chỉ được áp
dụng tiền tố trong từng lớp khung chứa thật sự đồng nhất.

\subsection{Lex-leader}

Với hai vector đại diện các đối tượng hoán vị được,
\[
  A=\langle a_1,\ldots,a_m\rangle,\qquad
  B=\langle b_1,\ldots,b_m\rangle,
\]
ràng buộc \(A\le_{\mathrm{lex}}B\) chọn thứ tự từ điển. Một phép biểu diễn dùng
các biến bằng nhau trên tiền tố \(e_i\):
\[
  e_i\leftrightarrow e_{i-1}\land(a_i\leftrightarrow b_i),
  \qquad e_0=\top,
\]
và cấm \(e_{i-1}\land a_i\land\neg b_i\). \emph{Lex-leader} (ràng buộc dẫn
đầu theo thứ tự từ điển) mạnh hơn tiền tố sử dụng nhưng tạo thêm một chuỗi biến
phụ cho mỗi cặp được so sánh.

\subsection{Neo hình học và phản xạ}

Với khung chứa hình chữ nhật, các phép phản xạ hoặc quay toàn bố cục có thể bảo
toàn nghiệm. Một vật chuẩn---thường là vật lớn nhất theo diện tích rồi theo
cạnh dài---được neo vào nửa trái hoặc một góc chuẩn:
\[
  2x_a+w_a\le W.
\]
Điều kiện phải khớp với nhóm đối xứng thật. Nếu khung chứa, miền quay hoặc hàm
mục tiêu không bất biến qua phản xạ thì ràng buộc này không hợp lệ.

\begin{keyidea}{Ràng buộc mạnh nhất không mặc nhiên là tốt nhất}
Mỗi ràng buộc phá đối xứng đổi cấu trúc mệnh đề, thứ tự quyết định và các xung
đột học được. Nên so sánh ít nhất ba cấu hình: không phá đối xứng, ràng buộc rẻ
và ràng buộc mạnh. Kết luận dựa trên số bài toán giải được và thời gian, không
chỉ trên số mô hình bị
loại.
\end{keyidea}

\section{Biến quan hệ cho các phép tuyển}

Nhiều bài toán có ràng buộc dạng
\[
  A\lor B\lor C\lor D.
\]
Trong bài toán xếp khối, hai hình \(i,j\) không chồng lấn khi ít nhất một quan hệ đúng:
\[
  L_{ij}\lor L_{ji}\lor B_{ij}\lor B_{ji},
\]
trong đó \(L_{ij}\) nghĩa là \(i\) nằm trái \(j\), còn \(B_{ij}\) nghĩa là
\(i\) nằm dưới \(j\). Trong lập lịch, hai tác vụ dùng cùng tài nguyên cần
\[
  P_{ij}\lor P_{ji},
\]
với \(P_{ij}\) nghĩa là \(i\) kết thúc trước khi \(j\) bắt đầu.

\emph{Witness variable} (biến làm chứng) là biến chỉ ra nhánh nào của phép
tuyển được chọn. Trong mô hình này, biến quan hệ đóng vai trò biến làm chứng
cho một nhánh của phép tuyển. Chúng rút
ràng buộc toàn cục thành các mệnh đề cục bộ quanh cặp đối tượng. Khi một quan
hệ được chọn, lựa chọn này kích hoạt một quan hệ suy diễn kéo theo trên tọa độ hoặc thời gian; khi
miền tọa độ loại quan hệ đó, mệnh đề tuyển buộc bộ giải chọn quan hệ còn lại.

\subsection{Biến quan hệ và giá trị hỗ trợ}

Từ góc nhìn chuyển CSP sang SAT, biến quan hệ nằm giữa biểu diễn trực tiếp và mã
hóa hỗ trợ. Mô hình không chọn toàn bộ giá trị của hai biến; phép biểu diễn chỉ chọn một
\emph{lớp} cặp giá trị,
chẳng hạn ``trái'', ``dưới'' hoặc ``trước''. Mỗi lớp sau đó được liên kết sang
các literal biên của miền. Đây là một phép phân tích thành nhân tử:
\[
  \text{bảng cặp giá trị hợp lệ}
  =
  \bigcup_{q\in Q}\text{lớp hỗ trợ}(q).
\]
Nếu \(|Q|\) nhỏ và mỗi lớp được mô tả bằng vài bất đẳng thức, các biến quan hệ
nén một bảng cấm lớn. Nếu các lớp chồng lấn quá nhiều hoặc không
có mô tả ngắn, biến quan hệ chỉ thêm một tầng trung gian.

\begin{figure}[tb]
\centering
\begin{tikzpicture}[satfig,x=8mm,y=7mm]
  \node[satlabel] at (1.7,3.2) {\(L_{ij}\)};
  \node[satlabel] at (6.7,3.2) {\(L_{ji}\)};
  \node[satlabel] at (1.65,1.3) {\(B_{ij}\)};
  \node[satlabel] at (6.65,1.3) {\(B_{ji}\)};
  \draw[Blue,fill=PaleBlue,thick] (0,2) rectangle (1.5,2.8);
  \node at (.75,2.4) {\(i\)};
  \draw[Teal,fill=PaleTeal,thick] (2,2) rectangle (3.4,2.8);
  \node at (2.7,2.4) {\(j\)};
  \draw[Teal,fill=PaleTeal,thick] (5,2) rectangle (6.4,2.8);
  \node at (5.7,2.4) {\(j\)};
  \draw[Blue,fill=PaleBlue,thick] (6.9,2) rectangle (8.4,2.8);
  \node at (7.65,2.4) {\(i\)};
  \draw[Blue,fill=PaleBlue,thick] (.9,-1) rectangle (2.4,-.2);
  \node at (1.65,-.6) {\(i\)};
  \draw[Teal,fill=PaleTeal,thick] (.9,.1) rectangle (2.4,.9);
  \node at (1.65,.5) {\(j\)};
  \draw[Teal,fill=PaleTeal,thick] (5.9,-1) rectangle (7.4,-.2);
  \node at (6.65,-.6) {\(j\)};
  \draw[Blue,fill=PaleBlue,thick] (5.9,.1) rectangle (7.4,.9);
  \node at (6.65,.5) {\(i\)};
\end{tikzpicture}
\caption{Bốn lớp hỗ trợ của quan hệ không chồng lấn giữa hai hình chữ nhật.
Một literal quan hệ làm chứng cho cả một lớp cặp tọa độ.}
\label{fig:four-spatial-relations}
\end{figure}

\section{Biểu diễn thứ tự và các literal so sánh}

Cho biến nguyên \(z\in\{0,\ldots,U\}\). Biểu diễn thứ tự dùng
\[
  o_{z,t}=1
  \quad\Longleftrightarrow\quad
  z\ge t,
  \qquad t=1,\ldots,U.
\]
Tính đơn điệu được biểu diễn bởi
\[
  o_{z,t+1}\rightarrow o_{z,t}.
\]
Một cận \(z\ge a\) là mệnh đề đơn vị \(o_{z,a}\); cận \(z\le b\) là
\(\neg o_{z,b+1}\). Phép so sánh dịch chuyển
\[
  x_i+w_i\le x_j
\]
trở thành các mệnh đề cấm mọi cặp ngưỡng mâu thuẫn. Khi \(w_i\) là hằng, cấu
trúc dịch chuyển này có tính địa phương cao và phù hợp với bài toán xếp khối (packing problem)
\parencite{soh2010binpacking,kieu2025strip,kieu2025bin}.

Với lập lịch, cùng kỹ thuật biểu diễn
\[
  s_i+p_i\le s_j.
\]
Nhờ vậy, biến quan hệ \(P_{ij}\) có thể liên kết trực tiếp với các literal thứ
tự của \(s_i,s_j\), thay vì tạo một mạch cộng tổng quát
\parencite{kieu2025scheduling}.

\begin{workedexample}{Liên kết biểu diễn trực tiếp--thứ tự cho một giá trị nguyên}
Cho \(z\in\{0,1,2,3,4\}\). Biểu diễn trực tiếp có \(d_a\Leftrightarrow(z=a)\);
biểu diễn thứ tự có \(o_t\Leftrightarrow(z\ge t)\). Hai góc nhìn được nối bởi
\[
  d_a\leftrightarrow o_a\land\neg o_{a+1},
  \qquad o_0\equiv\top,\quad o_5\equiv\bot.
\]
Nếu bộ giải suy ra \(o_3=1,o_4=0\), các mệnh đề liên kết buộc \(d_3=1\), rồi
ràng buộc đúng một trên \(d_0,\ldots,d_4\) loại mọi giá trị khác. Theo chiều
ngược, \(d_2=1\) buộc \(o_1=o_2=1\) và \(o_3=o_4=0\), khiến các ràng buộc cận
đọc được ngay \(z=2\). Ví dụ này cho thấy lợi ích của liên kết hai chiều không
nằm ở số mô hình mà ở đường truyền suy luận giữa hai giao diện.
\end{workedexample}

\section{Liên kết một chiều hay hai chiều}

Giả sử \(q\) là biến quan hệ và \(\Phi\) là công thức miền gốc biểu diễn quan
hệ đó. Có ba mức liên kết:
\[
  q\rightarrow\Phi,\qquad
  \Phi\rightarrow q,\qquad
  q\leftrightarrow\Phi.
\]
Chiều đầu đủ cho tính đúng nếu \(q\) chỉ được dùng làm biến làm chứng được chọn:
bộ giải không thể đặt \(q=1\) khi quan hệ hình học sai. Chiều ngược giúp miền
tọa độ suy ra \(q\); liên kết hai chiều cho phép lan truyền qua cả hai biểu
diễn.

Lựa chọn phụ thuộc vào cách mô-đun sử dụng đọc biến:

\begin{itemize}
  \item Nếu mệnh đề \(q_1\lor\cdots\lor q_r\) yêu cầu một biến làm chứng và không
  nơi nào đọc \(\neg q\), quan hệ kéo theo tiến thường đủ.
  \item Nếu \(\neg q\) được dùng để kết luận quan hệ đối nghịch, phải có
  đặc tả mạnh hơn.
  \item Nếu nhiều ràng buộc và hàm mục tiêu chia sẻ \(q\), quan hệ tương
  đương giúp
  tránh các trạng thái phụ không mang nghĩa.
\end{itemize}

\begin{example}
Với hai tác vụ cùng máy, đặt \(P_{ij}\lor P_{ji}\). Hai mệnh đề liên kết
\[
  P_{ij}\rightarrow(s_i+p_i\le s_j),
  \qquad
  P_{ji}\rightarrow(s_j+p_j\le s_i)
\]
đã bảo đảm chúng không chồng lấn. Nếu quan hệ trước--sau từ một ràng buộc khác cho
biết \(s_i+p_i\le s_j\), chiều ngược
\((s_i+p_i\le s_j)\rightarrow P_{ij}\) có thể truyền quan hệ này sang các
mô-đun khác.
\end{example}

\subsection{Lan truyền qua vòng liên kết}

Phép biểu diễn kết hợp tạo một vòng suy luận:
\[
  \text{lựa chọn chính xác}
  \rightarrow
  \text{các ngưỡng}
  \rightarrow
  \text{quan hệ}
  \rightarrow
  \text{ngưỡng của biến khác}.
\]
Mỗi chiều liên kết đóng một đường suy luận trong vòng trên. Phép đánh giá bắt
đầu bằng danh sách các suy luận miền mong muốn, sau đó tìm tập chiều tối thiểu
tạo được chúng. Đặc tả này có thể kiểm bằng các phép gán một phần và lan truyền
đơn vị như trong Chương 2.

Việc giữ nhiều góc nhìn cũng liên quan đến biểu diễn lại có cấu trúc: biến phụ có
giá trị khi đặt tên cho một lớp hỗ trợ được nhiều ràng buộc dùng lại, thay
vì chỉ thay một biểu thức xuất hiện một lần
\parencite{haberlandt2023aux}.

\begin{figure}[tb]
\centering
\begin{tikzpicture}[satfig,node distance=17mm]
  \node[satbox] (d) {lựa chọn trực tiếp\\\(d_a\)};
  \node[satbox,right=of d] (o) {literal thứ tự\\\(o_t\)};
  \node[satbox,right=of o] (q) {quan hệ\\\(q\)};
  \node[satbox,below=of o] (v) {miền của\\biến khác};
  \draw[satedge] (d)--node[above,satlabel]{giá trị \(\to\) cận}(o);
  \draw[satedge] (o)--node[above,satlabel]{hỗ trợ}(q);
  \draw[satedge] (q)--node[right,satlabel]{dịch chuyển}(v);
  \draw[satedge] (v)--node[left,satlabel]{loại cận}(d);
  \draw[satdash,bend left=20] (o) to node[below,satlabel]{chiều ngược}(d);
\end{tikzpicture}
\caption{Vòng liên kết trong phép biểu diễn kết hợp. Thiếu một chiều có thể làm
đứt một đường lan truyền đơn vị dù công thức vẫn đúng theo phép chiếu.}
\label{fig:channeling-loop}
\end{figure}

\section{Phép quay và kích thước có điều kiện}

Khi vật \(i\) được phép quay, biến \(\rho_i\) chọn hướng đặt:
\[
  (w_i(\rho_i),h_i(\rho_i))
  =
  \begin{cases}
    (w_i,h_i),&\rho_i=0,\\
    (h_i,w_i),&\rho_i=1.
  \end{cases}
\]
Các literal quan hệ được chia sẻ cho hai hướng đặt và được liên kết với phép
so sánh dịch chuyển dưới điều kiện gác \(\rho_i,\rho_j\). Ví dụ, \(L_{ij}\)
sinh bốn họ mệnh đề, mỗi họ ứng với một tổ hợp hướng đặt còn hợp lệ. Các tổ
hợp không vừa khung chứa bị loại bằng mệnh đề đơn vị hoặc mệnh đề nhị phân
trước khi sinh quan hệ cặp.

Cách tổ chức này tách ba tầng:
\[
  \text{hướng đặt}
  \longrightarrow
  \text{kích thước hiệu dụng}
  \longrightarrow
  \text{quan hệ không chồng lấn}.
\]
Cấu trúc này cũng cho phép tái sử dụng cùng biến \(\rho_i\) trong cận miền, quy tắc trội
và bộ giải mã.

\section{Đồ thị xung đột và tính địa phương}

Đồ thị xung đột \(H=(V,E)\) chọn chính xác các cặp cần biến quan hệ:
\(\{i,j\}\in E\) khi hai đối tượng có thể cạnh tranh cùng tài nguyên, chồng
lấn hình học hoặc vi phạm một cận cục bộ.

Trong lập lịch, cạnh xuất hiện khi hai tác vụ có thể cùng dùng một tài nguyên
và miền thời gian của chúng giao nhau. Trong bài toán xếp khối, mọi cặp vật cùng khung
chứa là ứng viên, nhưng cận miền và quan hệ trội có thể cố định trước một số
quan hệ. Trong phân bổ tần số, cạnh là cặp máy phát có yêu cầu khoảng cách nhãn.

Thiết kế theo đồ thị đem lại hai lợi ích: kích thước tỉ lệ với số tương tác
thật, và mệnh đề học được thường nằm quanh một lân cận nhỏ. Đây là một tiêu chí
cụ thể về mức độ phù hợp với bộ giải, thay vì chỉ nói phép biểu diễn ``gọn''.

\section{Kiến trúc kết hợp}

Một mô hình SAT lớn nên ghép các mô-đun theo thứ tự sau:

\begin{enumerate}
  \item chọn biến chính loại các chỉ số không mang nghĩa;
  \item tạo miền thứ tự/trực tiếp cho thời gian, tọa độ hoặc nhãn;
  \item tạo biến quan hệ trên đồ thị xung đột;
  \item liên kết mỗi quan hệ với các literal miền;
  \item thêm ràng buộc đếm hoặc bộ đếm chia sẻ cho tài nguyên;
  \item thêm các ràng buộc phá đối xứng có chứng minh quỹ đạo;
  \item gắn hàm mục tiêu bằng các giả định hoặc mệnh đề mềm.
\end{enumerate}

Các nghiên cứu về lập lịch, xếp khối dải/khung chứa và Social Golfer trong
\textcite{kieu2025thesis} cho thấy ba lựa chọn đầu---biến chính, quan hệ và
đối xứng---quyết định cấu trúc CNF trước khi chọn bộ giải. Các chương sau dùng
chính kiến trúc này cho bốn họ bài toán.

\section{Bài học thiết kế}

\begin{summarybox}
\begin{enumerate}
  \item Loại đối xứng biểu diễn trước khi thêm các mệnh đề phá đối xứng.
  \item Một breaker hợp lệ phải giữ ít nhất một đại diện của mọi quỹ đạo.
  \item Biến quan hệ là biến làm chứng có ngữ nghĩa, giúp cục bộ hóa phép tuyển.
  \item Liên kết mạnh đến đâu phụ thuộc vào cách biến quan hệ được đọc.
  \item Biểu diễn thứ tự làm các cận dịch chuyển trở thành mệnh đề đơn giản.
  \item Đồ thị xung đột xác định nơi thật sự cần biến phụ.
\end{enumerate}
\end{summarybox}

\subsection*{Câu hỏi nghiên cứu}
\begin{enumerate}
  \item Chứng minh tiền tố sử dụng giữ một đại diện khi các khung chứa hoàn
  toàn giống nhau.
  \item Viết liên kết một chiều cho bốn quan hệ không chồng lấn của hai hình.
  \item Cho một ràng buộc phản xạ trong bài toán xếp khối (packing problem) dải và nêu điều kiện
  để ràng buộc đó hợp lệ.
\end{enumerate}


\part{Các họ bài toán tối ưu tổ hợp}
\chapter{Lập lịch và cân bằng dây chuyền}
\chapterlead{Một phép biểu diễn cho bài toán thời gian chỉ thực sự rõ khi trả lời trực tiếp
ba câu hỏi: quyết định nằm ở mốc nào, điều gì phải liên tục giữa các mốc, và
cận mục tiêu được thay đổi ra sao trong quá trình giải.}

\index{scheduling}\index{non-preemption}\index{Nurse Rostering}
\index{lazy clause generation}\index{Dynamic Discretization Discovery}

\flowdiagram{Miền thời gian}{Biến quyết định}{Liên tục và tài nguyên}
  {Hàm mục tiêu}{Chứng nhận}

\section{Từ SAT theo mốc thời gian đến lazy clause generation}

Lập lịch là một trong những miền sớm cho thấy SAT có thể giải các bài toán
ngoài lôgic thuần túy. Các thí nghiệm SAT về lập lịch từ thập niên 1990 đã
dùng quyết định theo mốc thời gian và mệnh đề xung đột để so sánh với phương
pháp tìm kiếm chuyên dụng \parencite{crawford1994scheduling}. Từ đó, văn liệu
phát triển theo ba hướng:
\begin{enumerate}
  \item biên dịch toàn bộ ràng buộc gán, thời gian và tài nguyên sang SAT;
  \item dùng MaxSAT/PB để biểu diễn mức phạt, ưu tiên hoặc số ca vi phạm;
  \item tích hợp lan truyền miền với học mệnh đề, điển hình là
        \emph{lazy clause generation} (sinh mệnh đề lười, LCG).
\end{enumerate}

LCG giữ biến miền và bộ lan truyền của CP, đồng thời gắn mỗi phép loại giá trị
với một giải thích dạng mệnh đề; phân tích xung đột sau đó học trên các giải
thích này. Trên RCPSP/max, kiến trúc đó kết hợp lan truyền lập lịch toàn cục
với học mệnh đề \parencite{schutt2013lcg}. Ở một hướng khác, MaxSAT bộ phận có
trọng số cho phép tách độ phủ/tính hợp lệ thành các mệnh đề cứng và ưu tiên
trong lập lịch nhân viên thành các mệnh đề mềm
\parencite{demirovic2019staff}.

\begin{center}
\captionof{table}{Các mô hình giải chính xác cho bài toán thời gian.}
\label{tab:scheduling-paradigms}
\begin{tabularx}{\textwidth}{@{}>{\bfseries\sffamily}p{.18\textwidth}YYY@{}}
\toprule
Mô hình & Trạng thái thời gian & Cơ chế tài nguyên & Cận chủ đạo\\
\midrule
SAT theo mốc thời gian & Mốc rời rạc & Xung đột/ràng buộc đếm & UNSAT theo khung thời gian\\
SAT thứ tự & Ngưỡng thời gian & Quan hệ + cận dịch chuyển & Học mệnh đề\\
MaxSAT/PB & Như SAT + literal chi phí & Cứng/mềm hoặc PB & Lõi/cận chi phí\\
SMT & Biến nguyên/thực & Phép tuyển + lý thuyết & Lan truyền lý thuyết\\
CP/LCG & Miền và khoảng & Bộ lan truyền toàn cục & Lọc miền + mệnh đề\\
MIP & Biến nhị phân/liên tục & Bất đẳng thức tuyến tính & Nới lỏng quy hoạch tuyến tính\\
\bottomrule
\end{tabularx}
\end{center}

Vì vậy, một phép biểu diễn lập lịch được đánh giá trên cả kích thước, sức lan
truyền và biểu diễn đối chứng tương ứng: SAT theo mốc thời gian so với MIP mở
rộng theo thời gian; SAT thứ tự so với lôgic sai phân/SMT; bộ đếm chuỗi dùng
chung so với bộ lan truyền \textsc{Sequence}; và hàm mục tiêu MaxSAT so với PB
hay hàm mục tiêu tuyến tính.

\section{Kiến trúc chung của mô hình thời gian}

Các bài toán trong chương có bề mặt khác nhau: xếp tác vụ không ngắt
trên các máy đồng nhất, phân ca điều dưỡng, lập lịch cuộc họp B2B với cân bằng thời
gian chờ, cân bằng dây chuyền theo đỉnh công suất và điều chỉnh lịch tàu. Tuy
nhiên, phép biểu diễn của chúng đều có thể tách thành bốn lớp:
\begin{enumerate}
  \item \textbf{gán}: tác vụ, người hoặc chuyến tàu sử dụng tài nguyên nào;
  \item \textbf{thời gian/thứ tự}: thời điểm bắt đầu, kết thúc hoặc thứ tự trước--sau;
  \item \textbf{liên tục/chuỗi}: các quyết định cục bộ ghép thành một
        khoảng hoặc một chuỗi hợp lệ;
  \item \textbf{hàm mục tiêu}: thời điểm hoàn tất toàn lịch, đỉnh công suất,
        độ trễ hay số vi phạm.
\end{enumerate}

Phân tách kiến trúc theo lớp không chỉ giúp trình bày mạch lạc, mà còn cho phép thay thế một mô-đun mà không
đổi ngữ nghĩa của các mô-đun còn lại. Chẳng hạn, có thể giữ nguyên biến gán
máy và thay phép biểu diễn cho tính không ngắt; hoặc giữ nguyên mô hình khả thi
và thay vòng lặp tối ưu bằng \SAT tăng dần hay \MaxSAT.

\begin{designrule}{Nguyên tắc thiết kế}
Mỗi biến phụ phải thuộc về một lớp rõ ràng. Nếu một biến vừa được dùng để suy
thời điểm, vừa được diễn giải như trạng thái hoạt động, cần có liên kết hai
chiều hoặc một phép giải mã chứng minh được.
\end{designrule}

\section{Lập lịch không ngắt trên các tài nguyên đồng nhất}

\subsection{Mô hình}

Cho tập tác vụ \(J=\set{1,\ldots,n}\) và \(m\) tài nguyên đồng nhất. Tác vụ
\(i\) có thời điểm sẵn sàng \(r_i\), thời lượng \(e_i\) và hạn hoàn tất
\(d_i\). Thời điểm bắt đầu thuộc miền
\[
  T_i=\set{r_i,r_i+1,\ldots,d_i-e_i}.
\]
Nếu \(s_i\in T_i\), tác vụ chiếm tài nguyên trong nửa khoảng
\([s_i,s_i+e_i)\). Hai tác vụ được gán cùng tài nguyên không được có hai khoảng
giao nhau.

Một mô hình gọn dùng hai họ biến:
\[
  y_{ij}=1 \iff i\text{ được gán cho tài nguyên }j,
  \qquad
  z_{it}=1 \iff i\text{ hoạt động tại }t .
\]
Với mỗi tác vụ, các biến \(y_{ij}\) thỏa
\[
  \ExactlyOne(y_{i1},\ldots,y_{im}).
\]
Xung đột tài nguyên tại mốc \(t\) được loại bởi
\[
  \neg z_{it}\lor \neg y_{ij}\lor
  \neg z_{i't}\lor \neg y_{i'j},
  \qquad i<i'.
\]
Mệnh đề này thể hiện xung đột có điều kiện: mệnh đề chỉ kích hoạt khi hai tác vụ cùng
chọn tài nguyên \(j\).

\subsection{Từ ràng buộc đếm đến một khối liên tục}

Điều kiện
\[
  \sum_t z_{it}=e_i
\]
chỉ bảo đảm đủ số mốc hoạt động; cơ chế này không ngăn các mốc bị phân tách thành nhiều
đoạn. \emph{Non-preemption} (tính không ngắt: tác vụ phải chạy liên tục từ
lúc bắt đầu đến lúc kết thúc) đòi hỏi mạnh hơn:
\[
  z_{it}=1 \iff s_i\le t<s_i+e_i .
\]

Direct Encoding có thể tạo một biến \(b_{is}\) cho từng thời điểm bắt đầu
\(s\in T_i\). Trước hết phải đóng lựa chọn thời điểm bắt đầu bằng
\[
  \ExactlyOne\bigl(b_{is}:s\in T_i\bigr).
\]
Trong \CNF, phần ít nhất một là mệnh đề
\(\bigvee_{s\in T_i}b_{is}\); phần nhiều nhất một được sinh bằng biểu diễn từng
cặp, Sequential Counter hoặc một mô-đun AMO khác tùy kích thước \(T_i\). Sau đó
dùng
\[
  b_{is}\rightarrow z_{it}
  \quad\text{với }s\le t<s+e_i,
  \qquad
  b_{is}\rightarrow \neg z_{it}
  \quad\text{ở ngoài khoảng.}
\]
Cách này dễ hiểu nhưng lặp lại nhiều literal cho các thời điểm bắt đầu lân cận.
Phép biểu diễn gọn thay phần lặp bằng các biến tiền tố/hậu tố có ngữ nghĩa ``tất
cả bằng 0'' hoặc ``tất cả bằng 1''. Hai thời điểm bắt đầu kế tiếp chia sẻ gần
như toàn bộ cấu trúc; khác biệt của chúng chỉ nằm ở hai biên của khối.

\begin{proposition}[Đúng đắn của phép biểu diễn theo khối]
\label{prop:nonpreemptive-block}
Giả sử mô-đun khối bảo đảm: (i) tồn tại đúng một vị trí bắt đầu; (ii) từ vị
trí đó có đúng \(e_i\) giá trị liên tiếp bằng \(1\); và (iii) mọi vị trí trước
và sau khối đều bằng \(0\). Khi đó phép gán cho các biến \(z_{it}\) tương ứng
duy nhất với một khoảng xử lý không ngắt hợp lệ của tác vụ \(i\).
\end{proposition}

\begin{proof}
Từ (i), chọn vị trí bắt đầu \(s_i\). Điều kiện (ii) cho
\(z_{it}=1\) trên \([s_i,s_i+e_i)\); điều kiện (iii) loại mọi giá trị \(1\)
khác. Chiều ngược lại, từ một khoảng hợp lệ, đặt biến bắt đầu ở đầu khoảng và
gán các biến tiền tố/hậu tố theo định nghĩa của chúng. Mọi mệnh đề của mô-đun
đều được thỏa.
\end{proof}

\begin{figure}[tb]
\centering
\begin{tikzpicture}[satfig,x=9mm,y=8mm]
  \draw[->,Rule] (0,0)--(9,0) node[right,satlabel]{thời gian};
  \foreach \t in {0,...,8}{
    \draw[Rule] (\t,.12)--(\t,-.12);
    \node[below,satlabel] at (\t,-.15) {\(\t\)};
  }
  \node[left,satlabel] at (0,1.05) {hợp lệ};
  \draw[Blue,fill=PaleBlue,thick] (2,.55) rectangle (5,1.45);
  \node at (3.5,1) {tác vụ \(i\), \(e_i=3\)};
  \node[left,satlabel] at (0,2.55) {không hợp lệ};
  \draw[Gold,fill=PaleGold,thick] (1,2.1) rectangle (2,3);
  \draw[Gold,fill=PaleGold,thick] (4,2.1) rectangle (6,3);
  \node[satlabel,text=Gold] at (3.5,3.35)
    {đủ \(3\) mốc nhưng bị ngắt};
\end{tikzpicture}
\caption{Ràng buộc đếm \(\sum_tz_{it}=3\) không phân biệt một khối liên tục với
ba mốc rời; tính không ngắt phải biểu diễn vị trí bắt đầu và hai biên khối.}
\label{fig:nonpreemptive-block}
\end{figure}

\begin{workedexample}{Ba tác vụ trên hai máy}
Cho \(J_1=(r=0,e=3,d=6)\), \(J_2=(0,2,5)\), \(J_3=(1,2,6)\).
Một lịch hợp lệ là
\[
  M_1:\ J_1[0,3)\ \ J_3[3,5),
  \qquad
  M_2:\ J_2[0,2).
\]
Với \(t=1\), \(z_{1,1}=z_{2,1}=1\) nhưng hai tác vụ dùng máy khác nhau nên
mệnh đề xung đột có điều kiện vẫn thỏa. Nếu đổi \(y_{2,2}\) thành \(y_{2,1}\),
mệnh đề tại \(t=1,j=1\) trở thành
\(
\neg z_{1,1}\lor\neg y_{1,1}\lor
\neg z_{2,1}\lor\neg y_{2,1}
\)
và bị vi phạm. Ví dụ tách rõ hai lớp: \(z\) nói khi nào tác vụ chạy, \(y\)
nói tác vụ chạy ở vị trí nào.
\end{workedexample}

Kết quả của \textcite{kieu2025scheduling} cho thấy lợi ích chính của phép mã
hóa gọn nằm ở việc tránh lặp cấu trúc liên tục. Cấu hình kết hợp biểu diễn theo
khối với phá vỡ đối xứng và một phép biểu diễn giả Boolean thích hợp đạt tổng thời
gian thấp nhất trong nhiều nhóm bài toán của nghiên cứu đó. Điều cần rút ra cho
thiết kế là: \emph{tính không ngắt nên là một mô-đun độc lập}, thay vì được
ngầm suy từ một ràng buộc đếm.

\section{Ràng buộc chuỗi trong lập lịch điều dưỡng}

\subsection{Cửa sổ trượt}

Với dãy Boolean \(\Omega=(x_1,\ldots,x_n)\), một ràng buộc chuỗi cận dưới yêu
cầu mọi cửa sổ độ rộng \(w\) chứa ít nhất \(k\) giá trị đúng:
\[
  \operatorname{ALSC}(\Omega,w,k)=
  \bigwedge_{i=0}^{n-w}
  \left(\sum_{j=i+1}^{i+w}x_j\ge k\right).
\]
Nếu biểu diễn từng cửa sổ bằng một phép biểu diễn ràng buộc đếm độc lập, hai cửa sổ
kề nhau chia sẻ \(w-1\) biến gốc nhưng không chia sẻ biến phụ. Cấu trúc chuỗi
vì thế đặt ra bài toán \emph{tái sử dụng biểu thức con}.

Chia một cửa sổ thành hai phần \(A\) và \(B\). Ta có
\[
  \sum_{x\in A\cup B}x\ge k
  \iff
  \bigwedge_{q=1}^{k}
  \left[
    \left(\sum_{x\in A}x\ge q\right)
    \lor
    \left(\sum_{x\in B}x\ge k-q+1\right)
  \right].
\]
Mỗi mệnh đề ngưỡng trong ngoặc được biểu diễn bởi một thanh ghi chia sẻ.
Khi cửa sổ trượt sang phải, phần lớn thanh ghi vẫn còn nguyên.

\begin{example}
Xét \(x_3+x_4+x_5+x_6\ge2\), chia thành hai khối
\(A=(x_3,x_4)\) và \(B=(x_5,x_6)\). Nếu
\(R_{A,q}\) mang nghĩa ``\(A\) có ít nhất \(q\) giá trị đúng'', bộ nối chỉ
cần
\[
  (R_{A,1}\lor R_{B,2})
  \land
  (R_{A,2}\lor R_{B,1}).
\]
Các thanh ghi của \(B\) tiếp tục được dùng trong cửa sổ kế tiếp.
\end{example}

\begin{figure}[tb]
\centering
\begin{tikzpicture}[satfig,x=12mm,y=7mm]
  \foreach \d/\v in {1/1,2/0,3/1,4/0,5/1,6/0,7/1}{
    \ifnum\v=1
      \node[satbox,minimum width=10mm,fill=PaleTeal] (d\d) at (\d,0) {\(\v\)};
    \else
      \node[satbox,minimum width=10mm,fill=SoftGray] (d\d) at (\d,0) {\(\v\)};
    \fi
    \node[above,satlabel] at (d\d.north) {ngày \(\d\)};
  }
  \draw[Blue,very thick] (.55,-.75)--(4.45,-.75)
    node[midway,below,satlabel,text=Blue]{\(W_1\): tổng \(2\)};
  \draw[Teal,very thick] (1.55,-1.55)--(5.45,-1.55)
    node[midway,below,satlabel,text=Teal]{\(W_2\): tổng \(2\)};
  \draw[Gold,very thick] (2.55,-2.35)--(6.45,-2.35)
    node[midway,below,satlabel,text=Gold]{\(W_3\): tổng \(2\)};
\end{tikzpicture}
\caption{Ba cửa sổ ALSC \((w=4,k=2)\) trên bảy ngày. Mỗi cửa sổ kề nhau chỉ
thay một ngày, nên các ngưỡng ở phần giao có thể dùng lại.}
\label{fig:nurse-alsc}
\end{figure}

\subsection{Đưa vào mô hình phân ca}

Đặt \(x_{ids}=1\) khi điều dưỡng \(i\) làm loại ca \(s\) ở ngày \(d\). Mỗi cặp
\((i,d)\) nhận đúng một trạng thái. Những quy tắc như ``trong mọi \(w\) ngày có
ít nhất \(k\) ngày nghỉ'' trở thành ALSC; quy tắc at-most được đối ngẫu:
\[
  \sum_{j=1}^{w}x_j\le q
  \iff
  \sum_{j=1}^{w}(1-x_j)\ge w-q.
\]
Nhờ đó một bản cài đặt có thể dùng cùng một mô-đun cận dưới cho cả hai hướng.

Trong một mô hình phân ca đầy đủ, ràng buộc chuỗi chỉ là một lớp. Lớp độ phủ
bảo đảm số nhân viên mỗi ca; ràng buộc hợp đồng giới hạn tổng giờ và số cuối
tuần; ưu tiên tạo các chi phí mềm. Phép biểu diễn cửa sổ chia sẻ tác động trực
tiếp lên lớp chuỗi, còn MaxSAT có trọng số là một lựa chọn cho hàm mục tiêu của
toàn mô hình \parencite{demirovic2019staff}. Ma trận thiết kế tách các lớp để
đo riêng ảnh hưởng của ALSC, độ phủ và ưu tiên.

\begin{resultbox}{Kết quả tiêu biểu}
Trên 28 bài toán \emph{Nurse Rostering} (lập lịch điều dưỡng), với quy mô đến
120 điều dưỡng và 24 tuần,
ALSCE giải toàn bộ tập thử trong nghiên cứu của
\textcite{truong2025nurse}; bài toán lớn nhất được báo cáo ở mức
\(26{,}01\) giây. Kết quả này minh họa hiệu quả của kỹ thuật chia sẻ thanh ghi giữa
các cửa sổ chồng lấn.
\end{resultbox}

\section{Lập lịch cuộc họp B2B với cân bằng thời gian chờ}

\index{B2B meeting scheduling}\index{idle-time balancing}
\index{SparseSuffix}\index{domain filtering}

\subsection{Bài toán B2B và mục tiêu khoảng thời gian chờ}

Bài toán lập lịch cuộc họp doanh nghiệp (Business-to-Business Meeting Scheduling Problem, B2B) xếp lịch cho tập các cuộc họp đối tác \(M\) vào tập mốc thời gian rời rạc \(T=\set{1,\ldots,h}\) và \(K\) phòng họp/địa điểm dùng chung \parencite{pesant2015b2b,bofill2022b2b}. Mỗi cuộc họp \(m\in M\) được thực hiện bởi đúng hai đại biểu, ký hiệu \(\pi(m)\subseteq P\) với \(\card{\pi(m)}=2\), và có miền mốc khởi tạo \(D_m^0\subseteq T\) tương ứng với phiên làm việc, mốc cố định và mốc cấm của hai đại biểu. Đặt \(E\subseteq M\times M\) là tập các cung thứ tự tiên quyết giữa các cuộc họp. Một lịch xếp \(S: M\to T\) là hợp lệ nếu \(S(m)\in D_m^0\) và thỏa mãn các điều kiện:
\begin{align}
  S(m) &\ne S(m') && \forall m,m'\in M, m\ne m': \pi(m)\cap\pi(m')\ne\emptyset, \label{eq:b2b-participant}\\
  \card{\set{m\in M \suchthat S(m)=t}} &\le K && \forall t\in T, \label{eq:b2b-capacity}\\
  S(i) &< S(j) && \forall (i,j)\in E. \label{eq:b2b-precedence}
\end{align}

Với mỗi đại biểu \(p\in P\), đặt \(M_p = \set{m\in M \suchthat p\in\pi(m)}\) là tập cuộc họp của đại biểu đó. Đặt \(P^\star = \set{p\in P \suchthat \card{M_p}\ge 2}\) là tập các đại biểu có ít nhất hai cuộc họp. Đối với \(p\in P^\star\), mốc bận đầu tiên \(F_p(S)\) và mốc bận cuối cùng \(L_p(S)\) trong lịch \(S\) được xác định bởi:
\[
  F_p(S) = \min_{m\in M_p} S(m), \qquad L_p(S) = \max_{m\in M_p} S(m).
\]
Do các cuộc họp của đại biểu \(p\) không chồng lấn, số mốc thời gian trống nằm hoàn toàn giữa mốc đầu và mốc cuối---gọi là \emph{thời gian chờ nội bộ} (internal idle slots)---được tính bằng:
\begin{equation}
  I_p(S) =
  \begin{cases}
    L_p(S) - F_p(S) + 1 - \card{M_p}, & p\in P^\star,\\
    0, & \card{M_p}\le 1.
  \end{cases}
  \label{eq:b2b-internal-idle}
\end{equation}
Khác với mô hình truyền thống đếm tổng số nhóm nghỉ (\emph{break groups}) \parencite{bofill2022b2b}, chỉ số \(I_p(S)\) phân biệt rõ khoảng chờ dài (như chờ 5 mốc) với khoảng chờ ngắn (chờ 1 mốc). Hàm mục tiêu tối ưu hóa khoảng thời gian chờ nội bộ giữa các đại biểu:
\begin{equation}
  \Delta_I(S) = \max_{p\in P^\star} I_p(S) - \min_{p\in P^\star} I_p(S).
  \label{eq:b2b-idle-range}
\end{equation}
Mục tiêu này hướng tới sự cân bằng công bằng về thời gian chờ giữa các đại biểu, lấy cảm hứng từ các chuẩn cân bằng tải trong bài toán BACP \parencite{chiarandini2012}. Tổng thời gian chờ \(I_\Sigma(S)=\sum_{p\in P}I_p(S)\) được báo cáo như một chỉ số bổ trợ.

\subsection{Tiền xử lý bảo toàn nghiệm và lọc miền}

Trước khi sinh công thức \CNF, thuật toán lọc miền \parencite{nguyen2026compactsatmaxsatencodings} thu hẹp miền mốc của từng cuộc họp bằng cách áp dụng lặp cho đến điểm bất động ba quy tắc lọc:
\begin{enumerate}
  \item \textbf{Cận thứ tự tiên quyết}: Với mỗi quan hệ \((i,j,d)\) trong tập quan hệ \(\mathcal R_f\in\set{\mathcal R_E, \mathcal R_{E^\star}}\) (trong đó \(\mathcal R_E = \set{(i,j,1) \suchthat (i,j)\in E}\) và \(\mathcal R_{E^\star}\) là bao đóng bắc cầu gán khoảng cách đường đi dài nhất \(\ell_{ij}\)), loại \(t\in D_i\) nếu \(t+d > \max D_j\), và loại \(u\in D_j\) nếu \(\min D_i+d > u\).
  \item \textbf{Nhất quán ghép cặp đại biểu}: Với mỗi đại biểu \(p\), dựng đồ thị hai phía giữa \(M_p\) và miền mốc khả thi. Sử dụng thuật toán lọc \textsc{AllDifferent} dựa trên ghép cặp cực đại \parencite{regin1994alldiff} để xóa phân gán cuộc họp--mốc nào không thuộc bất kỳ ghép cặp hoàn hảo nào của \(M_p\).
  \item \textbf{Mốc bão hòa}: Nếu mốc \(t\) đã bị ép chọn bởi đúng \(K\) cuộc họp cố định, loại mốc \(t\) khỏi miền của tất cả cuộc họp còn lại.
\end{enumerate}

\begin{proposition}[Tính bảo toàn nghiệm của lọc miền B2B]
\label{prop:b2b-domain-preservation}
Thuật toán lọc miền dừng sau số bước hữu hạn, loại bỏ duy nhất các giá trị không thể thuộc bất kỳ lịch hợp lệ nào, bảo toàn trọn vẹn tập nghiệm khả thi, và chỉ thông báo vô nghiệm khi bài toán thực sự không có nghiệm.
\end{proposition}

\begin{proof}
Mỗi bước loại giá trị đều dựa trên các điều kiện cần về tính khả thi: vi phạm cận thứ tự tiên quyết, không có hỗ trợ trong ghép cặp đại biểu, hoặc vượt quá sức chứa địa điểm \(K\). Vì miền ban đầu là hữu hạn, quá trình xóa giá trị phải dừng sau tối đa \(\sum_{m}\card{D_m^0}\) bước.
\end{proof}

Phép biểu diễn \mbox{\textsc{Reduced}} áp dụng biến gán \(x_{m,t}\) chỉ trên miền mốc đã lọc \(A_m = D_m\), thay vì miền đầy đủ \(A_m = T\) của mô hình gốc \mbox{\textsc{ORG-MaxSAT}} \parencite{bofill2022b2b}.

\subsection{Biểu diễn thứ tự tiên quyết bằng hậu tố thưa chia sẻ (\textsc{SparseSuffix})}

Xét quan hệ tiên quyết \((i,j,d)\in\mathcal R_g\). Biểu diễn từng cặp (\mbox{\textnormal{\textsc{Pairwise}}}) thêm các mệnh đề nhị phân \(\neg x_{i,t}\lor\neg x_{j,u}\) với mọi \(t\in A_i, u\in A_j\) thỏa \(t+d>u\), có thể phát sinh \(O(\card{A_i}\card{A_j})\) mệnh đề cho một quan hệ.

Phép biểu diễn hậu tố thưa (\mbox{\textnormal{\textsc{SparseSuffix}}}) \parencite{nguyen2026compactsatmaxsatencodings} chỉ khởi tạo biến phụ tại các điểm cắt nội bộ thực sự xuất hiện trong miền tiền nhiệm. Với miền đã sắp thứ tự \(A_i = (a_{i,0},\ldots,a_{i,n_i-1})\), giả sử các điểm cắt xuất hiện là \(c_1 < c_2 < \cdots < c_s\). Đặt \(c_{s+1}=n_i\) và \(q_{i,n_i}=\bot\), biến đại diện hậu tố được định nghĩa quy nạp:
\begin{align}
  c(i,j,u) &= \min\set{r\in\set{0,\ldots,n_i-1} \suchthat a_{i,r} > u - d}, \notag\\
  q_{i,c_\ell} &\leftrightarrow \left(\bigvee_{r=c_\ell}^{c_{\ell+1}-1} x_{i,a_{i,r}}\right) \lor q_{i,c_{\ell+1}}, \label{eq:b2b-sparse-suffix}\\
  x_{j,u} &\rightarrow \neg q_{i,c(i,j,u)}. \label{eq:b2b-suffix-link}
\end{align}
Khi hai quan hệ tiên quyết cùng dẫn đến một điểm cắt \(c\) trên miền \(A_i\), biến hậu tố \(q_{i,c}\) được tái sử dụng hoàn toàn.

\begin{proposition}[Tương đương của SparseSuffix]
\label{prop:b2b-suffix-equivalence}
Phép biểu diễn \mbox{\textnormal{\textsc{SparseSuffix}}} và phép biểu diễn từng cặp \mbox{\textnormal{\textsc{Pairwise}}} mô tả chính xác cùng một tập phân gán cuộc họp--mốc thời gian.
\end{proposition}

\begin{proof}
Từ Công thức~\eqref{eq:b2b-sparse-suffix}, biến phụ \(q_{i,c}\) đúng khi và chỉ khi cuộc họp \(i\) được gán cho mốc có chỉ số ít nhất là \(c\). Mệnh đề kéo theo~\eqref{eq:b2b-suffix-link} do đó loại bỏ chính xác các cặp mốc \((t,u)\) vi phạm điều kiện \(t+d \le u\).
\end{proof}

\begin{figure}[tb]
\centering
\begin{tikzpicture}[satfig,x=8.5mm,y=7.5mm]
  %--- Diagram (a): Predecessor-domain cut ---
  \node[anchor=west,font=\rmfamily\bfseries\small] at (0,3.5) {(a) Điểm cắt miền tiền nhiệm};
  \node[anchor=east,font=\rmfamily\small] at (0.6,2.1) {$A_i$};
  
  \foreach \x/\v in {0.8/1, 1.8/3, 2.8/4, 3.8/7} {
    \draw[satbox] (\x,1.75) rectangle +(0.8,0.7);
    \node[font=\rmfamily\small] at (\x+0.4,2.1) {$\v$};
  }
  
  \draw[Gold,dashed,very thick] (2.7,1.5)--(2.7,2.7);
  \node[Gold,anchor=south,font=\rmfamily\bfseries\small] at (2.7,2.7) {$c=2$};
  
  \node[anchor=west,font=\rmfamily\small] at (4.8,2.1) {$x_{j,4}$};
  \node[anchor=west,font=\rmfamily\small] at (0.8,1.0) {$\neg x_{j,4} \lor \neg q_{i,2}$};
  \node[anchor=west,font=\rmfamily\small] at (0.8,0.3) {$q_{i,2} \leftrightarrow x_{i,4} \lor x_{i,7}$};

  %--- Separator ---
  \draw[Rule,thick] (6.5,0.1)--(6.5,3.6);

  %--- Diagram (b): Schedule span & idle-time threshold ---
  \node[anchor=west,font=\rmfamily\bfseries\small] at (7.2,3.5) {(b) Đánh giá ngưỡng chờ nội bộ};
  \node[anchor=east,font=\rmfamily\small] at (7.8,2.1) {mốc};
  
  \draw[satbox,fill=PaleBlue,draw=Blue,thick] (8.0,1.75) rectangle +(0.8,0.7);
  \node[font=\rmfamily\small] at (8.4,2.1) {$3$};
  
  \draw[satbox,fill=SoftGray] (9.0,1.75) rectangle +(0.8,0.7);
  \node[font=\rmfamily\small] at (9.4,2.1) {$4$};
  
  \draw[satbox,fill=PaleBlue,draw=Blue,thick] (10.0,1.75) rectangle +(0.8,0.7);
  \node[font=\rmfamily\small] at (10.4,2.1) {$5$};
  
  \draw[satbox,fill=SoftGray] (11.0,1.75) rectangle +(0.8,0.7);
  \node[font=\rmfamily\small] at (11.4,2.1) {$6$};
  
  \node[anchor=south,font=\rmfamily\small,text=Blue] at (8.4,2.55) {$f_{p,3}=1$};
  \node[anchor=south,font=\rmfamily\small,text=Blue] at (10.4,2.55) {$R_{p,5}=1$};
  
  \draw[satdim] (8.4,1.4)--(10.4,1.4);
  \node[anchor=north,font=\rmfamily\small] at (9.4,1.35) {$\card{M_p}=2 \implies I_p \ge 1$};
\end{tikzpicture}
\caption{Kỹ thuật biểu diễn trong B2B: (a) Biến hậu tố thưa \mbox{\textnormal{\textsc{SparseSuffix}}} tại điểm cắt $c=2$; (b) Xác định ngưỡng thời gian chờ nội bộ từ mốc bắt đầu và mốc kết thúc.}
\label{fig:b2b-encoding-ideas}
\end{figure}

\subsection{Mã hóa ngưỡng thời gian chờ từ nhịp hoạt động}

Để xác định \(I_p(S)\), gọi \(y_{p,t}\leftrightarrow\bigvee_{m\in M_p: t\in A_m}x_{m,t}\) là biến chỉ báo đại biểu \(p\) bận tại mốc \(t\). Định nghĩa biến tích lũy tiền tố \(P_{p,t}\leftrightarrow P_{p,t-1}\lor y_{p,t}\) và hậu tố \(R_{p,t}\leftrightarrow y_{p,t}\lor R_{p,t+1}\). Mốc bắt đầu đầu tiên của đại biểu \(p\) được nhận diện bằng:
\[
  f_{p,t} \leftrightarrow y_{p,t} \land \neg P_{p,t-1}.
\]
Gọi \(\theta_{p,k}\) là biến mệnh đề đại diện cho điều kiện \(I_p(S)\ge k\). Mối liên hệ giữa mốc đầu và mốc cuối được biểu diễn chính xác bằng:
\begin{equation}
  f_{p,t} \rightarrow \left(\theta_{p,k} \leftrightarrow R_{p, t + \card{M_p} + k - 1}\right)
  \qquad \forall t\in T, 1\le k\le U_p,
  \label{eq:b2b-span-threshold}
\end{equation}
trong đó \(U_p\) là cận trên hợp lệ của thời gian chờ đối với đại biểu \(p\).

Cuối cùng, với \(U=\max_{p\in P^\star} U_p\) và \(1\le k\le U\), biểu diễn định nghĩa biến ngưỡng cực đại \(\Theta_k^{\max}\leftrightarrow\bigvee_{p\in P^\star}\theta_{p,k}\), cực tiểu \(\Theta_k^{\min}\leftrightarrow\bigwedge_{p\in P^\star}\theta_{p,k}\), và biến mức chênh lệch:
\begin{equation}
  g_k \leftrightarrow \Theta_k^{\max} \land \neg \Theta_k^{\min}, \qquad \text{với } \sum_{k=1}^U g_k = \Delta_I(S).
  \label{eq:b2b-range-sum}
\end{equation}

\subsection{Phương pháp giải và kết quả thực nghiệm}

Phép biểu diễn B2B thu gọn được đánh giá thực nghiệm rộng rãi trên 126 bộ dữ liệu chuẩn chính thức (\emph{All-126}) và 100 bộ dữ liệu biến đổi mật độ thứ tự (\emph{Stress-100}) \parencite{nguyen2026compactsatmaxsatencodings}. Thử nghiệm so sánh ba phương pháp tối ưu dựa trên công thức SAT/MaxSAT (UWrMaxSAT, \textsc{Non-Incremental SAT}, và \textsc{Incremental SAT} sử dụng Totalizer kết hợp giả định assumption) với các công cụ thương mại hàng đầu đại diện cho MIP (Gurobi, CPLEX) và CP (CP Optimizer).

\begin{resultbox}{Kết quả tiêu biểu cho B2B Meeting Scheduling}
Trên tập dữ liệu chuẩn \emph{All-126}:
\begin{itemize}
  \item Phép biểu diễn thu gọn giảm trung vị \textbf{40.3\%} số mệnh đề, \textbf{55.9\%} dung lượng bộ nhớ đỉnh và \textbf{14.0\%} tổng thời gian chạy so với mô hình \mbox{\textsc{ORG-MaxSAT}} công bố \parencite{bofill2022b2b}.
  \item Riêng bước lọc miền (\mbox{\textsc{Reduced}}) đóng góp giảm \textbf{24.1\%} số biến gán và \textbf{16.2\%} số mệnh đề.
  \item Biểu diễn hậu tố \mbox{\textnormal{\textsc{SparseSuffix}}} kết hợp với bao đóng bắc cầu \mbox{\textnormal{\textsc{Closure-}}\ensuremath{E^\star}} giúp giảm số mệnh đề đến \textbf{5.5\%} ở các mức mật độ thứ tự cao.
  \item Cả ba phương pháp SAT và MaxSAT đều giải thành công 100\% các bài toán chính thức với thời gian trung vị vượt trội Gurobi MIP (\(0.521\) giây); trong đó \mbox{\textnormal{\textsc{Incremental SAT}}} đạt thời gian trung vị nhanh nhất (\textbf{0.227} giây).
\end{itemize}
\end{resultbox}

\section{Peak power: công suất đỉnh trong cân bằng dây chuyền}

\subsection{Mô hình đỉnh công suất}

Trong SALBP-3PM, tác vụ \(i\) có thời lượng \(t_i\), công suất \(w_i\), được gán
vào một trạm và bắt đầu trong chu kỳ dài \(c\). Nếu \(A_{i\tau}\) biểu diễn
``tác vụ \(i\) đang hoạt động tại \(\tau\)'', tổng công suất là
\[
  W(\tau)=\sum_i w_iA_{i\tau},
  \qquad
  P=\max_{\tau=0}^{c-1}W(\tau).
\]
Mục tiêu là tối thiểu hóa \(P\).

Các biến trực tiếp cho quyết định gán và bắt đầu thuận tiện cho phép giải mã:
\[
  X_{ik}=1 \iff i\text{ ở trạm }k,
  \qquad
  S_{it}=1 \iff i\text{ bắt đầu tại }t.
\]
Các biến thứ tự cung cấp lớp lan truyền:
\[
  R_{ik}=1 \iff \text{trạm}(i)\le k,
  \qquad
  T_{it}=1 \iff \text{bắt đầu}(i)\le t.
\]
Giá trị chính xác xuất hiện tại điểm chuyển của chuỗi thứ tự. Nhờ vậy, ràng
buộc AMO về phép gán giảm từ dạng từng cặp bậc hai xuống một chuỗi tuyến tính;
quan hệ trước--sau giữa trạm hoặc giữa mốc thời gian cũng được viết bằng quan
hệ kéo theo giữa các ngưỡng.

\subsection{Bốn giao diện tối ưu}

Cùng một phép biểu diễn tính khả thi có thể kết nối với hàm mục tiêu theo bốn cách:
\begin{center}
\captionof{table}{Bốn giao diện tối ưu cho công suất đỉnh.}
\label{tab:peak-interfaces}
\begin{tabularx}{\textwidth}{@{}>{\bfseries\sffamily}p{.23\textwidth}YY@{}}
\toprule
Cơ chế & Cách thay cận & Thành phần cần giữ cố định\\
\midrule
Chặn bằng mệnh đề & Cấm mọi đỉnh không tốt hơn nghiệm đương nhiệm &
Ngữ nghĩa của biến đỉnh\\
Tinh chỉnh PB & Siết \(\sum_iw_iA_{i\tau}\le B\) tại mọi \(\tau\) &
Phép biểu diễn giả Boolean\\
\MaxSAT có trọng số & Đưa mức đỉnh vào mệnh đề mềm &
Ánh xạ chi phí với giá trị \(P\)\\
Biến chỉ báo thứ tự & Kích hoạt cận bằng các giả định &
Mệnh đề cơ sở và chuỗi thứ tự\\
\bottomrule
\end{tabularx}
\end{center}

Biến chỉ báo thứ tự phù hợp với giải tăng dần: bộ giải được dựng một lần ở cận
trên, sau mỗi nghiệm chỉ thay các giả định hoặc thêm mệnh đề chặn. Trong nghiên
cứu của \textcite{kieu2025power}, phép biểu diễn tuần tự gọn kết hợp giải tăng dần
chứng nhận được nhiều bài toán hơn cấu hình đối chiếu trên tập 89 bài kiểm
chuẩn; lợi ích rõ nhất ở nhóm khó.

\begin{keyidea}{Tách tính khả thi khỏi tối ưu}
Mô-đun khả thi mô tả lịch hợp lệ qua một giao diện độc lập với tìm kiếm tuyến
tính, tìm kiếm nhị phân hay \MaxSAT. Lớp hàm mục tiêu đọc các biến đã có và
cung cấp giao diện cận \(B\). Cách tách này làm cho chứng minh tính đúng ngắn
hơn và thí nghiệm phân tích thành phần (ablation study) rõ hơn.
\end{keyidea}

\section{Điều chỉnh lịch tàu bằng MaxSAT--DDD}

\subsection{Tại sao cần rời rạc hóa động}

Với tàu \(i\) và tài nguyên \(r\) trên tuyến cố định, đặt \(t^{i,r}\) là thời
điểm tàu đi vào tài nguyên. Quan hệ trước--sau trên đường đi có dạng
\[
  t^{i,q}\ge t^{i,r}+\ell_i^r,
\]
trong khi hai tàu chia sẻ tài nguyên phải chọn một trong hai thứ tự an toàn.
Nếu dựng toàn bộ miền thời gian đến khung thời gian \(H\), số biến theo mốc thời gian có thể
rất lớn.

\emph{Dynamic Discretization Discovery} (khám phá rời rạc hóa động, DDD) khởi
đầu bằng tập mốc thưa \(\Lambda^{i,r}\). Mô hình hạn chế cho một cận dưới.
Nghiệm được giải mã trên trục thời gian đầy đủ; nếu nghiệm vi phạm thời gian di
chuyển hoặc khoảng cách an toàn, chỉ các mốc liên quan mới được thêm.

DDD không bắt nguồn từ SAT. Thuật toán được phát triển cho thiết kế mạng dịch
vụ thời gian liên tục bằng cách giải lặp các mạng mở rộng thời gian một phần
và tinh chỉnh những mốc cần thiết \parencite{boland2017ddd}. Khi chuyển ý tưởng
sang MaxSAT, điều phải bảo toàn là \emph{tính hợp lệ của cận dưới} trong mô
hình hạn chế và quy tắc tinh chỉnh hữu hạn; bộ giải bên trong có thể thay đổi.
Đây là một ví dụ tiêu biểu cho quá trình tích hợp kiến trúc từ Nghiên cứu vận trù
vào SAT, thay vì chỉ thay phép biểu diễn của một ràng buộc.

\subsection{Vòng lặp}

\begin{enumerate}
  \item Giải mô hình \MaxSAT hạn chế trên phép rời rạc hóa hiện tại.
  \item Giải mã lịch và kiểm tra mọi quan hệ trước--sau trên đường đi, xung
        đột tài nguyên và giá trị độ trễ.
  \item Nếu có vi phạm, thêm các mốc cần thiết và lặp lại.
  \item Nếu lịch khả thi và chi phí giải mã bằng giá trị tối ưu của mô hình
        hạn chế, cận dưới gặp cận trên; nghiệm là tối ưu.
\end{enumerate}

\begin{algorithm}[H]
\caption{Vòng lặp MaxSAT--DDD có kiểm chứng}
\label{alg:maxsat-ddd}
\begin{minipage}{.94\textwidth}\small
\textbf{Đầu vào:} tập mốc hữu hạn đầy đủ \(\Lambda^\star\), phép rời rạc hóa ban đầu
\(\Lambda^0\subseteq\Lambda^\star\), bộ giải mã–kiểm tra và hàm chi phí.\\
\textbf{Đầu ra:} lịch tối ưu hoặc chứng nhận vô nghiệm trong miền đầy đủ.
\begin{enumerate}
  \item Với vòng \(r\), sinh hoặc cập nhật mô hình hạn chế trên
  \(\Lambda^r\), giải tối ưu và ghi cận dưới \(L_r\).
  \item Nếu mô hình hạn chế vô nghiệm và là một mô hình nới lỏng hợp lệ của mô hình
  đầy đủ, kết luận mô hình đầy đủ vô nghiệm.
  \item Giải mã nghiệm. Bộ kiểm tra xác nhận một lịch đầy đủ có chi phí \(U_r\) nếu
  khả thi; đồng thời, khi lịch chưa khả thi hoặc \(L_r<U_r\), thuật toán phải trả về một
  bằng chứng \(C_r\) chỉ ra phần rời rạc hóa đang thiếu cho feasibility hoặc
  hàm mục tiêu.
  \item Nếu có lịch khả thi và \(U_r=L_r\), trả lịch cùng hai cận bằng nhau.
  \item Nếu chưa hội tụ, bước tinh chỉnh phải chọn ít nhất một mốc trong
  \(\Lambda^\star\setminus\Lambda^r\) được bằng chứng \(C_r\) yêu cầu; đặt
  \(\Lambda^{r+1}\supsetneq\Lambda^r\) và lặp.
\end{enumerate}
\end{minipage}
\end{algorithm}

\begin{proposition}[Hội tụ hữu hạn của DDD]
\label{prop:ddd-finite}
Giả sử (i) tập mốc đầy đủ \(\Lambda^\star\) hữu hạn; (ii) mỗi mô hình hạn chế
cho một cận dưới hợp lệ; (iii) bộ kiểm tra đúng; và (iv) mỗi vòng chưa kết thúc
thêm ít nhất một mốc mới từ \(\Lambda^\star\). Khi đó
Thuật toán~\ref{alg:maxsat-ddd} kết thúc sau không quá
\(\card{\Lambda^\star\setminus\Lambda^0}+1\) lần giải mô hình hạn chế và chỉ
trả nghiệm tối ưu hoặc kết luận vô nghiệm đúng.
\end{proposition}

\begin{proof}
Do chuỗi
\(\Lambda^0\subsetneq\Lambda^1\subsetneq\cdots\subseteq\Lambda^\star\)
tăng nghiêm ngặt trong một tập hữu hạn, số lần tinh chỉnh bị chặn bởi
\(\card{\Lambda^\star\setminus\Lambda^0}\). Nếu thuật toán dừng với
\(U_r=L_r\), \(L_r\) là cận dưới cho giá trị tối ưu đầy đủ còn \(U_r\) là chi
phí của một nghiệm đầy đủ được bộ kiểm tra xác nhận, nên cả hai bằng giá trị
tối ưu. Nếu đến \(\Lambda^\star\), mô hình hạn chế chính là mô hình đầy đủ; bộ giải quyết định
tối ưu hoặc vô nghiệm. Kết luận vô nghiệm sớm chỉ hợp lệ dưới giả thiết
nới lỏng ở (ii).
\end{proof}

\paragraph{Biên hội tụ và chiến lược tinh chỉnh.}
Ở biên xấu nhất, DDD khôi phục toàn bộ lưới \(\Lambda^\star\). Mỗi bước tinh
chỉnh vì thế phải thêm một mốc mới được bộ kiểm tra suy ra từ vi phạm cụ thể.
Với miền thời gian liên tục, một đối chứng thời gian liên tục hoặc một định lý
về tập ứng viên hữu hạn cung cấp điều kiện dừng. Ba thành phần này---mốc mới, bộ kiểm tra và
tập mốc đầy đủ---xác định đồng thời tính đúng và tốc độ hội tụ.

\begin{figure}[tb]
\centering
\begin{tikzpicture}[satfig,x=10mm,y=8mm]
  \draw[->,Rule] (0,0)--(10,0) node[right,satlabel]{thời gian};
  \foreach \t in {0,2,5,8,10}{
    \draw[Rule] (\t,.12)--(\t,-.12);
    \node[below,satlabel] at (\t,-.15) {\(\t\)};
  }
  \node[left,satlabel] at (0,1.1) {vòng \(r\)};
  \foreach \t in {0,5,10}
    \node[satstate,minimum size=6mm] at (\t,1.1) {};
  \draw[Blue,thick] (0,1.1)--(5,1.1)--(10,1.1);
  \draw[Gold,very thick] (5,1.1)--(8,1.1)
    node[midway,above,satlabel,text=Gold]{vi phạm được giải mã};
  \node[left,satlabel] at (0,2.5) {vòng \(r+1\)};
  \foreach \t in {0,5,10}
    \node[satstate,minimum size=6mm] at (\t,2.5) {};
  \node[satstate,satwarn,minimum size=6mm] at (8,2.5) {};
  \draw[Teal,thick] (0,2.5)--(5,2.5)--(8,2.5)--(10,2.5);
  \draw[satedge] (8,1.55)--(8,2.1)
    node[midway,right,satlabel]{thêm mốc \(8\)};
\end{tikzpicture}
\caption{DDD chỉ thêm mốc tại nơi nghiệm của mô hình hạn chế thiếu thông tin.
Lưới không được làm dày đồng loạt trên toàn bộ khung thời gian.}
\label{fig:ddd-refinement}
\end{figure}

Lan truyền quan hệ trước--sau có vai trò thu hẹp miền trước khi sinh \CNF. Nếu một
lựa chọn thứ tự \(i\prec j\) kéo theo một chuỗi thời điểm tối thiểu, các mốc
không đạt được bị loại ngay. AMO kết hợp tiếp tục chọn phép biểu diễn khác nhau
theo kích thước miền: từng cặp cho tập rất nhỏ, tuần tự hoặc chỉ huy cho tập lớn hơn.

\begin{resultbox}{Điểm rút ra từ Train MaxSAT--DDD}
Trong nghiên cứu của \textcite{kieu2026train}, hiệu quả thay đổi theo
dạng hàm độ trễ: \MaxSAT nổi bật với chi phí bậc thang và làm tròn, còn MIP cạnh
tranh hơn với chi phí liên tục. Giá trị thiết kế của DDD nằm ở khả năng chỉ mở
rộng miền tại nơi nghiệm hạn chế thực sự thiếu thông tin.
\end{resultbox}

\section{Chọn SAT, SMT, CP hay MIP}

Không có bộ giải thống trị mọi cấu trúc. Một quy tắc thực dụng là nhìn vào phần
khó nhất của mô hình:
\begin{itemize}
  \item Chọn \SAT/\MaxSAT khi lôgic rời rạc, cửa sổ và xung đột lặp lại; có lợi
        từ học mệnh đề và giải tăng dần.
  \item Chọn SMT khi bất đẳng thức thời gian nguyên là lõi, còn phần Boolean
        chỉ quyết định thứ tự hoặc chế độ \parencite{bofill2020rcpsp}.
  \item Chọn CP khi tích lũy/không chồng lấn là ràng buộc toàn cục tự nhiên và
        lan truyền miền có thể khai thác trực tiếp.
  \item Chọn MIP khi hàm mục tiêu tuyến tính mạnh và mô hình nới lỏng cung cấp cận tốt.
\end{itemize}

Phương pháp kết hợp nên được thiết kế như một phân chia nhiệm vụ suy diễn: mô-đun
nào tạo cận, mô-đun nào xử lý phép tuyển và bằng cách nào hai phía trao đổi
thông tin. Việc chỉ gọi nhiều bộ giải trên cùng một mô hình không tự tạo ra
một kiến trúc kết hợp.

\subsection{Một cây quyết định theo cấu trúc}

\begin{enumerate}
  \item Nếu khung thời gian nhỏ và mỗi tác vụ có ít giá trị bắt đầu, SAT/MaxSAT theo
        mốc thời gian là đối chứng minh bạch.
  \item Nếu khung thời gian lớn nhưng các ràng buộc trước--sau/sai phân dày, thử mã
        hóa thứ tự, SMT hoặc LCG trước khi mở toàn bộ trục thời gian.
  \item Nếu tài nguyên tích lũy là nút thắt và có suy luận năng lượng mạnh,
        CP/LCG có ưu thế cấu trúc mà xung đột từng cặp khó thay thế.
  \item Nếu hàm mục tiêu/mô hình nới lỏng tuyến tính tạo cận mạnh, MIP là đối
        chứng bắt buộc; SAT chỉ có lợi khi phép tuyển Boolean hoặc đối xứng chi phối.
  \item Nếu chỉ một phần nhỏ mốc thời gian quan trọng, xét DDD/tinh chỉnh và
        chứng minh rõ tính hợp lệ của cận dưới.
\end{enumerate}

\section{Bài học thiết kế}

\begin{summarybox}
\begin{enumerate}
  \item Phân biệt biến thời điểm, biến hoạt động và biến thứ tự; viết phép
        liên kết trước khi tối ưu mệnh đề.
  \item Tính không ngắt là một ràng buộc khối, không phải chỉ là ràng buộc đếm.
  \item Với các cửa sổ chồng lấn, hiệu quả đến từ kỹ thuật chia sẻ ngưỡng hoặc
        bộ đếm giữa các cửa sổ.
  \item Hàm mục tiêu nên giao tiếp với mô-đun khả thi qua một giao diện cận.
  \item Rời rạc hóa động phù hợp khi miền lớn nhưng nghiệm chỉ cần một số
        ít mốc quan trọng.
\end{enumerate}
\end{summarybox}

\subsection*{Câu hỏi nghiên cứu}
\begin{enumerate}
  \item Với cùng ngữ nghĩa khối, hãy so sánh biểu diễn bằng biến bắt đầu và biểu diễn
        bằng điểm chuyển về số biến, số mệnh đề và khả năng lan truyền đơn vị.
  \item Thiết kế một Shared Counter cho hai ALSC có cửa sổ chồng lấn nhưng cận
        \(k\) khác nhau.
  \item Viết điều kiện dừng chính xác cho một vòng DDD khi chi phí hạn chế là
        cận dưới và lịch giải mã cho cận trên.
\end{enumerate}

\chapter{Đóng gói và cắt hai chiều}
\chapterlead{Trong đóng gói, mệnh đề không làm việc trực tiếp với hình chữ
nhật; nó làm việc với tọa độ, hướng đặt và các quan hệ tách. Chất lượng phép mã hóa
phụ thuộc vào việc chọn đúng giao diện giữa ba lớp đó.}

\index{strip packing}\index{bin packing}\index{non-overlap}
\index{rotation}\index{coordinate compression}

\flowdiagram{Vật và hướng đặt}{Tọa độ thứ tự}{Quan hệ tách}
  {Cận khung chứa}{Bố trí và chứng nhận}

\section{Bản đồ các phương pháp chính xác cho đóng gói}

Hai chiều tạo khó khăn khác một chiều: không chồng lấn là một phép tuyển giữa
hai trục, còn tọa độ có nhiều giá trị đối xứng trong vùng trống. Văn liệu về
phương pháp chính xác vì thế không xoay quanh một mô hình duy nhất. Khảo cứu
của Iori và cộng sự phân biệt các mô hình theo vị trí đặt, vị trí tương đối,
sinh tập/cột, lập trình ràng buộc và các kỹ thuật nhánh--cận chuyên biệt
\parencite{iori2021packing}.

Đối với SAT, ba họ biểu diễn là đối chứng trực tiếp:
\begin{description}[style=nextline]
  \item[Chiếm chỗ]
  Biến cho từng ô mà vật chiếm. Không chồng lấn trở thành cục bộ nhưng số biến
  phụ thuộc diện tích lưới; tính liên tục của một hình cần các mệnh đề liên kết.

  \item[Vị trí đặt tuyệt đối]
  Mỗi biến chọn một vị trí khả dĩ của vật. Đồ thị xung đột nối hai vị trí chồng
  lấn; bài toán trở thành chọn đúng một vị trí cho mỗi vật. Các phương pháp vị
  trí tuyệt đối và đóng gói hoàn hảo khai thác nén tọa độ, tổng tập con và suy
  luận trên không gian trống
  \parencite{huang2013rectangle}.

  \item[Vị trí tương đối + tọa độ]
  Mỗi vật có tọa độ; mỗi cặp có biến chứng trái/phải/trên/dưới. Mã hóa thứ tự
  nén cận tọa độ, còn mệnh đề quan hệ xử lý phép tuyển.
\end{description}

Ngoài SAT, lập trình ràng buộc dùng \textsc{NoOverlap}/diffn, các phép chiếu
tích lũy và suy luận năng lượng; MIP dùng hằng số lớn \(M\), bất đẳng thức
tuyển hoặc mẫu phủ tập. Các phương pháp chính xác cho đóng gói trực giao còn
kết hợp cận, quan hệ trội và nhánh--cận chuyên dụng
\parencite{clautiaux2007orthogonal,clautiaux2008cp}.

\begin{center}
\captionof{table}{Các biểu diễn chính xác cho đóng gói hai chiều.}
\label{tab:packing-formulations}
\begin{tabularx}{\textwidth}{@{}>{\bfseries\sffamily}p{.2\textwidth}YYY@{}}
\toprule
Biểu diễn & Quy mô theo & Lan truyền tự nhiên & Phù hợp khi\\
\midrule
Chiếm chỗ & Diện tích lưới & Xung đột ô & Khung chứa nhỏ/rất rời rạc\\
Vị trí đặt & Số vị trí khả dĩ & Đồ thị xung đột & Miền đã nén mạnh\\
Tọa độ--thứ tự & Tổng độ rộng miền & Cận và quan hệ tách & Miền nguyên vừa\\
CP toàn cục & Số vật + bộ lan truyền & Không chồng lấn/tích lũy & Cận hình học mạnh\\
MIP/phủ tập & Biến/cột vị trí & Thư giãn và nhát cắt & Cận LP hữu ích\\
\bottomrule
\end{tabularx}
\end{center}

Ba trục tạo nên đóng góp SAT cho đóng gói là nén miền tọa độ, làm mạnh lan
truyền không chồng lấn và tái sử dụng mệnh đề đã học giữa các cận. Kiến trúc
mới được nhận diện qua thay đổi có thể đo trên ít nhất một trong ba trục này.

\section{Mô hình hình học rời rạc}

Mỗi vật \(i\) là hình chữ nhật kích thước \(w_i\times h_i\). Tọa độ
\((x_i,y_i)\) chỉ góc trái dưới. Khi cho phép quay \(90^\circ\), biến
\(\rho_i\) xác định kích thước hiệu dụng
\[
  (\widehat w_i,\widehat h_i)=
  \begin{cases}
    (w_i,h_i), & \rho_i=0,\\
    (h_i,w_i), & \rho_i=1.
  \end{cases}
\]
Trong khung chứa \(W\times H\), điều kiện nằm gọn là
\[
  0\le x_i\le W-\widehat w_i,
  \qquad
  0\le y_i\le H-\widehat h_i.
\]

Hai vật \(i\) và \(j\) không giao nhau khi ít nhất một trong bốn quan hệ đúng:
\[
  x_i+\widehat w_i\le x_j,\quad
  x_j+\widehat w_j\le x_i,\quad
  y_i+\widehat h_i\le y_j,\quad
  y_j+\widehat h_j\le y_i.
\]
Ta dùng các biến chứng \(l_{ij}\) và \(u_{ij}\):
\[
  l_{ij}\lor l_{ji}\lor u_{ij}\lor u_{ji},
\]
với quan hệ kéo theo từ từng biến chứng sang bất đẳng thức tương ứng. Các biến
chứng không cần loại trừ lẫn nhau: một vật có thể vừa ở bên trái vừa ở dưới vật
khác.

\begin{designrule}{Quy ước biên}
Tọa độ biểu diễn cạnh của hình chữ nhật, không phải tâm hay ô đã chiếm. Vì thế
dấu \(\le\) cho phép hai vật chạm cạnh. Quy ước này phải được giữ nguyên trong
điều kiện nằm gọn, không chồng lấn và bộ kiểm tra.
\end{designrule}

Hình \ref{fig:packing-layout} dùng một ví dụ tự dựng để minh họa cách tọa độ
thứ tự và quan hệ vị trí được dùng trong SAT cho đóng gói; cách tổ chức quan hệ
này tiếp nối mô hình đóng gói dải của \textcite{soh2010strip}.

\begin{figure}[tb]
\centering
\begin{tikzpicture}[satfig,x=8mm,y=8mm]
  \draw[step=1,Rule!35,very thin] (0,0) grid (8,6);
  \draw[Ink,very thick] (0,0) rectangle (8,6);
  \draw[Blue,fill=PaleBlue,thick] (0,0) rectangle (3,2);
  \node at (1.5,1) {\(A:3\times2\)};
  \draw[Teal,fill=PaleTeal,thick] (3,0) rectangle (5,4);
  \node[rotate=90] at (4,2) {\(B:2\times4\)};
  \draw[Gold,fill=PaleGold,thick] (5,0) rectangle (8,3);
  \node at (6.5,1.5) {\(C:3\times3\)};
  \draw[Blue!65!Teal,fill=SoftGray,thick] (0,4) rectangle (5,6);
  \node at (2.5,5) {\(D:5\times2\)};
  \foreach \t in {0,...,8}
    \node[below,satlabel] at (\t,0) {\(\t\)};
  \foreach \t in {1,...,6}
    \node[left,satlabel] at (0,\t) {\(\t\)};
  \draw[satdim] (0,6.45)--(8,6.45)
    node[midway,above,satlabel]{\(W=8\)};
  \draw[satdim] (8.45,0)--(8.45,6)
    node[midway,right,satlabel]{\(H=6\)};
\end{tikzpicture}
\caption{Một bố trí nguyên có bốn vật. Các cặp chạm cạnh vẫn hợp lệ; chẳng hạn
\(x_A+w_A=x_B=3\), nên biến chứng \(l_{AB}\) có thể nhận giá trị đúng.}
\label{fig:packing-layout}
\end{figure}

\begin{workedexample}{Từ một cặp hình đến literal quan hệ}
Trong \cref{fig:packing-layout}, \(A\) có \((x_A,y_A)=(0,0)\),
\((w_A,h_A)=(3,2)\); \(B\) có \((x_B,y_B)=(3,0)\),
\((w_B,h_B)=(2,4)\). Bốn bất đẳng thức tách là
\[
  0+3\le3,\quad 3+2\le0,\quad 0+2\le0,\quad 0+4\le0.
\]
Chỉ bất đẳng thức đầu đúng, nên \(l_{AB}=1\) là biến chứng tự nhiên. Nếu ép
\(x_B\le2\), miền tọa độ làm \(x_A+w_A\le x_B\) bất khả và mệnh đề liên kết
buộc \(\neg l_{AB}\); phép tuyển bốn quan hệ khi đó phải chọn một hướng khác
hoặc chứng minh cặp hình không thể cùng một bố trí.
\end{workedexample}

\section{Mã hóa thứ tự cho tọa độ}

\subsection{Ngữ nghĩa}

Với \(e\) trong miền ngang, đặt
\[
  p^x_{i,e}=1 \iff x_i\le e.
\]
Tương tự, \(p^y_{i,f}=1\iff y_i\le f\). Chuỗi phải đơn điệu:
\[
  p^x_{i,e}\rightarrow p^x_{i,e+1},
  \qquad
  p^y_{i,f}\rightarrow p^y_{i,f+1}.
\]
Giá trị tọa độ được xác định bởi điểm chuyển đầu tiên từ sai sang đúng. Một
bất đẳng thức như \(x_i+a\le x_j\) trở thành quan hệ kéo theo giữa hai ngưỡng:
\[
  p^x_{i,e}\ \text{và}\ p^x_{j,e+a-1}
\]
được nối theo hướng phù hợp với ngữ nghĩa \(\le\). Lợi ích là một literal
ngưỡng đại diện cho cả nửa miền, thay vì liệt kê từng cặp tọa độ.

\subsection{Hướng đặt và liên kết biểu diễn}

Khi kích thước phụ thuộc \(\rho_i\), quan hệ kéo theo
\[
  l_{ij}\rightarrow x_i+\widehat w_i\le x_j
\]
được tách theo các hướng đặt của \(i\) và \(j\). Các trường hợp không thể nằm
gọn được loại ở bước tiền xử lý. Với vật vuông, đặt
\(\rho_i=0\) để bỏ đối xứng quay.

Khi \(l_{ij}\) đóng vai trò biến chứng cho phép tuyển, mô-đun cần chiều
\[
  l_{ij}\rightarrow x_i+\widehat w_i\le x_j.
\]
Chiều ngược
\[
  x_i+\widehat w_i\le x_j\rightarrow l_{ij}
\]
có thể được bổ sung khi literal quan hệ cần đặc tả hai chiều. Bộ giải mã đọc
trực tiếp tọa độ và hướng đặt, nhờ đó bộ kiểm tra độc lập với lựa chọn liên kết.

\begin{proposition}[Tính đúng của mô-đun không chồng lấn]
Giả sử mỗi biến chứng \(l_{ij},u_{ij}\) kéo theo đúng bất đẳng thức hình học
của nó, và phép tuyển bốn biến chứng được thỏa cho mọi cặp vật trong cùng khung
chứa. Khi đó mọi bố trí giải mã từ mô hình \SAT đều không giao nhau.
\end{proposition}

\begin{proof}
Với một cặp \(i,j\), phép tuyển cho ít nhất một biến chứng đúng. Quan hệ kéo
theo của biến chứng đó cho một bất đẳng thức tách theo trục \(x\) hoặc \(y\). Vì vậy nội
thất hai hình chữ nhật rời nhau.
\end{proof}

\section{Bài toán đóng gói dải hai chiều}

\subsection{Từ tối ưu chiều cao sang kiểm tra tính khả thi}

Dải có chiều rộng cố định \(W\), còn chiều cao \(H\) cần tối thiểu hóa. Với
mỗi \(H\), dựng công thức
\[
  \Phi_{\mathrm{SPP}}(H)
  \iff
  \text{tồn tại bố trí trong }W\times H.
\]
Tính đơn điệu cho phép tìm kiếm trên cận:
\[
  \Phi_{\mathrm{SPP}}(H)=\SAT
  \Longrightarrow
  \Phi_{\mathrm{SPP}}(H')=\SAT
  \quad\forall H'\ge H.
\]
Cận dưới cơ bản là
\[
  LB=\max\left\{
    \max_i \min(h_i,w_i),
    \left\lceil\frac{\sum_iw_ih_i}{W}\right\rceil
  \right\},
\]
trong đó hạng đầu phải điều chỉnh theo hướng đặt thực sự vừa khung. Cận trên
đến từ một phép đặt theo kinh nghiệm.

Chỉ cận diện tích hiếm khi đủ sát. Các phương pháp chính xác còn dùng hàm khả
thi đối ngẫu, phép chiếu một chiều, tính không tương thích và tọa độ tổng tập
con. Trong SAT, mọi cận dưới hợp lệ đều có hai vai trò: thu hẹp số cận phải gọi
và thu hẹp miền tọa độ trước khi sinh mệnh đề. Nén tọa độ đặc biệt quan trọng:
nếu có thể chứng minh một vật chỉ cần bắt đầu tại các tổng tập con khả dĩ,
chuỗi thứ tự nên chạy trên tập mốc đã nén thay
vì mọi số nguyên từ \(0\) đến \(W-w_i\).

\subsection{Ba chế độ giải}

\begin{center}
\captionof{table}{Ba chế độ tối ưu chiều cao bằng SAT/MaxSAT.}
\label{tab:strip-solving-modes}
\begin{tabularx}{\textwidth}{@{}>{\bfseries\sffamily}p{.23\textwidth}YY@{}}
\toprule
Chế độ & Cấu trúc & Đặc điểm\\
\midrule
\SAT độc lập & Dựng lại \(\Phi(H)\) ở mỗi cận &
Đơn giản, mỗi cận có một \CNF riêng\\
\SAT tăng dần & Miền lớn cố định, kích hoạt \(H\) bằng các giả định &
Giữ mệnh đề đã học giữa các cận\\
\MaxSAT có trọng số & Chiều cao được biểu diễn bằng mệnh đề mềm &
Không cần vòng tìm kiếm ngoài\\
\bottomrule
\end{tabularx}
\end{center}

Mã hóa SAT thứ tự cho đóng gói dải đã được nghiên cứu trước đó cùng đối xứng,
quan hệ vị trí và tái sử dụng mệnh đề đã học
\parencite{soh2010strip}. Vì vậy nghiên cứu tình huống của nhóm được định vị
theo hai câu hỏi kế tiếp: phép quay/đối xứng làm đổi mô hình thế nào, và
SAT--MaxSAT--giải tăng dần có giữ lợi thế trên bộ kiểm chuẩn và bộ giải hiện
tại hay không. Phép so sánh hợp lý phải giữ cùng các cận, miền tọa độ và bộ
kiểm tra.

Trong nghiên cứu của \textcite{kieu2025strip}, \SAT không tăng dần với phép
quay và phá đối xứng chứng nhận 26 giá trị tối ưu trên 41 bài toán; phiên bản
tăng dần chứng nhận 21, còn CP chứng nhận 28. Kết quả này cho thấy giải tăng
dần chỉ có lợi khi công thức cơ sở và mệnh đề đã học còn phù hợp sau khi thay
cận. Với đóng gói dải, thay \(H\) có thể làm thay đổi mạnh miền tọa độ và
những xung đột hữu ích.

\begin{keyidea}{Chứng nhận tối ưu}
Một bố trí tại \(H^\star\) là cận trên. Công thức \(\Phi(H^\star-1)\) không
thỏa là chứng nhận cận dưới. Cặp SAT/UNSAT ở hai cận kề nhau mới xác định
\(H^\star\) là giá trị tối ưu.
\end{keyidea}

\section{Bài toán đóng gói vào khung chứa hai chiều}

\subsection{Hai kiến trúc}

Bài toán dùng nhiều khung chứa \(W\times H\) đồng nhất và tối thiểu hóa số
khung. Có hai kiến trúc chính.

\paragraph{Xếp chồng.}
Ghép \(k\) khung thành một dải \(W\times kH\). Tọa độ \(y_i\) đồng thời xác
định khung. Kiến trúc này tái sử dụng đóng gói dải nhưng phải ngăn vật vượt qua
đường biên \(qH\).

\paragraph{Không xếp chồng.}
Dùng biến \(a_{ib}=1\) khi vật \(i\) nằm trong khung \(b\), và tọa độ cục bộ
trong từng khung. Ta có
\[
  \ExactlyOne(a_{i1},\ldots,a_{ik}).
\]
Điều kiện không chồng lấn của cặp \(i,j\) trong khung \(b\) được gác bởi:
\[
  (a_{ib}\land a_{jb})
  \rightarrow
  (l_{ij}\lor l_{ji}\lor u_{ij}\lor u_{ji}).
\]
Biến sử dụng \(q_b\) nối với phép gán bằng \(a_{ib}\rightarrow q_b\).
Đối xứng giữa các khung được giảm bằng
\[
  q_{b+1}\rightarrow q_b.
\]

Các khung đã dùng tạo thành một tiền tố, vì thế hàm mục tiêu \(\sum_bq_b\)
đúng bằng số khung. Có thể thêm phép gán chuẩn cho một số vật đầu, nhưng mọi
ràng buộc đối xứng phải dựa trên phép hoán vị các khung đồng nhất, không dựa
trên một bố trí kinh nghiệm cụ thể.

\begin{figure}[tb]
\centering
\begin{tikzpicture}[satfig,x=7mm,y=7mm]
  \node[satlabel] at (2.5,4.9) {xếp chồng};
  \draw[Ink,thick] (0,0) rectangle (5,4);
  \draw[Rule,dashed] (0,2)--(5,2);
  \node[left,satlabel] at (0,1) {khung \(1\)};
  \node[left,satlabel] at (0,3) {khung \(2\)};
  \draw[Blue,fill=PaleBlue,thick] (.3,.3) rectangle (2.4,1.7);
  \draw[Teal,fill=PaleTeal,thick] (2.6,2.3) rectangle (4.7,3.7);
  \node[satlabel] at (9,4.9) {không xếp chồng};
  \draw[Ink,thick] (6.2,0) rectangle (10.2,2.6);
  \draw[Ink,thick] (10.8,0) rectangle (14.8,2.6);
  \node[below,satlabel] at (8.2,0) {khung \(1\)};
  \node[below,satlabel] at (12.8,0) {khung \(2\)};
  \draw[Blue,fill=PaleBlue,thick] (6.5,.3) rectangle (8.6,1.7);
  \draw[Teal,fill=PaleTeal,thick] (11.1,.3) rectangle (13.2,1.7);
  \node[satstate] at (8.2,3.5) {\(q_1\)};
  \node[satstate] at (12.8,3.5) {\(q_2\)};
  \draw[satdash] (12.25,3.5)--(8.75,3.5)
    node[midway,above,satlabel]{\(q_2\to q_1\)};
\end{tikzpicture}
\caption{Hai kiến trúc cho đóng gói vào khung: xếp chồng dùng một trục \(y\)
toàn cục; không xếp chồng dùng phép gán và tọa độ cục bộ, kèm tiền tố sử dụng.}
\label{fig:stacking-nonstacking}
\end{figure}

Trong mô hình theo vị trí đặt, một cách khác là xây đồ thị xung đột cho mọi cặp
vị trí--khung rồi thêm ràng buộc đúng một theo vật. Cách này tránh biến quan hệ
nhưng có thể tạo số vị trí rất lớn. Mô hình tọa độ không xếp chồng làm ngược
lại: số biến tọa độ nhỏ hơn toàn bộ tập vị trí, đổi lại mỗi cặp vật--khung cần
một phép tuyển có điều kiện. Hai mô hình là hai phía của phép phân tích nhân
tử; đánh giá thực nghiệm nên báo cả số vị trí bị loại ở tiền xử lý và số quan
hệ từng cặp còn lại.

\begin{resultbox}{So sánh kiến trúc}
Nghiên cứu của \textcite{kieu2025bin} so sánh các mô hình \SAT với MIP và CP
trên cùng họ bài toán kiểm chuẩn. Kết quả nhấn mạnh hai yếu tố hơn là một phép
mã hóa đơn lẻ: biến thứ tự giúp biểu diễn bất đẳng thức tọa độ, còn phá đối
xứng cho khung/bản sao quyết định đáng kể kích thước không gian tìm kiếm.
\end{resultbox}

\section{Bài toán cắt tấm một kích thước}

\subsection{Mô hình theo số tờ}

Cho tấm vật liệu kích thước \(W\times H\), nhiều loại vật \(t\), mỗi loại có
nhu cầu \(d_t\). Bài toán tối thiểu hóa số tấm. Nếu tạo từng bản sao riêng,
phép gán và không chồng lấn giống bài toán đóng gói vào khung. Nếu gộp theo
loại, cần thêm ràng buộc đếm để bảo đảm đúng nhu cầu.

Với một cận \(K\), công thức khả thi hỏi liệu có thể đáp ứng đủ nhu cầu trên
\(K\) tấm. Hàm mục tiêu là
\[
  \min\set{K\mid \Phi_{\mathrm{CSP}}(K)\text{ là SAT}}.
\]
Đối xứng xuất hiện ở hai tầng: các tấm đồng nhất và các bản sao cùng loại. Một
cách phá đối xứng bền vững là buộc biến sử dụng tấm theo tiền tố và buộc chỉ
số tấm của các bản sao cùng loại không giảm.

\subsection{SAT và MaxSAT}

Tìm kiếm SAT thay \(K\) ở vòng ngoài. \MaxSAT bộ phận có thể giữ tất cả các
tấm ứng viên, đặt mệnh đề mềm \(\neg q_b\) và tối thiểu
\(\sum_bq_b\). Hai cách dùng chung điều kiện nằm gọn, hướng đặt và không chồng
lấn; chúng chỉ khác ở lớp hàm mục tiêu.

\textcite{kieu2026cutting} cho thấy phép mã hóa gọn với tọa độ thứ tự và đối
xứng theo bản sao/tấm đạt kết quả cạnh tranh trên bộ bài toán cắt tấm một kích
thước. Điểm chuyển giao sang các bài toán đóng gói khác là chuỗi mô-đun
\[
  \texttt{Assignment}
  +\texttt{OrderCoordinate}
  +\texttt{ConditionalNonOverlap}
  +\texttt{UsageCounter}.
\]

\section{Giảm miền và đối xứng}

Trước khi sinh \CNF, bốn phép rút gọn thường có hiệu quả trực tiếp:
\begin{enumerate}
  \item loại hướng đặt không vừa khung;
  \item thu hẹp \(x_i\le W-\widehat w_i\) và
        \(y_i\le H-\widehat h_i\);
  \item loại một hướng tách nếu các cận chứng minh hướng đó bất khả thi;
  \item chuẩn hóa vật/bản sao và khung chứa theo một thứ tự cố định.
\end{enumerate}

Mỗi ràng buộc phá đối xứng được gắn với một nhóm đối xứng cụ thể:
\begin{itemize}
  \item \(\rho_i=0\) cho vật vuông: đối xứng hướng đặt;
  \item \(q_{b+1}\rightarrow q_b\): hoán vị khung/tấm;
  \item thứ tự bản sao cùng loại: hoán vị các bản sao đồng nhất;
  \item neo một vật lớn vào nửa khung chứa: phản xạ/tịnh tiến, nếu phép biến
        đổi giữ miền và hàm mục tiêu.
\end{itemize}
Các nguyên tắc tổng quát về đối xứng được trình bày trong
\textcite{walsh2006symmetry,gent2006symmetry}; phép mã hóa SAT cổ điển cho bài toán đóng gói khung 2D
Packing có rotation có thể xem ở \textcite{soh2010binpacking}.

\section{Tách ảnh hưởng của phép mã hóa, cận và quan hệ trội}

Ma trận thiết kế ghi riêng ảnh hưởng của các thành phần:
\begin{itemize}
  \item cận dưới theo diện tích, phép chiếu và đồ thị con đầy đủ của quan hệ
        không tương thích;
  \item loại hướng đặt và vị trí bị trội;
  \item cố định theo vật lớn hoặc thứ tự vật đồng nhất;
  \item nén tọa độ;
  \item khởi động ấm/cận trên từ phương pháp kinh nghiệm.
\end{itemize}

Ma trận thiết kế tách phần giảm mệnh đề do lọc miền khỏi phần giảm do chuyển
sang CNF. Phép so sánh CP, MIP và SAT dùng chung cách chuẩn hóa bài toán và cận
đầu vào; cận và quan hệ trội được ghi thành mô-đun riêng
\parencite{iori2021packing,clautiaux2007orthogonal}.

\section{Một thư viện mô-đun cho hình học}

\begin{summarybox}
\textbf{OrderCoordinate} cung cấp chuỗi ngưỡng và phép giải mã duy nhất.
\textbf{Orientation} trả về kích thước hiệu dụng.
\textbf{Containment} siết miền theo khung chứa đang chọn.
\textbf{RelativePosition} tạo biến chứng trái/phải/trên/dưới.
\textbf{ConditionalNonOverlap} gác phép tuyển bằng phép gán.
\textbf{UsageCounter} biến số khung chứa được dùng thành hàm mục tiêu.
\end{summarybox}

Thứ tự kết hợp đi từ hướng đặt đến điều kiện nằm gọn, vị trí tương đối, không
chồng lấn được gác bởi phép gán và bộ đếm sử dụng. Bộ kiểm tra đọc vật, hướng
đặt, tọa độ và khung chứa, độc lập với biến phụ.

\section{Bài học thiết kế}

\begin{enumerate}
  \item Chọn một quy ước tọa độ và giữ nó nhất quán trong toàn bộ phép mã hóa.
  \item Dùng biến chứng cho phép tuyển hình học; không buộc biến chứng trở thành
        quan hệ đầy đủ nếu phép giải mã không cần.
  \item Tách mô-đun khả thi khỏi hàm mục tiêu số khung/tấm/chiều cao.
  \item Phá đối xứng theo đúng phép hoán vị hoặc phản xạ của mô hình.
  \item Chứng nhận tối ưu gồm bố trí hợp lệ và một cận dưới tương ứng.
\end{enumerate}

\subsection*{Câu hỏi nghiên cứu}
\begin{enumerate}
  \item Dẫn xuất các mệnh đề thứ tự cho quan hệ kéo theo
        \(l_{ij}\rightarrow x_i+\widehat w_i\le x_j\) ở một hướng đặt cố định.
  \item So sánh số biến của xếp chồng và không xếp chồng trên \(n\) vật,
        \(k\) khung.
  \item Thiết kế ràng buộc phá đối xứng cho ba bản sao cùng loại và chứng minh
        mỗi quỹ đạo giữ một
        đại diện.
\end{enumerate}

\chapter{Bandwidth, antibandwidth và tô đa màu}
\chapterlead{Các bài toán trong chương đều nói về khoảng cách giữa nhãn, nhưng
không chia sẻ một ngữ nghĩa: antibandwidth gán một hoán vị, tô màu theo
bandwidth gán một màu, còn tô đa màu gán một tập màu có thứ tự.}

\index{antibandwidth}\index{cyclic antibandwidth}
\index{bandwidth coloring}\index{bandwidth multicoloring}
\index{forbidden interval}

\flowdiagram{Hoán vị hay tô màu}{Cửa sổ/thứ tự/khối}
  {Khoảng cách}{Tìm cận}{Gán nhãn và chứng nhận}

\section{Bố trí đồ thị, gán nhãn và tô màu}

Từ ``bandwidth'' được dùng cho nhiều bài toán có công thức gần nhau nhưng lịch
sử thuật toán khác nhau. Bandwidth đồ thị cổ điển tối thiểu hóa khoảng cách
lớn nhất trên cạnh của một hoán vị đỉnh. Antibandwidth tối đa hóa khoảng cách
nhỏ nhất. Tô màu theo bandwidth bỏ tính song ánh và cho trọng số khoảng cách
trên cạnh; tô đa màu cho mỗi đỉnh nhiều tần số. Văn liệu gán nhãn đồ thị rộng
hơn còn chứa hàng trăm quy tắc gán số khác nhau
\parencite{gallian2024labeling}.

Ba truyền thống giải liên quan trực tiếp:
\begin{itemize}
  \item \textbf{cận lý thuyết đồ thị}: bậc, số ổn định, số sắc,
        \emph{clique} (đồ thị con đầy đủ) và cấu
        trúc lớp đồ thị;
  \item \textbf{phương pháp kinh nghiệm/siêu kinh nghiệm}: tìm kiếm cục bộ,
        \emph{tabu search} (tìm kiếm kiêng kỵ), lân cận biến đổi và thứ tự kiến tạo để tìm
        nghiệm đương nhiệm cho
        cận trên/cận dưới;
  \item \textbf{phương pháp chính xác}: quy hoạch động, nhánh--cận, MIP, CP và
        SAT theo cận.
\end{itemize}

SAT giải hiệu quả bài toán quyết định ``có phép gán nhãn tại cận \(k\) không''
và kết hợp tự nhiên với các cận lý thuyết đồ thị. Các cận thu hẹp khoảng tìm
kiếm, tạo các giả định đầu tiên tốt hơn và đôi khi trực tiếp chứng nhận tối ưu.
Với BCP/BMCP, khảo cứu tô màu đỉnh và các biến thể cung cấp mô hình, cận theo
đồ thị con đầy đủ
và nguồn gốc bộ kiểm chuẩn cần dùng làm đối chứng
\parencite{malaguti2010coloring}.

\begin{center}
\captionof{table}{Bốn họ bài toán khoảng cách nhãn.}
\label{tab:distance-families}
\begin{tabularx}{\textwidth}{@{}>{\bfseries\sffamily}p{.2\textwidth}YYY@{}}
\toprule
Họ & Phép gán & Cận miền tự nhiên & Đối chứng chính xác\\
\midrule
Bandwidth & Hoán vị & Cận bố trí/bậc & DP, B\&B, MIP\\
Antibandwidth & Hoán vị & Cận ổn định/số sắc & MIP, CP, SAT\\
BCP & Một màu/đỉnh & Cận đồ thị con đầy đủ và tô màu & MIP, CP, SAT\\
BMCP & Nhiều màu/đỉnh & Cận ô màu/nhu cầu & MIP, tìm kiếm cục bộ, SAT\\
\bottomrule
\end{tabularx}
\end{center}

\section{Phân loại trước khi biểu diễn}

\begin{center}
\captionof{table}{Phân biệt ngữ nghĩa của ABP, CABP, No-hole Anti-\(k\)-labeling, BCP và BMCP.}
\label{tab:distance-semantics}
\begin{tabularx}{\textwidth}{@{}>{\bfseries\sffamily}p{.24\textwidth}YYY@{}}
\toprule
Bài toán & Nghiệm & Ràng buộc & Mục tiêu\\
\midrule
Antibandwidth & Hoán vị \(1,\ldots,n\) &
\(\card{f(u)-f(v)}\ge k\) & Tối đa \(k\)\\
Antibandwidth tuần hoàn & Hoán vị trên vòng &
\(D_c(f(u),f(v))\ge k\) & Tối đa \(k\)\\
No-hole Anti-\(k\)-labeling & Ánh xạ toàn ánh \(\set{1,\ldots,k}\) (\(k<n\)) &
\(\card{\psi(u)-\psi(v)}\ge w\) và \(\sum_i x_{ij}\ge 1\) & Tối đa \(w\)\\
Tô màu theo bandwidth & Một màu mỗi đỉnh &
\(\card{c(u)-c(v)}\ge d_{uv}\) & Tối thiểu độ trải\\
Tô đa màu theo bandwidth & \(w(u)\) màu mỗi đỉnh &
Khoảng cách nội/ngoại đỉnh & Tối thiểu màu lớn nhất\\
\bottomrule
\end{tabularx}
\end{center}

Với ABP và CABP, nhãn là một hoán vị nên ma trận gán cần
\(\ExactlyOne\) theo cả hàng và cột. Với BCP, hai đỉnh không kề nhau có thể
dùng cùng màu; chỉ hàng cần \(\ExactlyOne\). Với BMCP, một đỉnh có nhiều ô
màu; thứ tự các ô là một phần của mô hình.

\begin{designrule}{Kiểm tra ngữ nghĩa trước AMO}
Hai phép biểu diễn có cùng biến \(x_{v,\ell}\) vẫn có thể mô tả hai bài toán khác
nhau. Hãy viết phép giải mã và ràng buộc đếm theo hàng/cột trước khi thêm các
ràng buộc khoảng cách.
\end{designrule}

\section{Antibandwidth và cửa sổ staircase}

\subsection{Cận và phương pháp chính xác ngoài SAT}

Với đồ thị liên thông, các cận trên ABP có thể lấy từ số đỉnh/cạnh, bậc, số ổn
định \(\alpha(G)\) và số sắc \(\chi(G)\). Chẳng hạn
\[
  AB(G)\le \alpha(G),
  \qquad
  AB(G)\le
  \left\lfloor\frac{n-1}{\chi(G)-1}\right\rfloor.
\]
Tính chính xác \(\alpha,\chi\) cũng khó, nhưng một cận trên cho
\(\alpha\) hoặc cận dưới cho \(\chi\) vẫn có thể tạo cận hợp lệ. Các mô hình
MIP, nhánh--cắt và CP đã cải thiện đáng kể chứng nhận trên bộ Harwell--Boeing
\parencite{sinnl2021antibandwidth}. Do đó một nghiên cứu SAT trên cùng bộ kiểm
chuẩn cần báo cùng cận đầu vào và so cả giá trị nghiệm lẫn khoảng cách chứng
minh, không chỉ số bài toán có mô hình.

\subsection[Tính khả thi ở cận k]{Tính khả thi ở cận \(k\)}

Cho đồ thị \(G=(V,E)\), \(\card{V}=n\). Biến
\[
  x_{v,\ell}=1\iff f(v)=\ell
\]
mô tả một hoán vị:
\[
  \ExactlyOne(x_{v,1},\ldots,x_{v,n})
  \quad\forall v,
\]
\[
  \ExactlyOne(x_{1,\ell},\ldots,x_{n,\ell})
  \quad\forall\ell.
\]
Với một cận \(k\), mỗi cạnh \((u,v)\) cấm hai nhãn cách nhau dưới \(k\):
\[
  \neg x_{u,a}\lor\neg x_{v,b},
  \qquad \card{a-b}<k.
\]

Một cách nhìn khác cố định vị trí \(a\) của \(u\). Khi đó \(v\) không được
nằm trong cửa sổ
\[
  [a-k+1,a+k-1]\cap[1,n].
\]
Khi \(a\) trượt, các cửa sổ chồng lấn tạo một họ AMO bậc thang. Nếu biểu diễn
từng cửa sổ độc lập, phần lớn bộ đếm bị lặp.

\begin{figure}[tb]
\centering
\begin{tikzpicture}[satfig,x=8mm,y=7mm]
  \node[satlabel] at (2.7,2.3) {một phép gán nhãn cho \(P_4\)};
  \foreach \i/\x/\lab in {1/0/2,2/1.8/4,3/3.6/1,4/5.4/3}{
    \node[satstate] (v\i) at (\x,1) {\(v_{\i}\)};
    \node[below=2pt of v\i,satlabel,text=Blue] {\(f(v_{\i})=\lab\)};
  }
  \draw[Ink,thick] (v1)--(v2)--(v3)--(v4);
  \node[satlabel] at (9,2.3) {cặp nhãn bị cấm khi \(k=2\)};
  \foreach \a in {1,...,4}{
    \node[left,satlabel] at (7.8,{1.6-.55*\a}) {\(\a\)};
    \node[above,satlabel] at ({7.8+.55*\a},1.6) {\(\a\)};
    \foreach \b in {1,...,4}{
      \pgfmathtruncatemacro{\delta}{abs(\a-\b)}
      \ifnum\delta<2
        \draw[Gold,fill=PaleGold] ({7.8+.55*\b},{1.6-.55*\a})
          rectangle ++(.5,-.5);
      \else
        \draw[Rule,fill=white] ({7.8+.55*\b},{1.6-.55*\a})
          rectangle ++(.5,-.5);
      \fi
    }
  }
  \node[satlabel,text=Gold] at (9,-1.65)
    {ô tô màu: \(|a-b|<2\)};
\end{tikzpicture}
\caption{Ở trái, nhãn \((2,4,1,3)\) cho \(P_4\) đạt antibandwidth \(2\).
Ở phải, mỗi ô được tô biểu diễn một cặp nhãn bị cấm trên một cạnh.}
\label{fig:abp-path-conflicts}
\end{figure}

\subsection{Thang Sequential Counter}

SCL chia dãy biến thành các khối, xây Sequential Counter trong mỗi khối rồi nối
trạng thái ở biên. Ý tưởng cốt lõi là biến phụ mang nghĩa số đếm tiền tố:
\[
  s_{i,j}=1\iff
  \sum_{q=1}^{i}x_q\ge j.
\]
Các cửa sổ chia sẻ tiền tố/hậu tố thay vì dựng lại bộ đếm. Với trường hợp
\(k=1\), đây là phép biểu diễn cho AMO bậc thang; với \AMK, bộ đếm giữ nhiều mức
\(j=1,\ldots,k\).

\begin{proposition}[Tái sử dụng qua biên khối]
Nếu bộ đếm trong mỗi khối đúng với mọi tiền tố cục bộ và bộ nối bảo toàn tổng
tích lũy khi đi qua biên, thì ngưỡng của mọi cửa sổ ghép từ hậu tố khối trái
và tiền tố khối phải đúng với số literal đúng trong cửa sổ đó.
\end{proposition}

\begin{proof}
Tổng của cửa sổ bằng tổng trên hai phần rời nhau. Bộ đếm cục bộ biểu diễn đúng
hai tổng; bộ nối liệt kê các cặp ngưỡng có tổng vượt cận. Vì vậy một phép gán
vi phạm ràng buộc đếm kích hoạt ít nhất một mệnh đề nối, còn một phép gán hợp
lệ mở rộng được bằng các giá trị đếm tiền tố thật.
\end{proof}

\textcite{truong2025scamo} cho thấy SCL giảm cả biến phụ và mệnh đề so với
nhiều phép biểu diễn không chia sẻ trên họ AMO bậc thang. Hướng \AMK được phát triển trong
\textcite{truong2026ladderamk}; Duplex của \textcite{fazekas2020duplex} là một
đối chiếu quan trọng theo hướng BDD hai chiều. Một cách nhìn rộng hơn về biến
phụ có cấu trúc được thảo luận trong \textcite{haberlandt2023aux}.

SCL bổ sung cho các mô hình MIP/CP một cơ chế nén các khoảng xung đột trong
\CNF. Cận mạnh và SCL tương hỗ trực tiếp: \(k\) quyết định chiều rộng cửa sổ
cần biểu diễn, còn Shared Counter giảm phần biểu diễn lặp giữa các cửa sổ đó.

\section{Antibandwidth tuần hoàn}

\subsection{Khoảng cách vòng}

Trong CABP, nhãn nằm trên vòng \(C_n\):
\[
  D_c(a,b)=\min\{\card{a-b},\,n-\card{a-b}\}.
\]
Với cận \(k\), một cạnh cấm cửa sổ vòng bán kính \(k-1\). Cửa sổ có thể vắt
qua điểm cắt nên dùng chỉ số môđun:
\[
  W(a,k)=\set{a-k+1,\ldots,a+k-1}\pmod n.
\]
Xung đột theo nhãn chính xác vẫn đúng nhưng số mệnh đề lớn. Cửa sổ vòng tạo
AMO tuần hoàn, có thể biểu diễn bằng cách cắt vòng tại một vị trí neo rồi nối hai
đầu bằng bộ nối.

\begin{figure}[tb]
\centering
\begin{tikzpicture}[satfig]
  \foreach \a/\ang in {1/90,2/45,3/0,4/-45,5/-90,6/-135,7/180,8/135}{
    \ifnum\a<4
      \node[satstate,satwarn] (c\a) at (\ang:24mm) {\(\a\)};
    \else
      \ifnum\a>6
        \node[satstate,satwarn] (c\a) at (\ang:24mm) {\(\a\)};
      \else
        \node[satstate,sattrue] (c\a) at (\ang:24mm) {\(\a\)};
      \fi
    \fi
  }
  \draw[Rule,thick] (c1)--(c2)--(c3)--(c4)--(c5)--(c6)--(c7)--(c8)--cycle;
  \node[satlabel,align=center] at (0,0)
    {\(W(1,3)=\{7,8,1,2,3\}\)\\vắt qua điểm cắt};
  \draw[Gold,very thick,->] (c1) to[bend left=20] (c3);
  \draw[Gold,very thick,->] (c1) to[bend right=20] (c7);
\end{tikzpicture}
\caption{Cửa sổ cấm trên vòng \(C_8\) với tâm \(1\) và \(k=3\). Hai nhánh
\(7,8\) và \(2,3\) phải được nối khi tuyến tính hóa vòng.}
\label{fig:cabp-cyclic-window}
\end{figure}

\subsection{Đối xứng của vòng}

CABP bất biến dưới phép quay và phản xạ của nhãn. Ta có thể cố định một đỉnh
đại diện tại nhãn \(1\), rồi chọn hướng chuẩn cho một đỉnh thứ hai. Đây là phá
đối xứng ở cấp hoán vị; điều kiện này độc lập với SCL dùng cho các cửa sổ khoảng cách.

Tìm nghiệm tối ưu theo cận tăng:
\[
  F(k)=\SAT \Longrightarrow F(k-1)=\SAT,
  \qquad
  F(k)=\UNSAT \Longrightarrow F(k+1)=\UNSAT.
\]
Một nghiệm tại \(k^\star\) và \UNSAT tại \(k^\star+1\) chứng nhận tối ưu.

\begin{resultbox}{CABP và bộ kiểm chuẩn}
Nghiên cứu của \textcite{truong2026cabp} xây dựng các phép biểu diễn SAT cho CABP
và khai thác cửa sổ tuần hoàn cùng đối xứng của vòng. Một tập công thức khả thi được
chuẩn bị cho SAT Competition 2026 \parencite{truong2026competition}; về mặt
thiết kế, giá trị của bộ này là tách bài toán tối ưu thành chuỗi bài toán
\SAT/\UNSAT có cận \(k\) rõ ràng.
\end{resultbox}

\section{Bài toán No-hole Anti-\(k\)-labeling và nén khoảng cấm bằng biến tiền tố}

\index{No-hole Anti-k-labeling}\index{anti-k-labeling}
\index{prefix variables}\index{forbidden interval}

\subsection{Bài toán gán nhãn phi song ánh và điều kiện toàn ánh}

Trong quản lý phổ tần số và tài nguyên mạng không dây, số lượng kênh truyền dẫn khả thi \(k\) thường nhỏ hơn đáng kể so với số lượng trạm phát sóng \(n = \card{V}\) (\(k < n\)) do các giới hạn về mặt pháp lý và kỹ thuật. Bối cảnh này dẫn tới bài toán \emph{Anti-\(k\)-labeling} \parencite{guan2019anti}, một mở rộng phi song ánh tự nhiên của bài toán Antibandwidth cổ điển (nơi \(k=n\)). Khi \(k < n\), phép gán nhãn không còn là một hoán vị mà trở thành một phép phân hoạch các đỉnh vào các lớp nhãn.

Biến thể \emph{No-hole Anti-\(k\)-labeling} \parencite{wang2021anti,hai2026nohole} yêu cầu hàm gán nhãn \(\psi: V \to \set{1,\ldots,k}\) phải là một phép ánh xạ \emph{toàn ánh} (\emph{surjective mapping}), tức là mọi nhãn trong tập \(\set{1,\ldots,k}\) đều phải được gán cho ít nhất một đỉnh (không có nhãn trống/lỗ hổng). Mục tiêu của bài toán là tìm hàm gán nhãn \(\psi\) sao cho khoảng cách giữa hai đỉnh kề nhau là lớn nhất:
\begin{equation}
  \lambda_k^{nh}(G) = \max_{\psi} \min_{(u,v)\in E} \card{\psi(u) - \psi(v)}.
  \label{eq:nohole-objective}
\end{equation}

Mô hình quy hoạch tuyến tính số nguyên (ILP) cho bài toán quyết định ``liệu có tồn tại phép gán nhãn toàn ánh thỏa mãn khoảng cách tối thiểu \(w\) không'' sử dụng các biến nhị phân \(x_{ij} = 1 \iff \psi(i) = j\):
\begin{align}
  \sum_{j=1}^k x_{ij} &= 1 && \forall i\in\set{1,\ldots,n}, \label{eq:nohole-ilp-vertex}\\
  \sum_{i=1}^n x_{ij} &\ge 1 && \forall j\in\set{1,\ldots,k}, \label{eq:nohole-ilp-nohole}\\
  x_{uj} + \sum_{j'=\max(1, j-w+1)}^{\min(k, j+w-1)} x_{vj'} &\le 1 && \forall (u,v)\in E, j\in\set{1,\ldots,k}. \label{eq:nohole-ilp-distance}
\end{align}
Ràng buộc~\eqref{eq:nohole-ilp-vertex} bảo đảm mỗi đỉnh nhận đúng một nhãn; Ràng buộc~\eqref{eq:nohole-ilp-nohole} bảo đảm tính toàn ánh (\emph{no-hole}); và Ràng buộc~\eqref{eq:nohole-ilp-distance} cấm hai đỉnh kề nhau nhận các nhãn thuộc khoảng cấm \((j-w, j+w)\).

\subsection{Từ bùng nổ mệnh đề từng cặp đến biểu diễn khoảng bằng biến tiền tố}

Chuyển trực tiếp Ràng buộc khoảng cách~\eqref{eq:nohole-ilp-distance} sang \CNF bằng phép biểu diễn từng cặp (Direct Encoding - DE) sinh ra họ mệnh đề xung đột:
\begin{equation}
  \neg x_{ui} \lor \neg x_{vj} \qquad \forall (u,v)\in E, i\in\set{1,\dots,k}, j\in [\max(1,i-w+1),\min(k,i-w-1)].
  \label{eq:nohole-direct-cnf}
\end{equation}
Số lượng mệnh đề của DE phát triển theo cấp số \(O(\card{E}\cdot k\cdot w)\). Khi khoảng cách mục tiêu \(w\) tăng lên, số mệnh đề khoảng cách bùng nổ mạnh mẽ, gây tắc nghẽn bộ nhớ và làm suy giảm hiệu năng của bộ giải SAT.

Để giải quyết triệt để nút thắt này, phép biểu diễn \mbox{\textnormal{\textsc{PSE}}} (\emph{Proposed SAT Encoding}) \parencite{hai2026nohole} giới thiệu các biến phụ tiền tố \(l_{ij} \in \set{0,1}\) với ngữ nghĩa:
\[
  l_{ij} = 1 \iff \psi(i) \le j \iff \exists k' \le j : x_{ik'} = 1.
\]
Tính đơn điệu tự nhiên của biến tiền tố đòi hỏi \(l_{ij} \to l_{i(j+1)}\) cho mọi \(j < k\). Hệ thống mệnh đề liên kết biến gán gốc \(x_{ij}\) và biến tiền tố \(l_{ij}\) được định nghĩa bởi:
\begin{align}
  x_{ij} &\to l_{ij} && \forall i\in\set{1,\ldots,n}, j\in\set{1,\ldots,k}, \label{eq:nohole-link-1}\\
  l_{ij} &\to l_{i(j+1)} && \forall i\in\set{1,\ldots,n}, j\in\set{1,\ldots,k-1}, \label{eq:nohole-link-2}\\
  \neg x_{i1} &\to \neg l_{i1} && \forall i\in\set{1,\ldots,n}, \label{eq:nohole-link-3}\\
  (\neg x_{ij} \land \neg l_{i(j-1)}) &\to \neg l_{ij} && \forall i\in\set{1,\ldots,n}, j\in\set{2,\ldots,k}, \label{eq:nohole-link-4}\\
  x_{ij} &\to \neg l_{i(j-1)} && \forall i\in\set{1,\ldots,n}, j\in\set{2,\ldots,k}. \label{eq:nohole-link-5}
\end{align}

Nhờ cấu trúc biến tiền tố, toàn bộ khoảng cấm \((j-w, j+w)\) giữa hai đỉnh kề nhau \((u,v)\in E\) được nén lại chỉ trong **đúng một mệnh đề kéo theo đơn duy nhất**:
\begin{equation}
  x_{uj} \to \left( l_{v, \max(1, j-w)} \lor \neg l_{v, \min(k, j+w-1)} \right) \qquad \forall (u,v)\in E, j\in\set{1,\ldots,k}.
  \label{eq:nohole-prefix-distance}
\end{equation}
Mệnh đề~\eqref{eq:nohole-prefix-distance} khẳng định rằng: nếu đỉnh \(u\) nhận nhãn \(j\), thì đỉnh kề \(v\) bắt buộc phải nhận nhãn nhỏ hơn/bằng \(j-w\) (thỏa mãn literal \(l_{v, \max(1,j-w)}\)) hoặc phải nhận nhãn lớn hơn \(j+w-1\) (thỏa mãn literal \(\neg l_{v, \min(k,j+w-1)}\)). Điều này tự động loại bỏ tất cả các nhãn nằm trong khoảng cấm \([j-w+1, j+w-1]\) mà không cần liệt kê từng cặp nhãn vi phạm.

\begin{figure}[tb]
\centering
\begin{tikzpicture}[satfig,x=8.5mm,y=7.5mm]
  %--- Direct Encoding ---
  \node[anchor=west,font=\rmfamily\bfseries\small] at (0,3.5) {(a) Phép biểu diễn từng cặp (Direct Encoding)};
  \node[anchor=east,font=\rmfamily\small] at (0.6,2.1) {$x_{u,3}=1 \implies$};
  \foreach \j/\v/\c in {1/1/PaleTeal, 2/2/PaleGold, 3/3/PaleGold, 4/4/PaleGold, 5/5/PaleTeal} {
    \draw[satbox,fill=\c] (0.8+\j*1.1, 1.75) rectangle +(0.8,0.7);
    \node[font=\rmfamily\small] at (0.8+\j*1.1+0.4, 2.1) {$\v$};
  }
  \node[anchor=west,font=\rmfamily\small,text=Gold!80!black] at (0.8, 0.8) {Liệt kê 3 mệnh đề cấm độc lập: $\neg x_{v,2}$, $\neg x_{v,3}$, $\neg x_{v,4}$};

  %--- Separator ---
  \draw[Rule,thick] (7.2,0.1)--(7.2,3.6);

  %--- Proposed Prefix Encoding ---
  \node[anchor=west,font=\rmfamily\bfseries\small] at (7.8,3.5) {(b) Phép biểu diễn tiền tố khoảng cấm (PSE)};
  \node[anchor=east,font=\rmfamily\small] at (8.4,2.1) {$x_{u,3}=1 \implies$};
  
  \draw[satbox,fill=PaleTeal,draw=Teal,thick] (8.6,1.75) rectangle +(0.8,0.7);
  \node[font=\rmfamily\small] at (9.0,2.1) {$1$};
  \node[anchor=south,font=\rmfamily\small,text=Teal] at (9.0,2.55) {$l_{v,1}=1$};

  \draw[satbox,fill=SoftGray,dashed] (9.6,1.75) rectangle +(3.0,0.7);
  \node[font=\rmfamily\small,text=black!60] at (11.1,2.1) {Khoảng cấm $[2, 4]$};

  \draw[satbox,fill=PaleTeal,draw=Teal,thick] (12.8,1.75) rectangle +(0.8,0.7);
  \node[font=\rmfamily\small] at (13.2,2.1) {$5$};
  \node[anchor=south,font=\rmfamily\small,text=Teal] at (13.2,2.55) {$\neg l_{v,4}=1$};

  \node[anchor=west,font=\rmfamily\small,text=Teal] at (8.6,0.8) {Một mệnh đề đơn duy nhất: $x_{u,3} \to (l_{v,1} \lor \neg l_{v,4})$};
\end{tikzpicture}
\caption{So sánh cơ chế nén khoảng cấm khi $k=5, w=2$: (a) Phép biểu diễn từng cặp liệt kê $O(w)$ mệnh đề; (b) Phép biểu diễn tiền tố PSE loại toàn bộ khoảng cấm bằng một mệnh đề kéo theo duy nhất.}
\label{fig:nohole-prefix-interval}
\end{figure}

\begin{proposition}[Độ phức tạp và tính đúng đắn của PSE]
\label{prop:nohole-pse-complexity}
Phép biểu diễn \mbox{\textnormal{\textsc{PSE}}} biểu diễn chính xác bài toán No-hole Anti-\(k\)-labeling. Tổng số biến nhị phân là \(2nk\), tổng số mệnh đề khoảng cách là \(2\card{E}k\), và các mệnh đề liên kết là \(n(4k-2)\). Phép biểu diễn hoàn toàn triệt tiêu sự phụ thuộc của số mệnh đề vào tham số khoảng cách mục tiêu \(w\).
\end{proposition}

\begin{proof}
Mỗi đỉnh \(i\) có \(k\) biến gán gốc \(x_{ij}\) và \(k\) biến tiền tố \(l_{ij}\), tổng cộng \(2nk\) biến. Các Mệnh đề~\eqref{eq:nohole-link-1}--\eqref{eq:nohole-link-5} bảo đảm ngữ nghĩa \(l_{ij} = \sum_{k'\le j}x_{ik'}\) và ép buộc đỉnh \(i\) nhận đúng một nhãn mà không cần thêm các mệnh đề \textsc{ExactlyOne} bổ sung. Mệnh đề khoảng cách~\eqref{eq:nohole-prefix-distance} chỉ tạo ra 2 mệnh đề cho mỗi cặp cạnh định hướng và nhãn \(j\), cho tổng số \(2\card{E}k\) mệnh đề.
\end{proof}

\subsection{Phá vỡ đối xứng phản xạ và kết quả thực nghiệm}

Không gian nghiệm của bài toán No-hole Anti-\(k\)-labeling chứa đối xứng phản xạ: nếu \(\psi(v)\) là một nghiệm hợp lệ, thì nghiệm đối xứng \(\psi'(v) = k + 1 - \psi(v)\) cũng hợp lệ và giữ nguyên khoảng cách \(\card{\psi'(u) - \psi'(v)} = \card{\psi(u) - \psi(v)}\) \parencite{hai2026nohole}. Để phá vỡ đối xứng này và thu hẹp không gian tìm kiếm, phép biểu diễn cố định một đỉnh chốt \(p\) (thường chọn đỉnh có chỉ số 0) thuộc nửa dưới miền nhãn:
\begin{equation}
  \bigwedge_{i=\lceil k/2 \rceil + 1}^k \neg x_{pi}.
  \label{eq:nohole-symmetry-breaking}
\end{equation}

\begin{resultbox}{Kết quả tiêu biểu cho No-hole Anti-\(k\)-labeling}
Trên 23 bộ dữ liệu chuẩn Harwell--Boeing \parencite{hai2026nohole}:
\begin{itemize}
  \item Phép biểu diễn \mbox{\textnormal{\textsc{PSE}}} kết hợp với bộ giải song song PAINLESS giải thành công \textbf{100\% các bài toán}, vượt trội hoàn toàn so với mô hình \mbox{\textnormal{\textsc{DE}}} (bị bùng nổ mệnh đề và không tìm được nghiệm khả thi trên 3 bài toán lớn U, V, X trong thời gian 1800 giây).
  \item So với các công cụ thương mại hàng đầu đại diện cho MIP (Gurobi, CPLEX) và CP (CPLEX CP)---vốn hoàn toàn thất bại không tìm được bất kỳ nghiệm khả thi nào trên 8 bài toán quy mô lớn từ M đến X---phép biểu diễn \mbox{\textnormal{\textsc{PSE}}} cải thiện trung bình \textbf{19\%} chất lượng nghiệm mục tiêu và giảm trên \textbf{82\%} thời gian chạy trên các bài toán giải được.
\end{itemize}
\end{resultbox}

\section{Tô màu theo bandwidth}

\subsection{Quan hệ với tô màu đỉnh}

Khi \(d_{uv}=1\) trên mọi cạnh, BCP trở về tô màu đỉnh theo độ trải. Khi các
trọng số lớn hơn, một đồ thị con đầy đủ không chỉ yêu cầu các màu khác nhau mà
còn yêu cầu một phép đóng gói một chiều của các khoảng cách. Vì vậy kích thước
đồ thị con đầy đủ
vẫn cho cận dưới nhưng không mô tả đủ phép tách có trọng số. Các mô hình MIP
theo phép gán, đại diện hoặc phủ tập trong tô màu đỉnh là điểm khởi đầu; BCP
phải bổ sung các xung đột có xét khoảng cách
\parencite{malaguti2010coloring}.

Từ góc nhìn biểu diễn, biến trực tiếp mở rộng mô hình gán; biến thứ tự nén một dãy
giá trị xung đột thành hai biên; biến khối tạo các đại diện khoảng. Sáu họ dưới
đây khảo sát có hệ thống trục biểu diễn đó và kết nối với các cận lý thuyết đồ
thị đã nêu.

\subsection{Mô hình và sáu họ biểu diễn}

BCP cho đồ thị có trọng số cạnh \(d_{uv}>0\). Mỗi đỉnh nhận màu
\(c(u)\in[1,H]\) và
\[
  \card{c(u)-c(v)}\ge d_{uv}
  \quad\forall uv\in E.
\]
Mục tiêu là tối thiểu
\[
  \chi_B(G)=\min_c\max_{u\in V}c(u).
\]

Nghiên cứu hệ thống của \textcite{nguyen2026bcp} tổ chức không gian thiết kế thành sáu
họ:
\begin{center}
\captionof{table}{Sáu họ biểu diễn cho tô màu theo bandwidth.}
\label{tab:bcp-encodings}
\begin{tabularx}{\textwidth}{@{}>{\bfseries\sffamily}p{.14\textwidth}p{.31\textwidth}Y@{}}
\toprule
Họ & Biến & Cơ chế khoảng cách\\
\midrule
1G & \(g_{u,j}\iff c(u)\ge j\) & Một chuỗi thứ tự hướng lớn hơn\\
1L & \(l_{u,j}\iff c(u)\le j\) & Một chuỗi thứ tự hướng nhỏ hơn\\
2G & Giá trị đúng \(x_{u,j}\) và ngưỡng \(g_{u,j}\) &
Giá trị chính xác kích hoạt hai biên khoảng cấm\\
2L & Giá trị đúng \(x_{u,j}\) và ngưỡng \(l_{u,j}\) &
Như 2G với hướng thứ tự đảo\\
X & Biến khối & Chia sẻ cấu trúc khoảng cách theo khối\\
\(X_a\) & Khối và phép trừ phụ &
Tái sử dụng hiệu/khoảng qua biến phụ\\
\bottomrule
\end{tabularx}
\end{center}

\subsection{Khoảng cấm bằng hai biên}

Giả sử \(x_{u,a}=1\). Khi đó màu của \(v\) phải nằm ngoài
\[
  [a-d_{uv}+1,\;a+d_{uv}-1].
\]
Với ngữ nghĩa \(g_{v,t}=1\iff c(v)\ge t\), điều kiện được viết bằng một mệnh đề
biên:
\[
  \neg x_{u,a}
  \lor
  \neg g_{v,a-d_{uv}+1}
  \lor
  g_{v,a+d_{uv}}.
\]
Hai literal cuối lần lượt nói \(c(v)<a-d_{uv}+1\) hoặc
\(c(v)\ge a+d_{uv}\). Số mệnh đề không còn phụ thuộc vào độ rộng khoảng cấm;
mỗi giá trị chính xác tạo một mệnh đề cho mỗi hàng xóm.

Biểu diễn theo khối chia miền màu thành các đoạn và dùng biến phụ để tóm tắt quan
hệ giữa các đoạn. CNF có thể lớn hơn Order Encoding, nhưng lan truyền trên bài
toán khó có thể mạnh hơn vì bộ giải nhìn thấy các cấu trúc khoảng được chia sẻ.

\begin{workedexample}{Một khoảng cấm được nén thành hai biên}
Cho \(H=7\), \(x_{u,4}=1\) và \(d_{uv}=2\). Đỉnh \(v\) không được nhận
\[
  [4-2+1,\,4+2-1]=[3,5].
\]
Với \(g_{v,t}\Leftrightarrow c(v)\ge t\), chỉ cần mệnh đề
\[
  \neg x_{u,4}\lor\neg g_{v,3}\lor g_{v,6}.
\]
Khi \(x_{u,4}=1\), hai literal còn lại nói đúng hai miền hợp lệ:
\(c(v)\le2\) hoặc \(c(v)\ge6\). Direct Encoding từng cặp phải cấm riêng
\((4,3),(4,4),(4,5)\); biểu diễn theo biên giữ số mệnh đề không đổi khi khoảng
cấm rộng hơn, miễn tập cấm vẫn là một khoảng.
\end{workedexample}

\begin{resultbox}{Hiệu ứng tương tác}
Trên tập GEOM và MS-CAP của \textcite{nguyen2026bcp}, cấu hình theo khối giải
được các bài toán khó mà Order Encoding không hoàn tất trong cùng giới hạn.
Giải tăng dần và đối xứng không có hiệu ứng đồng đều: lợi ích phụ thuộc vào họ
biểu diễn. Vì vậy thí nghiệm phân tích thành phần (ablation study) nên dùng đơn vị
\(\text{phép biểu diễn}\times\text{đặc trưng}\), không chỉ so trung bình từng
đặc trưng.
\end{resultbox}

\section{Tô đa màu theo bandwidth}

\subsection{Mô hình ô màu và giảm miền}

BMCP gán cho đỉnh \(u\) một tập \(w(u)\) màu. Viết chúng thành các ô có thứ
tự
\[
  c(u)_1<c(u)_2<\cdots<c(u)_{w(u)}.
\]
Hai ô cùng đỉnh cách nhau ít nhất \(d(u,u)\); hai ô ở hai đỉnh kề nhau cách ít
nhất \(d(u,v)\). Với cận màu lớn nhất \(k\), miền ô được siết:
\[
  LB_{u,i}=1+(i-1)d(u,u),
\]
\[
  UB_{u,i}=k-(w(u)-i)d(u,u).
\]
Thứ tự ô loại đối xứng hoán vị màu bên trong một đỉnh và đồng thời tạo lan
truyền cận.

Mô hình ô là một phép mở rộng miền, trong khi mô hình tần số là một mô hình
tập. Hai cách có đối xứng và cận dưới khác nhau. Trong mô hình ô, thứ tự nội bộ
loại \(w(u)!\) hoán vị; trong mô hình tập, đối xứng đó không được tạo ra nhưng
ràng buộc ``chọn đúng \(w(u)\) màu'' cần ràng buộc đếm. So sánh FE/SBE/OBE vì
vậy phải bao gồm cả số trạng thái gán trước các mệnh đề khoảng cách.

\subsection{FE, SBE và OBE}

\paragraph{Flat Encoding (biểu diễn phẳng, FE).}
Biến \(x_{u,m}\) cho biết màu \(m\) được dùng tại \(u\). Ràng buộc đếm bảo đảm
đúng \(w(u)\) màu; xung đột từng cặp cấm các cặp quá gần.

\paragraph{Slot-Based Encoding (biểu diễn theo ô màu, SBE).}
Biến \(x_{u,i,m}\) nói ô \(i\) bằng \(m\). Mỗi ô có ràng buộc đúng một và các
cặp ô vi phạm khoảng cách bị cấm trực tiếp.

\paragraph{Order-Based Encoding (biểu diễn theo thứ tự, OBE).}
Thêm \(g_{u,i,m}=1\iff c(u)_i\ge m\). Giá trị chính xác tại một ô kích hoạt
hai biên khoảng cấm của ô ở đỉnh khác:
\[
  x_{u,i,m}\rightarrow
  \left(
    g_{v,j,m+d(u,v)}
    \lor
    \neg g_{v,j,m-d(u,v)+1}
  \right).
\]
Số mệnh đề xung đột giảm một bậc theo miền màu so với SBE từng cặp.

\begin{proposition}[Đúng đắn của OBE]
Nếu liên kết giá trị chính xác--thứ tự xác định cùng một giá trị cho mỗi ô,
chuỗi thứ tự đơn điệu và mọi giá trị chính xác kích hoạt mệnh đề khoảng cấm ở
trên, thì mọi mô hình
OBE giải mã thành một phép tô đa màu theo bandwidth hợp lệ.
\end{proposition}

\begin{proof}
Ràng buộc đúng một cho mỗi ô một màu; các mệnh đề cận và thứ tự giữ các ô trong
miền và đúng thứ tự. Với mỗi cặp ô bị ràng buộc, giá trị chính xác \(m\) của ô
thứ nhất buộc ô thứ hai nằm dưới \(m-d+1\) hoặc từ \(m+d\) trở lên. Vì vậy
khoảng cách tuyệt đối ít nhất \(d\).
\end{proof}

Tìm kiếm tăng dần từ trên xuống bắt đầu từ một cận trên khả thi. Sau khi tìm
mô hình có độ trải \(s\), bộ giải cấm các màu từ \(s\) trở lên và tiếp tục.
Lần \UNSAT đầu tiên tại \(s-1\), cùng mô hình tại \(s\), chứng nhận tối ưu.
\textcite{nguyen2026bmcp} dùng cơ chế này để chứng minh tối ưu cho BMCP.

\section{Một trừu tượng chung: giá trị chính xác và khoảng cấm}

ABP, BCP và BMCP gặp cùng một mẫu:
\begin{enumerate}
  \item một biến chính xác chọn giá trị \(a\);
  \item lựa chọn đó tạo một khoảng giá trị bị cấm cho đối tượng khác;
  \item biến thứ tự nén khoảng thành hai literal biên.
\end{enumerate}
Khác biệt nằm ở nguồn của khoảng và số lựa chọn trong phép gán. ABP có ràng
buộc khác nhau toàn cục; BCP không có; BMCP có nhiều ô trên một đỉnh. Vì vậy
mô-đun \texttt{ForbiddenInterval} có thể dùng lại, còn
\texttt{AssignmentSemantics} phải được tham số hóa theo bài toán.

\section{Quy trình chính xác kết hợp}

Một quy trình có giá trị tham khảo cho cả bốn bài toán là:
\begin{enumerate}
  \item dùng cận lý thuyết đồ thị và phương pháp kinh nghiệm để tạo khoảng
        \([LB,UB]\);
  \item chọn mô hình đúng ngữ nghĩa: hoán vị, một màu hoặc nhiều ô;
  \item lọc miền bằng cận cục bộ hoặc đồ thị con đầy đủ trước khi tạo biến;
  \item chọn Direct Encoding/thứ tự/khối theo hình dạng tập cấm;
  \item giải theo quy trình cận đơn điệu và lưu nghiệm làm chứng;
  \item tại cận cuối, kiểm tra phép gán nhãn miền và chứng minh \UNSAT độc lập.
\end{enumerate}

Kiến trúc này đặt SAT trong một khung giải chính xác thay vì coi bộ giải SAT là
toàn bộ thuật toán. Cấu trúc này cũng làm rõ vị trí đóng góp của một nghiên cứu mới: cận, mô
hình, phép biểu diễn, quy trình tìm kiếm hay chứng nhận.

\begin{summarybox}
\textbf{Bài học trung tâm.} Chọn Direct Encoding, thứ tự hay theo khối không
chỉ là chọn kích thước \CNF. Biến trực tiếp làm phép giải mã và xung đột cục bộ
rõ; biến thứ tự nén bất đẳng thức và khoảng; biến khối tạo điểm chia sẻ mới cho
lan truyền. Một kiến trúc tốt thường kết hợp cả ba ở những lớp khác nhau.
\end{summarybox}

\subsection*{Câu hỏi nghiên cứu}
\begin{enumerate}
  \item Với ABP trên đường \(P_4\), dựng \(F(k)\) tại \(k=2\) và chỉ rõ
        exactly-one theo hàng lẫn cột.
  \item Dẫn xuất mệnh đề biên BCP cho ba trường hợp: khoảng nằm trong
        miền, chạm biên trái và chạm biên phải.
  \item Thiết kế thí nghiệm phân tích thành phần (ablation study) \(2\times2\) cho giải tăng dần và đối xứng
        trên hai họ biểu diễn BCP.
\end{enumerate}

\chapter{Gán nhãn và phân bổ tần số}
\chapterlead{\emph{Radio labeling} (gán nhãn radio),
\emph{\(k\)-safe labeling} (gán nhãn an toàn \(k\)) và
\emph{frequency assignment} (phân bổ tần số) cùng chia sẻ một mô-đun: một lựa
chọn chính xác tạo ra \emph{forbidden interval} (khoảng cấm) trên miền của đối
tượng khác. Sự khác nhau nằm ở cách tạo khoảng cấm và cách định nghĩa hàm mục
tiêu.}

\index{radio labeling}\index{k-safe labeling}
\index{frequency assignment}\index{forbidden interval}
\index{arc-consistency preprocessing}

\flowdiagram{Giá trị chính xác}{Chuỗi thứ tự}{Khoảng cấm}
  {Độ trải hoặc số giá trị}{Nghiệm chứng và cận dưới}

\section{Hai văn liệu nguồn: gán nhãn đồ thị và phân bổ tần số}

Các bài toán trong chương đi vào SAT từ hai văn liệu khác nhau. Gán nhãn đồ
thị nghiên cứu các ánh xạ số trên đỉnh/cạnh dưới nhiều quy tắc khoảng cách;
gán nhãn radio là một nhánh trong hệ phân loại rất rộng đó
\parencite{gallian2024labeling}. Phân bổ tần số bắt nguồn từ mô hình nhiễu vô
tuyến: mỗi máy phát có miền, ràng buộc phân cách, đôi khi có giá trị gán trước,
phân cực, nhiều loại nhiễu và nhiều hàm mục tiêu.

Khảo cứu của Aardal và cộng sự phân loại FAP theo miền, mô hình nhiễu và hàm
mục tiêu, đồng thời tổng hợp MIP, cận dưới lý thuyết đồ thị, tối ưu ràng buộc
và phương pháp kinh nghiệm \parencite{aardal2007frequency}. Từ góc nhìn SAT,
điểm quan trọng là không đồng nhất ba hàm mục tiêu:
\begin{itemize}
  \item \textbf{độ trải nhỏ nhất}: giảm hiệu giữa tần số lớn nhất và nhỏ nhất;
  \item \textbf{số giá trị nhỏ nhất}: giảm số tần số phân biệt được dùng;
  \item \textbf{nhiễu/chi phí nhỏ nhất}: cho phép vi phạm hoặc mức nhiễu có
        trọng số.
\end{itemize}
Chúng có thể dùng cùng các ràng buộc khả thi nhưng cần biến hàm mục tiêu và lập
luận cận dưới khác nhau.

\begin{center}
\captionof{table}{Ba lớp bài toán gán nhãn và tần số.}
\label{tab:labeling-fap-classes}
\begin{tabularx}{\textwidth}{@{}>{\bfseries\sffamily}p{.2\textwidth}YYY@{}}
\toprule
Lớp & Miền & Ràng buộc & Hàm mục tiêu điển hình\\
\midrule
Gán nhãn radio & Liên tục \([1,B]\) & Từ khoảng cách đồ thị & Độ trải/nhãn lớn nhất nhỏ nhất\\
Đơn ánh an toàn \(k\) & Liên tục \([1,B]\) & Khác nhau + khoảng cách cạnh &
Nhãn lớn nhất nhỏ nhất\\
FAP & Rời rạc theo máy phát & Phân cách nhiễu &
Độ trải, số giá trị hoặc chi phí\\
\bottomrule
\end{tabularx}
\end{center}

Các công trình của nhóm trong chương được dùng để minh họa ba phép chuyển:
xung đột từng cặp sang biên khoảng; miền liên tục sang miền rời rạc; và cận độ
trải sang bộ đếm giá trị phân biệt. Ba phép chuyển này kết nối hệ phân loại FAP
và gán nhãn đồ thị với các mô-đun SAT có thể dùng lại.

\section{Bổ đề khoảng cấm}

Cho \(x_{u,a}=1\) khi đối tượng \(u\) nhận giá trị \(a\). Với đối tượng \(v\),
đặt
\[
  g_{v,t}=1\iff f(v)\ge t.
\]
Giả sử lựa chọn \(a\) cấm khoảng \([L(a),U(a)]\) cho \(v\). Điều kiện hợp lệ
là
\[
  f(v)<L(a)
  \quad\lor\quad
  f(v)>U(a).
\]
Trên miền nguyên liên tục, ràng buộc trở thành
\[
  \neg x_{u,a}
  \lor
  \neg g_{v,L(a)}
  \lor
  g_{v,U(a)+1}.
  \tag{11.1}\label{eq:boundary}
\]

\begin{lemma}[Mệnh đề biên]\label{lem:boundary}
Giả sử chuỗi thứ tự \(g_{v,t}\) đúng với ngữ nghĩa
\(f(v)\ge t\). Khi \(x_{u,a}=1\), mệnh đề
\eqref{eq:boundary} được thỏa khi và chỉ khi giá trị của \(v\) nằm ngoài
\([L(a),U(a)]\).
\end{lemma}

\begin{proof}
Nếu \(x_{u,a}=1\), literal đầu sai. Literal
\(\neg g_{v,L(a)}\) đúng khi \(f(v)<L(a)\); literal
\(g_{v,U(a)+1}\) đúng khi \(f(v)\ge U(a)+1\). Phép
tuyển vì thế tương đương với điều kiện nằm ngoài khoảng.
\end{proof}

Trên miền rời rạc \(D_v\), không thể dùng trực tiếp chỉ số \(L,U\). Ta thay
chúng bằng hai biên theo giá trị:
\[
  \lambda_v(a)=\min\set{q\mid D_{v,q}\ge L(a)},
\]
\[
  \rho_v(a)=\min\set{q\mid D_{v,q}>U(a)}.
\]
Các literal biên sau đó dùng \(\lambda_v(a)\) và \(\rho_v(a)\) trong chuỗi thứ
tự theo chỉ số miền.

\subsection{Quan hệ với Support Encoding}

Trong Direct Encoding/hỗ trợ, lựa chọn \(x_{u,a}\) yêu cầu \(v\) nhận một
trong các giá trị còn hỗ trợ:
\[
  \neg x_{u,a}
  \lor
  \bigvee_{b\in D_v:\ b\notin[L(a),U(a)]}x_{v,b}.
\]
Mệnh đề biên nén phép tuyển này thành hai ngưỡng khi tập hỗ trợ là hợp của một
tiền tố và một hậu tố. Đây không phải một mẹo riêng của radio/FAP; đó là phép
phân tích mệnh đề hỗ trợ dựa trên tính lồi của tập cấm. Nếu tập cấm có nhiều
khoảng rời nhau, cần nhiều cặp biên hoặc một
biểu diễn BDD/MDD khác.

\begin{figure}[tb]
\centering
\begin{tikzpicture}[satfig,x=11mm,y=8mm]
  \foreach \a in {1,...,9}{
    \ifnum\a<4
      \node[satbox,minimum width=9mm,fill=PaleBlue] (a\a) at (\a,0) {\(\a\)};
    \else
      \ifnum\a<7
        \node[satbox,minimum width=9mm,fill=PaleGold,draw=Gold] (a\a) at (\a,0)
          {\(\a\)};
      \else
        \node[satbox,minimum width=9mm,fill=PaleTeal] (a\a) at (\a,0) {\(\a\)};
      \fi
    \fi
  }
  \draw[Blue,very thick] (.55,-.75)--(3.45,-.75)
    node[midway,below,satlabel,text=Blue]{\(\neg g_{v,4}\): tiền tố hợp lệ};
  \draw[Gold,very thick] (3.55,.75)--(6.45,.75)
    node[midway,above,satlabel,text=Gold]{khoảng cấm \([4,6]\)};
  \draw[Teal,very thick] (6.55,-.75)--(9.45,-.75)
    node[midway,below,satlabel,text=Teal]{\(g_{v,7}\): hậu tố hợp lệ};
\end{tikzpicture}
\caption{Một khoảng cấm liên tục được nén thành hai literal biên:
\(\neg g_{v,4}\lor g_{v,7}\).}
\label{fig:forbidden-interval-boundaries}
\end{figure}

Phép so sánh lan truyền theo dõi đường suy luận của từng biểu diễn. Mệnh đề hỗ
trợ trực tiếp loại giá trị chính xác; mệnh đề biên suy qua chuỗi thứ tự và liên
kết giá trị chính xác--thứ tự. Hai phép biểu diễn đạt cùng mức cắt tỉa khi đặc tả của các
chuỗi đó cung cấp đủ các bước suy luận
\parencite{gent2002arc,bacchus2007gac}.

\begin{keyidea}{Mô-đun dùng lại}
\texttt{ExactValue} chọn \(a\); \texttt{OrderChain} biểu diễn miền của \(v\);
\texttt{ForbiddenInterval} chỉ cần một hàm trả về hai biên. Gán nhãn radio,
gán nhãn an toàn \(k\) và MO-FAP khác nhau ở hàm này, không khác nhau ở lôgic của
\cref{lem:boundary}.
\end{keyidea}

\section[Radio k-labeling]{Radio \(k\)-labeling: gán nhãn radio \(k\)}

\subsection{Mô hình}

Cho đồ thị liên thông \(G=(V,E)\), khoảng cách đường đi ngắn nhất \(d(u,v)\) và
tham số \(k\). Phép gán nhãn radio \(k\) là ánh xạ
\(f:V\to\mathbb{Z}_{>0}\) thỏa
\[
  \card{f(u)-f(v)}
  \ge
  \Delta_{uv},
  \qquad
  \Delta_{uv}=k+1-d(u,v),
\]
cho mọi cặp có \(\Delta_{uv}>0\). Mục tiêu là tối thiểu màu lớn nhất
\[
  \min\max_{v\in V}f(v).
\]
Trong bài toán gán nhãn radio chuẩn, người ta thường xét
\(k=\operatorname{diam}(G)\).

Khoảng cách đồ thị tạo một đồ thị ràng buộc dày hơn đồ thị gốc: mọi cặp có
\(d(u,v)\le k\) đều bị ràng buộc, với mức phân cách giảm khi khoảng cách đồ
thị tăng. Bước tìm đường đi ngắn nhất cho mọi cặp vì thế không chỉ dành cho bộ
bộ kiểm tra xác định toàn bộ cấu trúc liên kết của phép biểu diễn. Các cặp có
\(\Delta_{uv}\le0\) được bỏ, còn các đỉnh trong cùng quỹ đạo có thể hỗ trợ phá
đối xứng.

\begin{figure}[tb]
\centering
\begin{tikzpicture}[satfig,node distance=19mm]
  \node[satstate] (v1) {\(v_1\)};
  \node[satstate,right=of v1] (v2) {\(v_2\)};
  \node[satstate,right=of v2] (v3) {\(v_3\)};
  \node[satstate,right=of v3] (v4) {\(v_4\)};
  \draw[Ink,very thick] (v1)--(v2)--(v3)--(v4);
  \draw[Blue,dashed,bend left=35] (v1) to node[above,satlabel]{\(d=2\)} (v3);
  \draw[Blue,dashed,bend left=35] (v2) to node[above,satlabel]{\(d=2\)} (v4);
  \draw[Gold,dotted,very thick,bend right=48] (v1) to
    node[below,satlabel]{\(d=3\)} (v4);
  \node[below=4pt of v1,satlabel,text=Teal] {\(f=1\)};
  \node[below=4pt of v2,satlabel,text=Teal] {\(f=4\)};
  \node[below=4pt of v3,satlabel,text=Teal] {\(f=7\)};
  \node[below=4pt of v4,satlabel,text=Teal] {\(f=2\)};
\end{tikzpicture}
\caption{Một phép gán nhãn radio \(3\) khả thi cho \(P_4\). Hai đỉnh kề nhau
cần chênh ít nhất \(3\), cặp cách hai cạnh cần \(2\), và hai đầu mút cần \(1\).}
\label{fig:radio-p4}
\end{figure}

\subsection{SAT trực tiếp}

Với một cận trên \(B\), đặt
\[
  x_{v,\ell}=1\iff f(v)=\ell,
  \qquad \ell\in[1,B].
\]
Mỗi đỉnh nhận đúng một nhãn. Với cặp \(u,v\), mọi cặp nhãn
\((a,b)\) có \(\card{a-b}<\Delta_{uv}\) tạo xung đột
\[
  \neg x_{u,a}\lor\neg x_{v,b}.
\]
Direct Encoding có phép giải mã đơn giản nhưng số xung đột tăng theo độ rộng
khoảng cấm.

\subsection{SAT thứ tự}

Thêm
\[
  g_{v,\ell}=1\iff f(v)\ge\ell.
\]
Biến chính xác được xác định tại điểm chuyển:
\[
  x_{v,\ell}\iff
  g_{v,\ell}\land\neg g_{v,\ell+1}.
\]
Nếu \(u\) nhận \(a\), \(v\) không được nhận nhãn trong
\([a-\Delta_{uv}+1,a+\Delta_{uv}-1]\). Mệnh đề biên là
\[
  \neg x_{u,a}
  \lor
  \neg g_{v,a-\Delta_{uv}+1}
  \lor
  g_{v,a+\Delta_{uv}},
\]
sau khi bỏ các literal nằm ngoài miền. Số mệnh đề khoảng cách trở thành tuyến tính
theo \(B\) cho mỗi cặp đỉnh, thay vì phụ thuộc thêm vào \(\Delta_{uv}\).

\subsection{Tối thiểu hóa độ trải theo cách tăng dần}

Bộ giải được dựng một lần ở cận trên \(B\). Sau khi tìm mô hình có giá trị lớn
nhất \(\beta\), thêm
\[
  \neg g_{v,\beta}
  \qquad\forall v
\]
để buộc mọi nhãn không vượt quá \(\beta-1\). Vòng lặp nhảy theo nghiệm đương
nhiệm thật, không cần giảm từng đơn vị. Nếu lần gọi ở \(\beta-1\) trả \UNSAT,
mô hình tại \(\beta\) cùng cận dưới này chứng nhận tối ưu.

Đối xứng có thể khai thác theo tự đẳng cấu của đồ thị. Với chu trình \(C_n\),
sau khi dịch nhãn nhỏ nhất về \(1\), phép quay đưa đỉnh mang nhãn \(1\) về một
đỉnh neo cố định. Quy tắc neo như vậy phải dựa trên quỹ đạo của đồ thị.

\begin{resultbox}{Gán nhãn radio}
Trên 146 bài toán từ chín họ đồ thị, \textcite{thanh2026radio} báo 38 số radio
mới, khớp hoặc cải thiện \BKS trên 130 bài toán và chứng nhận 109 giá trị tối
ưu khi kết hợp các khung SAT với ILP. SAT thứ tự đặc biệt hiệu quả ở những họ
có đường kính tăng theo kích thước; ILP bổ sung tốt ở các họ có đường kính bị
chặn.
\end{resultbox}

\section[Gán nhãn an toàn k chính xác]{Gán nhãn an toàn \(k\) chính xác}

\subsection{Mô hình được sử dụng}

Trong mô hình của \textcite{thanh2026ksafe}, nhãn
\[
  L:V\rightarrow\set{1,\ldots,B}
\]
thỏa hai điều kiện:
\[
  L_u\ne L_v \qquad\forall u\ne v,
\]
\[
  \card{L_u-L_v}\ge k \qquad\forall uv\in E.
\]
Mục tiêu là tối thiểu \(\max_vL_v\). Điều kiện đầu là tính đơn ánh toàn cục:
mỗi nhãn được dùng nhiều nhất một lần.

\begin{designrule}{Quy ước của chương}
Chương giữ đúng mô hình đơn ánh ở trên. Vì vậy, khi \(k=1\), bài toán vẫn là
gán nhãn đơn ánh; đây không phải bài toán tô màu đồ thị chuẩn. Biến thể không đơn ánh là
một mô hình khác.
\end{designrule}

\subsection{Direct, forward order và reverse order encoding}

\paragraph{Direct Encoding (Direct Encoding, DE).}
\(K_{v,\ell}\) nói \(L_v=\ell\). Ràng buộc đúng một theo đỉnh và AMO theo nhãn
thực hiện phép gán đơn ánh. Mỗi cạnh cấm các cặp nhãn cách dưới \(k\).

\paragraph{Forward Order Encoding (Order Encoding xuôi, FE).}
\(X_{v,j}=1\iff L_v\ge j\). Giá trị chính xác xuất hiện ở điểm chuyển và kích
hoạt hai biên của khoảng \([j-k+1,j+k-1]\).

\paragraph{Reverse Order Encoding (Order Encoding ngược, RE).}
\(X_{v,j}=1\iff L_v\le j\). Ngữ nghĩa đối xứng với FE; hướng của chuỗi thứ tự
và các literal cận được đảo.

FE và RE có cùng tập phép gán nhãn nhưng có thể dẫn bộ giải đi theo các nhánh
khác nhau. Đối xứng phản xạ
\[
  L'_v=S+1-L_v
\]
giữ mọi khoảng cách và tính đơn ánh trên miền \([1,S]\). Có thể chọn một đỉnh
đại diện và đặt biến đại diện ở nửa dưới miền để giữ một đại diện của cặp phản xạ.

\begin{resultbox}{So sánh ba phép biểu diễn}
Trên 16 đồ thị Erdős--Rényi của \textcite{thanh2026ksafe}, FE và RE tìm nghiệm
cho toàn bộ tập và chứng nhận 13 giá trị tối ưu; DE tìm nghiệm cho 15 và chứng
nhận 8. Kết quả cho thấy biên thứ tự phù hợp với ràng buộc khoảng cách, còn
hướng FE/RE tạo khác biệt nhỏ hơn lựa chọn trực tiếp/thứ tự.
\end{resultbox}

\section{Phân bổ tần số với số giá trị nhỏ nhất}

\subsection{Miền rời rạc và hàm mục tiêu}

Mỗi máy phát \(v\) có miền tần số đã sắp
\[
  D_v=\set{D_{v,1}<D_{v,2}<\cdots<D_{v,m_v}}.
\]
Các cạnh mang khoảng cách nhiễu; một số máy phát có giá trị gán trước. Hàm mục
tiêu không phải tần số lớn nhất mà là số giá trị phân biệt được dùng:
\[
  \min\card{\set{f_v\mid v\in V}}.
\]
Do miền có lỗ và khác nhau giữa các máy phát, khoảng cấm phải được
tính theo giá trị thực của tần số rồi ánh xạ về chỉ số miền.

Trong hệ phân loại FAP, đây là mô hình số giá trị nhỏ nhất: hai nghiệm dùng cùng
số giá trị phân biệt có cùng giá trị mục tiêu dù độ trải khác nhau. Do đó chuỗi
thứ tự theo \emph{chỉ số miền cục bộ} chỉ phục vụ tính khả thi; biến \(A_a\)
theo \emph{giá trị toàn cục} mới phục vụ hàm mục tiêu. Trộn hai loại thứ tự này là một
lỗi mô hình phổ biến.

\begin{workedexample}{Cùng độ trải nhưng khác số giá trị}
Xét ba máy phát có miền chứa \(\{1,4,7\}\), và không có cạnh nhiễu giữa máy
phát thứ hai với thứ ba. Hai phương án
\[
  f=(1,4,7),\qquad f'=(1,7,7)
\]
đều có tần số lớn nhất \(7\), nhưng lần lượt dùng \(3\) và \(2\) giá trị phân
biệt. Tối thiểu hóa độ trải xem chúng ngang nhau; FAP số giá trị nhỏ nhất ưu
tiên \(f'\). Vì vậy biến hàm mục tiêu phải là \(A_1,A_4,A_7\), không phải chỉ
một ngưỡng ``tần số lớn nhất không vượt \(7\)''.
\end{workedexample}

\subsection{Tiền xử lý nhất quán cung}

Với cặp \((v,w)\), một giá trị \(a\in D_v\) phải có ít nhất một giá trị hỗ trợ
\(b\in D_w\) thỏa khoảng cách yêu cầu. Nếu không có, \(a\) bị loại. Quá trình
lặp đến \emph{fixpoint} (điểm bất động, tức không còn giá trị nào bị loại):
\[
  D_v\leftarrow
  \set{a\in D_v\mid
    \exists b\in D_w:\card{a-b}\ge d_{vw}}
\]
cho mọi cung ràng buộc. Tiền xử lý này không thay giá trị tối ưu vì chỉ loại giá
trị không thể xuất hiện trong bất kỳ nghiệm khả thi nào.

Đây là AC-3 trên các ràng buộc phân cách nhị phân. Khả năng suy diễn lan truyền của AC-3 bị giới hạn ở
hỗ trợ từng cạnh: một giá trị có thể có hỗ trợ riêng trên mọi đỉnh kề nhưng
không có hỗ trợ đồng thời toàn cục. Sau tiền xử lý, học mệnh đề SAT xử lý các
tổ hợp xung đột còn lại. Vì vậy nên báo kích thước miền tại điểm bất động và
thời gian AC riêng; không gán toàn bộ mức giảm thời gian chạy cho POSE.

\subsection{DSE và POSE}

Biểu diễn SAT trực tiếp dùng \(x_{v,q}=1\iff f_v=D_{v,q}\), ràng buộc đúng một
theo máy phát và xung đột từng cặp cho các cặp giá trị vi phạm.

Biểu diễn SAT thứ tự bộ phận thêm
\[
  g_{v,q}=1\iff
  \text{chỉ số tần số của }v\text{ không nhỏ hơn }q.
\]
Nếu \(x_{v,q}=1\), tần số \(D_{v,q}\) tạo khoảng cấm
\([D_{v,q}-d+1,D_{v,q}+d-1]\) trong miền \(D_w\). Hai biên
\(\lambda_w,\rho_w\) của miền rời rạc được đưa vào mệnh đề biên
\[
  \neg x_{v,q}
  \lor
  \neg g_{w,\lambda_w(q)}
  \lor
  g_{w,\rho_w(q)}.
\]
POSE vì thế loại sự phụ thuộc vào số giá trị nằm bên trong khoảng cấm.

\begin{figure}[tb]
\centering
\begin{tikzpicture}[satfig,x=13mm,y=9mm]
  \node[left,satlabel] at (0,1) {\(D_v\)};
  \foreach \i/\a in {1/1,2/4,3/7,4/10}{
    \ifnum\a=7
      \node[satbox,fill=PaleBlue,draw=Blue] at (\i,1) {\(\a\)};
    \else
      \node[satbox] at (\i,1) {\(\a\)};
    \fi
  }
  \node[left,satlabel] at (0,0) {\(D_w\)};
  \foreach \i/\a in {1/2,2/5,3/8,4/11}{
    \ifnum\i=1
      \node[satbox,fill=PaleTeal] at (\i,0) {\(\a\)};
    \else
      \ifnum\i=4
        \node[satbox,fill=PaleTeal] at (\i,0) {\(\a\)};
      \else
        \node[satbox,fill=PaleGold,draw=Gold] at (\i,0) {\(\a\)};
      \fi
    \fi
  }
  \draw[Gold,decorate,decoration={brace,amplitude=4pt}] (1.55,-.85)--(3.45,-.85)
    node[midway,below=4pt,satlabel,text=Gold]
      {bị cấm bởi \(f_v=7,d=4\)};
  \node[satlabel,align=left] at (7.4,.5)
    {\(\lambda_w=2\) tại giá trị \(5\)\\
     \(\rho_w=4\) tại giá trị \(11\)\\
     hỗ trợ còn lại: \(\{2,11\}\)};
\end{tikzpicture}
\caption{Ánh xạ khoảng giá trị \([4,10]\) vào hai biên của miền rời rạc
\(D_w=\{2,5,8,11\}\).}
\label{fig:fap-discrete-boundaries}
\end{figure}

\begin{workedexample}{Mệnh đề biên trên miền có lỗ}
Với các miền trong \cref{fig:fap-discrete-boundaries}, chọn giá trị thứ ba của
\(D_v\), tức \(x_{v,3}=1\) và \(f_v=7\). Khoảng cách \(d=4\) cấm
\([4,10]\) ở \(D_w\). Mệnh đề POSE là
\[
  \neg x_{v,3}\lor\neg g_{w,2}\lor g_{w,4}.
\]
Khi \(x_{v,3}=1\), \(\neg g_{w,2}\) cho phép chỉ số \(1\) (tần số \(2\));
\(g_{w,4}\) cho phép chỉ số \(4\) (tần số \(11\)). Hai giá trị \(5,8\) nằm
giữa bị loại mà không cần hai mệnh đề xung đột riêng.
\end{workedexample}

\subsection{Đếm số tần số phân biệt}

Với mỗi giá trị tần số toàn cục \(a\), tạo biến
\[
  A_a=1\iff
  \exists v:\ f_v=a.
\]
Liên kết \(x_{v,q}\rightarrow A_{D_{v,q}}\) bảo đảm mọi tần số được dùng đều
được đánh dấu. Hàm mục tiêu trở thành một ràng buộc đếm:
\[
  \min\sum_aA_a.
\]
Sequential Counter tăng dần được dựng một lần đến cận trên ban đầu. Sau khi mô hình
dùng \(q\) tần số, lần gọi kế tiếp kích hoạt cận \(q-1\) bằng các giả định. Lần
\UNSAT tại \(q-1\) cùng mô hình tại \(q\) chứng nhận tối ưu.

\begin{resultbox}{MO-FAP}
Trên 35 bài toán CALMA/DUTtest hoặc được mô hình hóa lại,
\textcite{dao2026mofap} báo cấu hình POSE kết hợp bộ đếm tăng dần chứng nhận
trạng thái khả thi hoặc vô nghiệm của toàn bộ tập thử. Nhất quán cung
giảm mạnh tổng thời gian của cấu hình này. Kết quả cho thấy vai trò bổ sung của
ba mô-đun: lọc miền, nén khoảng cấm và đếm giá trị phân biệt.
\end{resultbox}

\section{So sánh theo cấu trúc}

\begin{center}
\captionof{table}{So sánh cấu trúc Radio \(k\), \(k\)-Safe và MO-FAP.}
\label{tab:labeling-structure}
\begin{tabularx}{\textwidth}{@{}>{\bfseries\sffamily}p{.2\textwidth}YYY@{}}
\toprule
Đặc trưng & Radio \(k\) & \(k\)-Safe & MO-FAP\\
\midrule
Miền & \([1,B]\) liên tục & \([1,B]\) liên tục & \(D_v\) rời rạc\\
Khoảng cấm & Từ khoảng cách đồ thị & \(k\) trên cạnh & Từ nhiễu\\
Cặp bị ràng buộc & Mọi cặp có \(\Delta>0\) &
Cạnh và tính đơn ánh & Cạnh ràng buộc\\
Hàm mục tiêu & Nhãn lớn nhất nhỏ nhất & Nhãn lớn nhất nhỏ nhất & Số giá trị khác nhau nhỏ nhất\\
Lớp mục tiêu & Cận độ trải & Cận độ trải & Bộ đếm giá trị phân biệt\\
\bottomrule
\end{tabularx}
\end{center}

Ba bài toán chia sẻ mô-đun khả thi ở mức khoảng, nhưng lớp hàm mục tiêu khác
nhau. Radio và \(k\)-safe thay cận trên của nhãn. MO-FAP giữ các
giá trị số cố định và tối thiểu số giá trị được kích hoạt. Đây là lý do không
nên đồng nhất ``số giá trị nhỏ nhất'' với ``độ trải nhỏ nhất''.

\section{Đối chứng ngoài SAT và cận dưới}

Đối với độ trải gán nhãn, MIP dùng biến nhãn nguyên và hướng nhị phân cho trị
tuyệt đối; CP dùng ràng buộc khác nhau và khoảng cách; các phép dựng lý thuyết
đồ thị cho cận hoặc nghiệm trên từng họ đồ thị. Đối với FAP, MIP/sinh cột, tối
ưu ràng buộc, tìm kiếm cục bộ và \emph{tabu search} (tìm kiếm kiêng kỵ) là các đối
chứng đã được phát triển lâu
dài \parencite{aardal2007frequency}. Một đánh giá SAT
có giá trị nên phân biệt:
\begin{enumerate}
  \item bài toán mà SAT cải thiện nghiệm đương nhiệm;
  \item bài toán mà SAT đóng khoảng cách bằng chứng nhận \UNSAT;
  \item bài toán mà phương pháp khác cung cấp cận dưới/cận trên đầu vào;
  \item bài toán vô nghiệm do miền/giá trị gán trước, có thể phát hiện trước
        bước tối ưu.
\end{enumerate}

Với số giá trị nhỏ nhất, cận dưới cơ bản có thể đến từ cấu trúc đồ thị con đầy
đủ hoặc cấu trúc nhiễu
hoặc số giá trị bắt buộc bởi các giá trị gán trước; cận trên đến từ phương án
tần số kinh nghiệm. Bộ đếm giá trị phân biệt SAT tiếp nhận hai cận và lần lượt đóng
khoảng \([LB,UB]\).

\section{Bộ kiểm tra và chứng nhận}

Bộ kiểm tra khả thi đọc trực tiếp nghiệm:
\begin{itemize}
  \item radio: kiểm miền, đường đi ngắn nhất cho mọi cặp và mọi
        \(\Delta_{uv}>0\);
  \item \(k\)-safe: kiểm miền, tính đơn ánh và phân cách trên cạnh;
  \item MO-FAP: kiểm \(f_v\in D_v\), giá trị gán trước, mọi ràng buộc
        và số tần số phân biệt.
\end{itemize}
Bộ kiểm tra không đọc các biến thứ tự phụ. Điều này giúp bộ kiểm tra giữ kích thước nhỏ gọn và độc lập với
phép biểu diễn.

Chứng nhận tối ưu gồm một nghiệm làm chứng tại cận trên và bằng chứng cận dưới:
\UNSAT ở độ trải/số giá trị kề trước, hoặc một cận dưới toán học bằng cận trên.
Với các giả định tăng dần, có thể hiện thực hóa (reification) các giả định cuối thành mệnh đề đơn
vị để tạo một \CNF độc lập cho bộ kiểm tra chứng minh.

\begin{summarybox}
\begin{enumerate}
  \item Khoảng cấm là trừu tượng chung của cả ba bài toán.
  \item Trên miền liên tục, khoảng dùng hai ngưỡng trực tiếp; trên miền
        rời rạc, phải tìm biên theo giá trị.
  \item Tối thiểu hóa độ trải và tối thiểu hóa số giá trị phân biệt cần hai mô
        đun hàm mục tiêu khác nhau.
  \item Biến trực tiếp phục vụ lựa chọn chính xác và giải mã; biến thứ tự nén
        khoảng; Sequential Counter đếm số giá trị được dùng.
\end{enumerate}
\end{summarybox}

\subsection*{Câu hỏi nghiên cứu}
\begin{enumerate}
  \item Chứng minh \cref{lem:boundary} khi khoảng chạm một biên của miền.
  \item Với \(D_v=\set{10,30,40}\) và khoảng \([15,35]\), xác định hai
        literal biên của POSE.
\end{enumerate}


\chapter{Kết luận và hướng nghiên cứu}
\index{encoding!lựa chọn}
\index{bằng chứng thực nghiệm}
\index{giới hạn phương pháp}
\index{chứng nhận tối ưu}

\chapterlead{Chương kết nối các kết quả kỹ thuật thành một luận điểm chung:
một phép biểu diễn SAT tốt là một giao diện có ngữ nghĩa, được chọn theo cấu trúc
bài toán và được đánh giá bằng một chuỗi bằng chứng có thể kiểm tra. Từ đó,
chương tổng hợp miền hiệu quả của từng họ phương pháp và xây dựng chương trình
nghiên cứu tiếp theo.}

\section{Những kết luận bền vững}

Các họ bộ biểu diễn tạo thành một biên Pareto (Pareto frontier) thay vì một thứ hạng duy nhất. Quy
trình lựa chọn có thể dùng lại gồm bốn lớp.

\paragraph{Lớp ngữ nghĩa.}
Trước khi đếm biến và mệnh đề, phải viết đặc tả cho biến chính, biến phụ, bộ
giải mã và hàm mục tiêu. Tính đúng, tính đầy đủ theo phép chiếu và khả năng dùng
literal phụ ở mô-đun khác là ba yêu cầu khác nhau. Một biến Tseitin chỉ được
xem như ngưỡng, biến làm chứng hay trạng thái khi các mệnh đề thật sự bảo đảm ngữ
nghĩa đó.

\paragraph{Lớp cấu trúc.}
Phép biểu diễn nên khớp với hình học của ràng buộc. Cửa sổ chồng lấn gợi ý trạng
thái chia sẻ; miền hữu hạn có thứ tự gợi ý literal thứ tự; tổng Boolean gợi
ý bộ đếm, cây tổng, mạng đếm hoặc BDD; phép tuyển hình học gợi ý biến quan hệ.
Đối xứng nên được loại trước hết ở cách chọn biến, sau đó mới bằng ràng buộc có
phép ánh xạ bảo toàn nghiệm.

\paragraph{Lớp giải.}
Đánh giá ở lớp giải kết hợp số mệnh đề với lan truyền đơn vị, đường truyền qua
biến phụ, khả năng giữ mệnh đề đã học khi thay cận, chi phí của bộ giải mã/bộ
kiểm tra và tương tác với tiền xử lý. Cùng một phép biểu diễn khả thi có thể cư xử
khác nhau dưới SAT độc lập, SAT tăng dần, MaxSAT hay một kiến trúc kết hợp.

\paragraph{Lớp bằng chứng.}
Phát biểu về tính đúng cần chứng minh hoặc kiểm tra vét cạn trên miền nhỏ; phát
biểu về kích thước cần công thức và bộ đếm từ gói tái lập; phát biểu về hiệu
năng cần quy trình, đối chứng, giới hạn thời gian và dữ liệu từng lần chạy.
Tối ưu chỉ được chứng nhận khi cận trên từ một nghiệm khả thi gặp cận dưới hợp
lệ, thường là một kết quả
\UNSAT ở cận kế tiếp.

\section{Tổng hợp các nghiên cứu tình huống}

Bảng~\ref{tab:reported-evidence-map} tổng hợp thiết kế đánh giá, kết quả và cơ
chế biểu diễn nổi bật của các nghiên cứu tình huống.

\begin{longtable}{@{}
>{\raggedright\arraybackslash}p{.19\textwidth}
>{\raggedright\arraybackslash}p{.23\textwidth}
>{\raggedright\arraybackslash}p{.39\textwidth}@{}}
\caption{Tổng hợp các nghiên cứu tình huống.}
\label{tab:reported-evidence-map}\\
\toprule
Mô-đun/bài toán & Thiết kế đánh giá & Kết quả chính\\
\midrule
\endfirsthead
\toprule
Mô-đun/bài toán & Thiết kế đánh giá & Kết quả chính\\
\midrule
\endhead
SCL/SCAMO và Anti-Bandwidth &
Phân tích kích thước SCAMO; 24 đồ thị kiểm chuẩn &
Thang chia sẻ giảm phần bộ đếm lặp trên AMO bậc thang và cạnh tranh trong
cấu hình Anti-Bandwidth đã thử; độ chồng lấn khoảng là đặc trưng quyết định
\parencite{truong2025scamo}.\\

NSC/Xếp nhóm chơi golf &
66 bài toán; NSC, ADDER, BDD, CARD, BINOMIAL &
NSC giải \(56/66\); ADDER, BDD và CARD cùng \(54/66\), BINOMIAL \(32/66\).
Kết quả định vị NSC như một lựa chọn hiệu quả cho ràng buộc đúng \(K\) trong mô hình
bài toán xếp nhóm chơi golf của nguồn \parencite{kieu2026golfer}.\\

ALSCE/Lập lịch điều dưỡng &
28 bài toán, đến 120 điều dưỡng và 24 tuần &
ALSCE giải toàn bộ tập; kết quả hỗ trợ việc chia sẻ thanh ghi cho các cửa sổ
chuỗi trong mô hình này \parencite{truong2025nurse}.\\

SAT/Đóng gói dải &
41 bài toán; SAT thường, SAT tăng dần và CP &
Số giá trị tối ưu được chứng nhận lần lượt là \(26\), \(21\) và \(28\). SAT
tăng dần phát huy lợi thế khi thay cận giữ được phần lớn miền và mệnh đề đã học
\parencite{kieu2025strip}.\\

Gán nhãn radio &
146 bài toán từ chín họ đồ thị; các khung SAT và ILP &
Công trình báo 109 giá trị tối ưu được chứng nhận, 38 số radio mới, và khớp
hoặc cải thiện BKS trên 130 bài toán; SAT và ILP có vai trò bổ sung theo cấu
trúc đồ thị
\parencite{thanh2026radio}.\\

Gán nhãn an toàn \(k\) &
16 đồ thị Erdős--Rényi; FE, RE và DE &
FE và RE tìm nghiệm cho \(16/16\), chứng nhận 13 giá trị tối ưu; DE tìm nghiệm
cho 15 và chứng nhận 8. Biên thứ tự phù hợp hơn biểu diễn trực tiếp trong tập thử
này \parencite{thanh2026ksafe}.\\

POSE/MO-FAP &
35 bài toán CALMA/DUTtest hoặc được mô hình hóa lại &
POSE kết hợp bộ đếm tăng dần quyết định toàn bộ tập thử; kết quả hỗ trợ
kiến trúc lọc miền, nén khoảng cấm và đếm số tần số khác nhau
\parencite{dao2026mofap}.\\
\bottomrule
\end{longtable}

\section{Miền hiệu quả và điểm chuyển thiết kế}

\paragraph{Bộ đếm chia sẻ và Ladder AMK.}
Bộ đếm chia sẻ được ưu tiên khi các cửa sổ chồng lấn mạnh, cận \(k\) nhỏ và
mỗi cửa sổ cắt ít khối. Khi số biên tăng, bộ đếm ghép thay bộ nối hai ngôi;
với các tập gần rời nhau, một phép biểu diễn cục bộ thường gọn hơn.

\paragraph{NSC.}
NSC có kích thước \(O(nk)\) và tiết kiệm vùng trạng thái tam giác bất khả. Cấu trúc này
phù hợp nhất khi vùng này chiếm tỷ lệ đáng kể và bài toán dùng một cận cố định.
Cây tổng/mạng đếm phù hợp với nhiều cận tăng dần; khi \(k\) gần \(n\), đếm
literal bù chuyển bài toán về ngưỡng \(n-k\) nhỏ hơn.

\paragraph{Phá vỡ đối xứng.}
Mỗi ràng buộc phá vỡ đối xứng được suy từ một nhóm tự đẳng cấu bảo toàn miền, ràng
buộc và hàm mục tiêu. Sức chứa, miền quay và cách hàm mục tiêu đọc chỉ số quyết
định những hoán vị nào thuộc nhóm này. Mức phá vỡ đối xứng được chọn cùng độ dài
mệnh đề và chiến lược phân nhánh.

\paragraph{Biểu diễn thời gian và đóng gói.}
SAT theo mốc thời gian phù hợp với chân trời vừa phải và miền bắt đầu đã được
lọc. Tọa độ thứ tự nén tốt các so sánh trên miền nguyên. DDD mở rộng lựa chọn
cho chân trời lớn khi có tập mốc đầy đủ hữu hạn, cận dưới hợp lệ và bước tinh
chỉnh tăng nghiêm ngặt.

\paragraph{Khoảng cấm và biên thứ tự.}
Hai literal biên nén hiệu quả một khoảng cấm trên miền có thứ tự và đặc tả thứ
tự hai chiều. Miền rời rạc dùng ánh xạ phần tử trước/sau; mỗi khoảng rời nhau
bổ sung một cặp biên. Nhất quán cung tạo nhiều giá trị nhất trước SAT khi miền
lớn hoặc đồ thị ràng buộc đủ dày.

\section{Hướng nghiên cứu tiếp theo}

\begin{enumerate}
  \item \textbf{Lựa chọn bộ biểu diễn theo bài toán.} Xây bộ chọn có thể giải thích
  quyết định bằng các đặc trưng như \(n,k\), độ chồng lấn, độ thưa của miền và
  số cận sẽ truy vấn; đánh giá bộ chọn so với bộ giải tốt nhất ảo.
  \item \textbf{Đặc tả và kết hợp.} Phát triển thư viện mô-đun mà mỗi literal
  đầu ra có ngữ nghĩa máy đọc được, cho phép kiểm tự động việc ghép bộ đếm,
  liên kết biểu diễn, hàm mục tiêu và bộ giải mã.
  \item \textbf{Chứng nhận toàn trình.} Nối ghi chứng minh của SAT/MaxSAT với
  bộ kiểm tra miền và chứng nhận cận dưới để giá trị tối ưu không phụ thuộc vào niềm tin ở
  bộ sinh.
  \item \textbf{Trạng thái chia sẻ tổng quát.} Mô hình hóa việc chọn khối,
  cây ghép và ngưỡng cần hiện thực hóa (reification) như một bài toán tối ưu riêng; tìm cận dưới
  cho số trạng thái/bộ nối dưới các yêu cầu lan truyền khác nhau.
  \item \textbf{Phương pháp kết hợp có phân công rõ.} Dùng CP/MIP để lọc miền
  hoặc tạo cận, SAT/MaxSAT để xử lý phép tuyển, nhưng đánh giá bằng thí nghiệm
  phân tích thành phần (ablation) để biết lợi ích đến từ mô-đun nào.
  \item \textbf{Gói tái lập bền vững.} Dùng vùng chứa phần mềm, bản kê dữ liệu,
  mã băm, giấy phép và quy trình tái lập; tách kết luận đã công bố khỏi kết quả
  thăm dò.
\end{enumerate}

\section{Kết luận cuối}

Giá trị của phép biểu diễn SAT nằm ở khả năng chọn đúng trạng thái trung gian để
CDCL nhìn thấy cấu trúc của bài toán. Bộ đếm chia sẻ, thanh ghi tam giác,
biên thứ tự, biến quan hệ và rời rạc hóa động là các biểu hiện khác nhau của
cùng một nguyên tắc: chỉ hiện thực hóa (reification) thông tin vừa đủ, đặt đặc tả rõ cho thông tin
đó, rồi kiểm bằng cả chứng minh lẫn gói tái lập. Ba bước này chuyển một
nghiên cứu tình huống thành một phương pháp có thể kiểm chứng, so sánh và dùng
lại.

\backmatter

\renewcommand{\thetable}{T.\arabic{table}}
\renewcommand{\theHtable}{T.\arabic{table}}
\setcounter{table}{0}
\chapter*{Thuật ngữ Anh--Việt}
\addcontentsline{toc}{chapter}{Thuật ngữ Anh--Việt}
\markboth{Thuật ngữ Anh--Việt}{Thuật ngữ Anh--Việt}
\begin{longtable}{@{}>{\bfseries\sffamily}p{.3\textwidth}p{.62\textwidth}@{}}
\caption{Thuật ngữ Anh--Việt dùng trong sách.}
\label{tab:glossary}\\
\toprule
Thuật ngữ tiếng Anh & Tương đương và cách dùng trong sách\\
\midrule
\endfirsthead
\toprule
Thuật ngữ tiếng Anh & Tương đương và cách dùng trong sách\\
\midrule
\endhead
arc consistency & nhất quán cung: mỗi giá trị còn ít nhất một giá trị hỗ trợ
trên mỗi cung ràng buộc\\
assumption & giả định: literal tạm thời cho một lần gọi bộ giải\\
at-least-\(k\) & ràng buộc có ít nhất \(k\) literal đúng\\
at-most-\(k\) & ràng buộc có nhiều nhất \(k\) literal đúng\\
auxiliary variable & biến phụ có đặc tả ngữ nghĩa\\
benchmark & bộ kiểm chuẩn: tập bài toán thử và quy trình đánh giá\\
binary decision diagram & biểu đồ quyết định nhị phân, viết tắt BDD\\
bound tightening & siết cận miền hoặc cận mục tiêu\\
branch-and-bound & tìm kiếm phân nhánh và cận\\
cardinality constraint & ràng buộc đếm số literal đúng\\
certified optimum & nghiệm tối ưu có cận dưới tương ứng\\
channeling & liên kết biểu diễn: các mệnh đề nối hai cách biểu diễn cùng một
quyết định\\
clause learning & học mệnh đề từ xung đột\\
conflict-driven clause learning & học mệnh đề theo xung đột, viết tắt CDCL\\
consistency & mức nhất quán hoặc lọc miền mà phép mã hóa/bộ lan truyền đạt được\\
constraint programming & lập trình ràng buộc, viết tắt CP\\
contract & đặc tả ngữ nghĩa mà một mô-đun cung cấp cho mô-đun sử dụng\\
decoder & bộ giải mã mô hình Boolean thành nghiệm của bài toán gốc\\
direct encoding & mã hóa trực tiếp: một literal giá trị cho mỗi giá trị miền\\
dominance rule & quy tắc trội: loại nghiệm bị trội mà vẫn giữ ít nhất một
nghiệm tối ưu\\
encoding & phép mã hóa một ràng buộc hoặc mô hình thành \CNF\\
exactly-\(k\) & ràng buộc có đúng \(k\) literal đúng\\
feasibility & tính tồn tại nghiệm thỏa mọi ràng buộc cứng\\
feasible solution & nghiệm khả thi\\
forbidden interval & khoảng cấm do một lựa chọn giá trị tạo ra\\
global constraint & ràng buộc toàn cục có ngữ nghĩa và bộ lan truyền riêng\\
incremental SAT & SAT tăng dần: dùng lại trạng thái bộ giải qua nhiều lần gọi\\
lazy clause generation & sinh mệnh đề lười: kết hợp lan truyền miền với lời
giải thích dạng mệnh đề\\
literal & biến Boolean hoặc phủ định của biến đó\\
lower bound & cận dưới hợp lệ của giá trị mục tiêu\\
model checker & bộ kiểm tra trực tiếp nghiệm đã giải mã\\
non-preemption & tác vụ chạy trong một khoảng liên tục, không bị ngắt\\
order encoding & mã hóa thứ tự: chuỗi literal ngưỡng biểu diễn thứ tự trong miền\\
performance profile & đồ thị tỷ lệ hiệu năng chuẩn hóa theo cấu hình tốt nhất\\
proof logging & ghi chứng minh: xuất chứng cứ mệnh đề để bộ kiểm tra xác minh
\UNSAT\\
propagation complete & đầy đủ lan truyền trong phạm vi đã định nghĩa\\
pseudo-Boolean constraint & bất đẳng thức tuyến tính trên biến Boolean\\
reification & vật hóa chân trị: gắn một literal với tính đúng/sai của một ràng
buộc\\
relaxation & mô hình nới lỏng dùng để tạo cận hoặc phân tích\\
restart & khởi động lại nhánh tìm kiếm nhưng giữ các mệnh đề đã học\\
selection network & mạng lựa chọn: mạng bộ so sánh chỉ tạo các đầu ra cần thiết\\
shifted geometric mean & trung bình nhân dịch chuyển để tổng hợp thời gian chạy\\
sorting network & mạng sắp xếp xây từ các bộ so sánh Boolean\\
soundness & tính đúng: mọi mô hình Boolean giải mã thành nghiệm hợp lệ\\
support encoding & mã hóa hỗ trợ: liệt kê các giá trị hỗ trợ cho một lựa chọn\\
symmetry breaking & phá đối xứng: giữ đại diện chuẩn tắc của mỗi quỹ đạo\\
symmetry breaking predicate & mệnh đề giữ ít nhất một đại diện của mỗi quỹ đạo\\
timeout & quá thời hạn: lần chạy vượt giới hạn thời gian\\
totalizer & cây tổng dạng đơn phân cho ràng buộc đếm/PB\\
unit propagation & lan truyền đơn vị\\
upper bound & cận trên đến từ một nghiệm khả thi\\
virtual best solver & bộ giải tốt nhất ảo: phép chọn hậu nghiệm cấu hình nhanh
nhất cho từng bài toán thử\\
witness & biến chứng: literal chỉ ra một nhánh của phép tuyển\\
\bottomrule
\end{longtable}

\printbibliography[heading=bibintoc,title={Tài liệu tham khảo}]
\printindex

\end{document}
